\documentclass[11pt,letterpaper]{article}

% --- Encoding & fonts ---
\usepackage[utf8]{inputenc}
\usepackage[T1]{fontenc}
% Note: lmodern intentionally omitted — the default Computer Modern font
% has full ec-lm coverage in this build environment.
\usepackage[expansion=false,protrusion=false]{microtype}
\usepackage{textcomp}                 % needed by some Unicode passthroughs

% --- Math ---
\usepackage{amsmath}
\usepackage{amssymb}
\usepackage{amsthm}
\usepackage{mathtools}

% --- Tables (pandoc emits longtable) ---
\usepackage{longtable}
\usepackage{booktabs}
\usepackage{array}
\usepackage{calc}

% --- Lists ---
\usepackage{enumitem}
\setlist{nosep,leftmargin=*,topsep=2pt}

% --- Layout ---
\usepackage{geometry}
\geometry{margin=1in}
\usepackage{parskip}                  % paragraph spacing rather than indent

% --- Pandoc compatibility helpers ---
\providecommand{\tightlist}{%
  \setlength{\itemsep}{0pt}\setlength{\parskip}{0pt}}
\providecommand{\pandocbounded}[1]{#1}

% --- Hyperlinks (last among style packages) ---
\usepackage[colorlinks=true,linkcolor=blue,citecolor=blue,urlcolor=blue]{hyperref}

% --- Theorem environments ---
% The source carries its own theorem numbers (e.g. "Theorem 5.2",
% "Lemma C.10"). Use a custom theoremstyle that places the qualifier
% inline as the head note, and \newtheorem*{...} so amsthm doesn't add
% its own counter-based number.
%
% The default amsthm rendering of `\begin{theorem*}[Foo]` is "Theorem (Foo)."
% with the qualifier in parens. The source's qualifiers are like "5.2" or
% "5.2 (Dimension-transcendence)" — so the parens nest: "Theorem (5.2
% (Dimension-transcendence))." which reads badly. Use a custom theoremstyle
% to render the qualifier with a leading space and a period: "Theorem 5.2."
% or "Theorem 5.2 (Dimension-transcendence)."
\newtheoremstyle{paperstyle}%   name
  {\topsep}%                    spaceabove
  {\topsep}%                    spacebelow
  {\itshape}%                   bodyfont
  {0pt}%                        indent
  {\bfseries}%                  headfont
  {.}%                          headpunct
  { }%                          headspace
  {\thmname{#1}\thmnumber{ #2}\thmnote{ #3}}% headspec  -> "Kind <num> <note>"

\newtheoremstyle{paperstyledef}%
  {\topsep}{\topsep}{\normalfont}{0pt}{\bfseries}{.}{ }%
  {\thmname{#1}\thmnumber{ #2}\thmnote{ #3}}

\newtheoremstyle{paperstylerem}%
  {\topsep}{\topsep}{\normalfont}{0pt}{\itshape}{.}{ }%
  {\thmname{#1}\thmnumber{ #2}\thmnote{ #3}}

\theoremstyle{paperstyle}
\newtheorem*{theorem}{Theorem}
\newtheorem*{lemma}{Lemma}
\newtheorem*{proposition}{Proposition}
\newtheorem*{corollary}{Corollary}
\newtheorem*{conjecture}{Conjecture}

\theoremstyle{paperstyledef}
\newtheorem*{definition}{Definition}
\newtheorem*{example}{Example}

\theoremstyle{paperstylerem}
\newtheorem*{remark}{Remark}
\newtheorem*{observation}{Observation}

% --- Title metadata ---
\title{Dimension-Transcendent Attractors in\\
       Generalized Kaprekar Routines}
\author{Clay Elmore\\
        Independent Researcher\thanks{Correspondence: \texttt{clayelmore@att.net}.}}
\date{April 2026}

\begin{document}
\maketitle

\begin{abstract}
The classical Kaprekar routine at four digits iterates
$K(n) = \alpha(n) - \beta(n)$, where $\alpha(n)$ and $\beta(n)$
are the integers formed by arranging $n$'s digits in descending and
ascending order. For every non-repdigit four-digit input, the orbit
converges in at most seven steps to the fixed point $6174$.
Generalizing the construction to arbitrary digit length $d$, the
analogous rule converges to $495$ at $d = 3$ but fails at $d = 5$:
orbits enter cycles rather than fixed points.

This failure is commonly described as a peculiarity of five-digit
arithmetic. We argue that it is instead a peculiarity of the specific
subtraction rule. The classical ``descending minus ascending''
construction is one specific point in a space of $d! \cdot (d! - 1)$
ordered permutation-pair rules at each digit length, and its behavior
across $d$ is not uniform: the algebraic rank of the classical rule is
$d$ at $d = 4$ but strictly less than $d$ at $d = 3$ and
$d = 5$ because of forced middle-digit cancellation. The failure at
$d = 5$ reflects the classical rule being rank-$4$ there, not a
failure of Kaprekar convergence.

When we classify all permutation-pair rules at each digit length,
universal full-variable fixed points exist at $d = 4$, $d = 5$, and
$d = 6$. We complete this classification exhaustively at $d \leq 6$,
identifying $33$ universal fixed points at $d = 5$ and $506$ at
$d = 6$. We then ask which fixed points extend from one digit length
to the next: if $F$ is a universal full-variable fixed point at digit
length $d$, does some rule at $d + 1$ also have $F$ as a universal
fixed point?

Among the $33$ universal fixed points at $d = 5$, exactly one ---
$F = 60714$ --- is a universal full-variable fixed point at $d = 6$.
The remaining $32$ are dimension-locked at $d = 5$: no full-variable
rule at $d = 6$ attracts them universally. The selection of $60714$
from this population is the central discovery; the lifting machinery is
the test it uniquely passes, not a construction tailored to fit it.
Given this selection, we construct an explicit family of
permutation-pair rules at each $d \geq 5$, connected by a
coefficient-preserving lifting recipe, and prove:

\begin{theorem}
$60714$ is a fixed point at every digit length
$d \geq 5$ under the coefficient-preserving lifting family constructed
in §5. The corresponding dynamics are universal on $A_d \setminus E_d$
at every $d \geq 5$, where $A_d$ is the admissible input set and
$E_d$ is an exactly-characterized escape class. The escape class is
empty at $d = 5, 6$, and at $d \geq 7$ has the same structural form
--- block-aligned multisets that collapse to $0$ under the rule --- as
the classical Kaprekar escape class at $d = 3, 4$ (the ten repdigits
at $d = 3$ and the ten multiples of $1111$ at $d = 4$, which under
the classical rule each map to $0$ rather than to the nontrivial fixed
point).
\end{theorem}

This is, to our knowledge, the first explicitly constructed attractor in
the generalized Kaprekar family for which a fixed-point equation
persists across infinitely many dimensions, with universal convergence
on $A_d \setminus E_d$ at every $d \geq 5$ and a structurally
characterized escape class continuous with the classical case (§7.6).
The classical Kaprekar constant $6174$ exhibits a related but distinct
pattern, proven here: universal at $d = 4$ (native), algebraically
obstructed at $d = 5$, and universal on $A_d \setminus E_d^{6174}$
at every $d \geq 6$ under coefficient-preserving lifting, where
$E_d^{6174}$ is a precisely characterized $45$-input class at
$d \geq 8$. $60714$, $6174$, and $60417$ share the digit
multiset $\{7, 6, 4, 1\}$, and within this thread the empirical
invariant $\mathrm{sum\_locked\_spans}$ orders cross-dimensional
behavior: $60714$ ($\mathrm{sls} = 5$) is universal on
$A_d \setminus E_d$ at every $d \geq 5$ with $E_d = \emptyset$ at
$d = 5, 6$, while $6174$ ($\mathrm{sls} = 8$) is universal on
$A_d \setminus E_d^{6174}$ at every $d \geq 6$. The invariant orders
within-thread behaviors but does not generalize to a cross-multiset
principle.
\end{abstract}

\section{Introduction}\label{introduction}

\subsection{Kaprekar's routine and the middle-digit
cancellation}\label{kaprekars-routine-and-the-middle-digit-cancellation}

In 1949, D. R. Kaprekar observed a striking property of four-digit
integers in base ten \cite{Kaprekar1955}. For any four-digit $n$ with
at least two distinct digits, let $\alpha(n)$ be the integer formed by
arranging $n$'s digits in descending order, and $\beta(n)$ the
integer formed by arranging them in ascending order. Iterating

\[K(n) = \alpha(n) - \beta(n)\]

produces a sequence that converges to $6174$ within at most seven
steps, regardless of starting value.

At three digits, the analogous construction converges to $495$
\cite{Trigg1974}. At five digits, it fails: every non-repdigit orbit
enters one of three cycles rather than reaching a fixed point
\cite{Prichett1981}. This failure is typically described as a
peculiarity of five-digit arithmetic. We argue that it reveals something
deeper about the classical rule itself.

Expand $K$ at $d = 3$ algebraically. For sorted-descending digits
$(a, b, c)$ with $a \geq b \geq c$ and at least two distinct digits,
$K(a, b, c) = \overline{abc} - \overline{cba}$. The subtraction's
borrow pattern is forced: the ones place $c - a$ borrows; the middle
place becomes $b - b - 1 = -1$ after the borrow and must borrow again;
the hundreds place becomes $a - c - 1$ after the second borrow. The
result simplifies to

\[K(a, b, c) = 99 \, (a - c),\]

which depends on only $a$ and $c$. The middle digit $b$ cancels
identically. Algebraically, $K$ at $d = 3$ has coefficient vector
$(99, \, 0, \, -99)$: the coefficient on the middle position is
exactly zero.

This is not a computational observation --- it is a structural fact
about the classical construction. At every odd digit length, the
classical rule's center-position coefficient vanishes through the same
forced borrow chain. The classical rule at $d = 3$ is \emph{not} a
rank-$3$ object; it is a rank-$2$ object wearing rank-$3$ clothes.
The classical rule at $d = 5$ is likewise rank-$4$, not rank-$5$.
The failure of the classical rule to produce fixed-point convergence at
$d = 5$ is therefore not surprising once one asks about rank-$5$
rules: the classical rule \emph{isn't one}.

This observation opens the question we pursue in this paper. If the
classical rule is a specific algebraic construction whose rank varies
with $d$, and if convergence properties depend on rank, then the
natural object of study is not the classical rule across $d$ but
\emph{the space of all rearrangement-subtraction rules} at each $d$.
The classical rule is one specific point in that space. The question
becomes: which integers are universal fixed points, which rules attract
them, and how do the rules relate across digit lengths?

\subsection{Permutation-pair rules}\label{permutation-pair-rules}

Let $n$ be a positive integer padded to $d$ digits, and let
$\mathrm{sort\_desc}(n) = (x_0, x_1, \ldots, x_{d-1})$ denote its
sorted-descending digit sequence, with
$x_0 \geq x_1 \geq \cdots \geq x_{d-1}$. Given any ordered pair of
permutations $\pi, \sigma \in S_d$ with $\pi \neq \sigma$, the
\textbf{permutation-pair rule} is

\[K_{\pi, \sigma}(n) = \bigl| \pi \cdot \mathrm{sort\_desc}(n) - \sigma \cdot \mathrm{sort\_desc}(n) \bigr|,\]

where $\pi \cdot (x_0, \ldots, x_{d-1})$ denotes the integer
$\sum_{i} 10^{d-1-i} x_{\pi_i}$ --- the digits rearranged by $\pi$.
The classical rule is $\pi = (0, 1, \ldots, d-1)$ and
$\sigma = (d-1, d-2, \ldots, 0)$ (descending minus ascending), which
we denote $K_0$.

Every permutation-pair rule expands linearly in the sorted-descending
digits:

\[K_{\pi, \sigma}(n) = \left| \sum_{i=0}^{d-1} c_i \, x_i \right|,\]

where the integer coefficients $c_i$ depend only on $(\pi, \sigma)$
and satisfy $\sum_i c_i = 0$. Specifically,
$c_i = 10^{d-1-\pi^{-1}(i)} - 10^{d-1-\sigma^{-1}(i)}$. Define the
\textbf{algebraic rank} (or \textbf{surviving-variable count}) of the
rule:

\[\mathrm{sv}(K_{\pi, \sigma}) = \bigl|\{ i : c_i \neq 0 \}\bigr|.\]

The rule is \textbf{full-variable} at digit length $d$ if
$\mathrm{sv} = d$ --- every sorted-descending position participates
nontrivially in the output. The classical rule is full-variable at
$d = 4$ (coefficient vector $(999, 90, -90, -999)$) but not at odd
$d$, where the middle-position coefficient vanishes by the
cancellation of §1.1.

A rule $K$ is \textbf{universal} for a fixed point $F$ at digit
length $d$ if $K(F) = F$ and every non-repdigit $d$-digit input
iterates to $F$ under $K$. We adopt the standard basin convention
excluding repdigits (multisets with one distinct digit) and
near-repdigits (multisets where one digit appears $d - 1$ or $d$
times); these degenerate inputs are excluded from Kaprekar-style
convergence analysis throughout the literature
\cite{Prichett1981,Dahl2026}.

\subsection{The organizing question}\label{the-organizing-question}

A note on the order of discovery, before the formalism. The
investigation reported in this paper did not begin with the integer
$60714$ as a target. It began with a structural puzzle suggested by
§1.1: the classical operator $K_0$ has $6174$ as a fixed point at
$d = 4$, but at $d = 5$ it has no fixed point at all, and there is
no analogue of the classical construction that lifts $K_0$ into the
rank-$5$ space. We asked instead whether the \emph{space of all
permutation-pair rules} at $d = 5$ contains rules with universal fixed
points, and whether any such fixed points persist across digit lengths.
Exhaustive enumeration at $d = 5$ produced $33$ universal
full-variable fixed points (Theorem 3.3); the exhaustive
cross-dimensional check (Theorem 4.1) found that exactly one of these
--- $60714$ --- extends to $d = 6$ as a universal full-variable
fixed point of \emph{some} rule there. The remaining $32$ are
dimension-locked at $d = 5$.

The coefficient-preserving lifting framework introduced in §5 is the
\emph{mechanism} we discovered as the explanation for this
cross-dimensional persistence. The recipe is uniform --- applied without
modification across all $33$ candidates --- and $60714$ is the
integer that survives the test. Once the lifting is identified and
$60714$ is selected, Theorem 5.2 proves that the lifted family is
universally convergent on $A_d \setminus E_d$ at every $d \geq 5$,
where $E_d$ is the structurally characterized escape class of §5.2.
The escape class is empty at $d = 5, 6$. At $d \geq 7$ it is
non-empty but takes the same structural form as the classical Kaprekar
escape class at $d = 3, 4$: a set of inputs whose orbit reaches $0$
in finitely many applications of the rule, with the inputs characterized
by an algebraic block-alignment condition (single-block repdigits at
$d = 3$, single-block multiples of $1111$ at $d = 4$ --- both of
which collapse in a single step --- and multi-block aligned multisets at
$d \geq 7$, which collapse in at most a few steps; for $d \leq 11$,
exhaustive enumeration confirms collapse in at most $4$ iterations).
See §7.6 for the full continuity argument. Universal convergence on
$A_d \setminus E_d$ is a substantive dynamical claim that does not
follow from the lifting's algebraic preservation property; it requires
the proof developed in §5 and Appendix C. The discovery sequence ---
observation of the rank gap, classification at low $d$,
cross-dimensional check, structural mechanism, dynamical theorem --- is
the methodological spine of §1.4 and is recapitulated in detail there.

Fix an integer $F$. Define its \textbf{native digit length} $d_F$ to
be the smallest $d$ at which $F$ is a universal full-variable fixed
point of some rule $K_{\pi, \sigma}$ at $d$. If $F$ has no native
digit length in this sense, it is outside the scope of this paper.

The organizing question is then:

\begin{quote}
\emph{For a fixed point $F$ with native digit length $d_F$, does
$F$ remain a universal full-variable fixed point at any digit length
$d > d_F$? If so, what is the structural relationship between the rule
at $d_F$ and the rule at $d$?}
\end{quote}

The first part of this question has an immediate affirmative answer for
one class of fixed points: the dimension-agnostic family (including
$45$, $495$, $450$) whose rules have $\mathrm{sv} < d$ at every
$d$. These fixed points persist across digit lengths precisely because
their rules cancel interior digits, reducing to lower-rank operations.
We discuss them in §2.5 and exclude them from the full-variable
classification by the $\mathrm{sv} = d$ constraint.

Among \emph{full-variable} fixed points, the question is considerably
sharper. The classical result says: at $d = 4$, the classical rule
$K_0$ is full-variable (rank $4$) and universal for $6174$. But
$K_0$ at $d = 5$ has rank $4$, not $5$, so it is not even in the
$\mathrm{sv} = 5$ space at $d = 5$; the question of whether $6174$
is a universal \emph{full-variable} fixed point at $d = 5$ is not the
question of whether $K_0$ converges there. It is the question of
whether any rank-$5$ rule, anywhere in the $14{,}280$-rule space at
$d = 5$, has $6174$ (padded to $06174$) as a universal fixed
point. The answer, exhaustively verified in §3 and §4, is no: no
rank-$5$ rule even has $6174$ as a fixed point at $d = 5$.

The full-variable universal fixed points at $d = 5$ are thirty-three
specific integers --- none of them $6174$. Among these thirty-three,
the question of §1.3 becomes: which extend to $d = 6$ as universal
rank-$6$ fixed points? Which further extend to $d = 7$, $d = 8$,
and beyond?

\subsection{On method: from brute enumeration to structural
proof}\label{on-method-from-brute-enumeration-to-structural-proof}

A word on how this investigation proceeded. The classical Kaprekar
routine has an elegance accessible to anyone who can subtract: arrange
the digits in descending order, arrange them in ascending order,
subtract, iterate. Convergence to $6174$ is a surprise that anyone can
verify with a calculator, and the rule requires no definition beyond
ordinary arithmetic. Any generalization in the spirit of this result
should ideally preserve some comparable transparency.

Our approach proceeds in the opposite order. We did not begin with a
structural hypothesis about which integers ought to be
dimension-transcendent attractors. We began with brute-force
enumeration: at each digit length $d$, compute the universal
full-variable fixed points by testing every permutation-pair rule
exhaustively against every non-repdigit input. At $d = 3$, $4$,
$5$, $6$ this enumeration is tractable --- the largest, $d = 6$,
requires testing $190{,}800$ full-variable rules against $4{,}905$
admissible digit multisets ($999{,}450$ admissible padded six-digit
strings), a search that runs in hours on commodity hardware but is
impractical to organize, verify, and iterate on without automated
scaffolding. From the resulting lists of fixed points, we asked the
cross-dimensional question empirically: which fixed points at $d$
reappear at $d + 1$?

The enumerations themselves, the data-structuring of the resulting
fixed-point catalogues, and the systematic probe of structural
hypotheses against the empirical evidence were carried out using a
methodological pattern --- exhaustive brute search, pattern recognition
on the output, proposed structural explanation, adversarial re-testing
of the proposal against further enumeration, repeat --- that we believe
is increasingly relevant to exploratory mathematics. The computer
assistance used throughout the project --- including AI tools alongside
conventional scripting --- is documented in the Acknowledgments.

\subsubsection{The 54 → 60714 discovery
arc}\label{the-54-60714-discovery-arc}

To make the method concrete, we describe the specific discovery process
that led to this paper's main theorem. An early candidate for
dimension-transcendence was $F = 54$. At $d = 2$ under the classical
rule, $54$ is a universal fixed point. At $d = 5$, the full-variable
classification (§3) finds $54$ as one of $33$ universal fps under a
specific rank-$5$ rule --- the rule
$\pi = (1, 2, 4, 3, 0), \sigma = (2, 0, 3, 4, 1)$ with coefficient
vector $(-90, 99, 9000, -9000, -9)$. Initial empirical work suggested
$54$ as a natural candidate for the dimension-transcendence story,
since it appears at both $d = 2$ and $d = 5$.

Closer examination revealed a structural complication. $54$'s
sorted-descending form at $d = 5$ is $(5, 4, 0, 0, 0)$ --- three of
the five digits are zero. In the native rule at $d = 5$, three of the
five coefficients are paired with these zero digits and contribute
nothing to $K(F) = F$. The effective rank of the rule at $F$ ---
what we formalize in §2.6 as $\mathrm{sv}_F$ --- is $2$, not $5$.
$54$ is algebraically an sv$=5$ fixed point, but structurally it
inherits its fixed-point status from a rank-$2$ algebraic operation on
$(5, 4)$ with the other three positions absorbing coefficients that
cancel internally. The fixed-point equation $K(54) = 54$ at $d = 5$
is effectively the same equation as at $d = 2$, dressed up to occupy
$d = 5$ without adding new structural content.

This observation sharpened the question. A genuine
dimension-transcendent fixed point would need $\mathrm{sv}_F = d_F$
--- all coefficients of its native rule contributing to $K(F) = F$
through nonzero digits. Among the $33$ fixed points at $d = 5$, we
searched for fixed points with no zero digits (so that
$\mathrm{sv}_F = 5$ automatically). Exhaustive enumeration at
$d = 5 \to d = 6$ then identified the unique fixed point satisfying
this criterion that also extends to $d = 6$ via a
coefficient-preserving lifting: $F = 60714$. The existence of a
lifting at $d = 6$ is automatic by a structural argument (Proposition
5.1); the nontrivial content is that the lifting is \emph{universal} at
$d = 6$ and, we prove, at every $d \geq 5$.

The paper's main theorem is thus the answer to a question that the
$54$ counterexample forced us to formulate precisely: among universal
full-variable fixed points with $\mathrm{sv}_F = d_F$ at their native
digit length, which admit universal coefficient-preserving liftings at
every higher digit length?

\subsubsection{Conclusion}\label{conclusion}

This methodological sequence --- exhaustive empirical search,
identification of counterexamples that refine the question, structural
proof on the refined target --- is explicit rather than disguised. The
classical Kaprekar routine is elegant because its rule is stated in a
single sentence; the result we present here is structural in a different
sense. The rule varies with $d$ (via coefficient-preserving lifting),
the construction is explicit, the proof is finite-state-plus-algebraic,
and the fixed point $60714$ is distinguished not by aesthetic priority
but by being the one case where the proof closes out uniformly across
all $d \geq 5$.

\subsection{The main results}\label{the-main-results}

This paper answers the following. (Throughout this section, theorems are
presented as overview statements with section pointers; the body
restates them with section-numbered identifiers --- Theorem 1 here
corresponds to Theorems 3.1--3.4 in §3, Theorem 2 to Theorem 4.1 in §4,
Theorem 3 to Theorem 5.2 in §5, and Theorem 4 to Theorem 6.1 in §6.)

\begin{theorem}[1 (Classification, §3)]
The universal
full-variable fixed points at digit lengths $d \leq 6$ are as
follows:
\end{theorem}

\begin{itemize}
\tightlist
\item
  \emph{$d = 3$: no universal full-variable fixed points. ($12$
  full-variable rules enumerated exhaustively.)}
\item
  \emph{$d = 4$: four universal full-variable fixed points ---
  $1746$, $2538$, $5382$, $6174$. ($216$ full-variable rules
  enumerated exhaustively.)}
\item
  \emph{$d = 5$: thirty-three universal full-variable fixed points.
  ($5{,}280$ full-variable rules enumerated exhaustively.)}
\item
  \emph{$d = 6$: five hundred and six universal full-variable fixed
  points. ($190{,}800$ full-variable rules enumerated exhaustively.)}
\end{itemize}

The classifications at each $d$ are complete, reproducible, and
definitive. Full tables appear in Appendix A.

\begin{theorem}[2 (Cross-dimensional cross-check, §4)]
Of the
thirty-three universal full-variable fixed points at $d = 5$, exactly
one --- $F = 60714$ --- is a universal full-variable fixed point at
$d = 6$ under some rule. The remaining thirty-two are dimension-locked
at $d = 5$: no rank-$6$ rule at $d = 6$ is universal for any of
them.
\end{theorem}

The proof is by exhaustive computation: for each of the thirty-three
fixed points, we enumerate all rank-$6$ rules fixing $F$ (padded to
six digits), test each for universality, and observe that $60714$
alone admits a universal rule.

\begin{theorem}[3 (Dimension-transcendence of $60714$, §5)]
$F = 60714$ is a universal full-variable fixed point at every
digit length $d \geq 5$, under an explicit family of
coefficient-preserving lifting rules derived from its native rule at
$d = 5$.
\end{theorem}

The lifting family consists of two ladders: an odd ladder
$d = 5, 7, 9, \ldots$ and an even ladder $d = 6, 8, 10, \ldots$. At
each step of each ladder, the rule at $d$ is obtained from the rule at
$d - 2$ by appending a pair of coefficients summing to zero. The
coefficients at positions corresponding to $60714$'s nonzero digits
are preserved throughout; the coefficients at positions corresponding to
$60714$'s zero digits absorb the structural content of the lifting.

The proof reduces to two lemmas. Lemma 5.1 (structural) establishes that
the set of inputs with two trailing zeros in their sorted-descending
form is invariant under any zero-sum pair lifting, with the lifted rule
acting on this set as the base rule at $d - 2$ on the non-tail digits.
Lemma 5.2 (the main analytic content) establishes that every
non-near-repdigit orbit at $d \geq 7$ reaches the tail-two-zeros
absorbing set in a bounded number of iterations. Lemma 5.2 is proved by
two complementary techniques: exhaustive finite-state enumeration over
digit multisets at $d = 7, 8, \ldots, 16$ (approximately $25$
million multisets, zero exceptions), and a $d$-independent algebraic
argument at $d \geq 15$ (odd ladder, one-step $T_d$ closure) and a
bounded-reaching-time argument at $d \geq 16$ (even ladder, at most
two-step $T_d$ closure), based on core non-negativity under
sorted-descending inputs (Lemmas 5.3 and 5.4; full proofs in Appendix
C). Together with direct verification at the ladder roots $d = 5$ and
$d = 6$, the two lemmas yield universality at every $d \geq 5$ by
induction along each ladder.

\begin{theorem}[4 (The $\{7, 6, 4, 1\}$-thread, §6)]
The digit
multiset $\{7, 6, 4, 1\}$ (with zero padding as needed) is the support
of universal full-variable fixed points at three distinct digit lengths:
$d = 4$ ($6174$, $1746$), $d = 5$ ($60714$, $60417$), and
$d = 6$ ($60714$, $146070$, $170460$, $607140$). The classical
Kaprekar constant $6174$ has a verified cross-dimensional pattern:
\end{theorem}

\begin{itemize}
\tightlist
\item
  \emph{$d = 4$: strict-universal (native).}
\item
  \emph{$d = 5$: algebraic obstruction --- no rank-$5$ rule fixes
  $6174$.}
\item
  \emph{$d = 6$: dynamic obstruction --- four rank-$6$ rules fix
  $6174$, best basin $0.969$.}
\item
  \emph{$d = 7$: strict-universal via the coefficient-preserving
  lifting $(efgabcd, fgedcba)$.}
\item
  \emph{$d = 8$, $d = 9$: near-universal, with a constant escape
  class of $45$ multisets of the form $(X, X, X, X, 0, 0, \ldots)$;
  basin asymptotes to $1$ as $d$ grows.}
\end{itemize}

\emph{$60714$, $6174$, and $60417$ are unified under a structural
invariant, the $\mathrm{sum\_locked\_spans}$ of their native rules ---
taking values $5$, $8$, $7$ respectively --- which empirically
controls the number of absorbing (zero-digit) positions at which a
coefficient-preserving lifting succeeds. This unification is empirical,
not proven.}

\subsection{\texorpdfstring{How the framework selects
$60714$}{How the framework selects 60714}}\label{how-the-framework-selects-60714}

The result of this paper proceeds in two stages, with importantly
different epistemic statuses. \textbf{The first stage is discovery.} At
digit length $d = 5$, exhaustive enumeration over the $5{,}280$
full-variable permutation-pair rules identifies $33$ universal
full-variable fixed points (Theorem 3.3). At $d = 6$, the analogous
enumeration identifies $506$ such fixed points (Theorem 3.4). These
classifications are computational facts, not constructions. Theorem 4.1
then asks a natural cross-dimensional question: of the $33$ universal
fixed points at $d = 5$, how many admit a coefficient-preserving
lifting to $d = 6$ that is also universal at $d = 6$? The answer is
\textbf{exactly one}: $60714$. The other $32$ are obstructed ---
$17$ by algebra (no rank-$6$ rule satisfies $K(F) = F$), $15$ by
dynamics (the rules that fix them have basin strictly less than $1$).
This selection --- one survivor from a population of $33$ --- is the
central discovery, and it is not a construction. We did not choose the
lifting machinery to fit $60714$; $60714$ is the integer that
survived the lifting test.

\textbf{The second stage is extension.} Given $60714$'s success at
$d = 5 \to d = 6$ under the discovered lifting recipe, the same recipe
applied at every higher $d$ produces a coefficient-preserving lifted
rule whose fixed-point equation $K^{(d)}(60714) = 60714$ holds by
construction. The recipe is uniform across the ladder --- an odd ladder
$d = 5, 7, 9, \ldots$ and an even ladder $d = 6, 8, 10, \ldots$,
each obtained from the previous by appending a pair of coefficients
summing to zero (§5). Two clarifications about what is and is not
engineered here. \emph{What is uniform across all candidates}: the
recipe itself, which was applied identically to all $33$ d=5 fixed
points and selected $60714$ as its unique strict-universal survivor at
$d = 6$. \emph{What is specific to $60714$}: the appended
coefficient pairs at higher $d$ vanish on $60714$'s zero-digit
positions (which is what makes the fixed-point equation persist along
the ladder), but this is a consequence of the same recipe applied to the
same fixed point, not a separate per-$d$ design choice. The
non-trivial content of Theorem 5.2 is therefore the dynamical claim
about basins --- that $K^{(d)}$ is universally convergent at every
$d \geq 5$ on $A_d \setminus E_d$, with $E_d$ exactly
characterized --- and this claim does not follow from algebraic
preservation alone. It requires the structural induction of §5 and the
support lemmas of Appendix C.

The escape class $E_d$ at $d \geq 7$ is continuous with the
classical Kaprekar phenomenon: at $d = 3, 4$ classical Kaprekar
dynamics already exhibit an escape class --- the ten repdigits at
$d = 3$ and the ten multiples of $1111$ at $d = 4$ --- each of
which iterates to $0$ rather than to the nontrivial fixed point. The
classical literature treats these as ``inadmissible.'' The escape class
at $d \geq 7$ is the same phenomenon at higher digit lengths: a
structurally characterized set of inputs whose orbit collapses to $0$.
What is new is not the existence of an escape class, but the explicit
closed form for $|E_d^{(1)}|$ (Lemma 5.5) and the conjecture that
$|E_d|/|A_d| \to 0$ across the lifting family (Conjecture 7.6). See
§7.6 for the full continuity argument.

\textbf{The shape of the result.} Other universal fixed points exhibit
cross-dimensional transcendence in partial forms:

\begin{itemize}
\tightlist
\item
  $6174$ transcends at $d = 4$ (Kaprekar's classical result), and is
  near-universal at every $d \geq 6$ with monotonically improving
  basin under coefficient-preserving lifting. At $d = 5$ it is
  algebraically obstructed (no sv$=5$ rule fixes it). At $d = 8, 9$,
  the escape class consists of $45$ multisets of form
  $(X, X, X, X, Y, Y, Y, Y)$ for $X > Y \geq 0$, all collapsing to
  $0$.
\item
  $146070$ and $607140$ (both in the $\{7, 6, 4, 1\}$-thread,
  native at $d = 6$) transcend to $d = 7$ via the same
  coefficient-preserving machinery; $146070$'s ladder extension is
  verified through $d = 12$.
\item
  Among the $506$ universal fixed points at $d = 6$, empirical
  surveys indicate transcendence to $d = 7$ is \emph{common}, not
  rare, at higher native digit lengths. These surveys are tabulated in
  Appendix D but not proven as theorems in this paper.
\end{itemize}

$60714$'s distinction is therefore not ``uniqueness as a transcendent
fixed point'' --- many other fixed points transcend partially --- but
\textbf{uniqueness as the case that survives Theorem 4.1's selection
test among the $d = 5$ universal full-variable fixed points}. The
classification machinery (Theorems 3.3 and 3.4) and the
cross-dimensional selection (Theorem 4.1) are the durable contributions;
the integer $60714$ is what those tools select.

\subsection{Relation to existing
literature}\label{relation-to-existing-literature}

The literature on Kaprekar routines generalizes the classical
construction along three principal axes, each distinct from ours.

The first varies the base while fixing the classical rule. Prichett,
Ludington, and Lapenta (1981) \cite{Prichett1981} and subsequent work by
Peterson and Pulapaka (2008) \cite{Peterson2008}, Devlin and Zeng (2021)
\cite{Devlin2021}, and Kay and Downes-Ward (2022, 2024)
\cite{Kay2022,Kay2024} classify classical-rule fixed points and cycles
across integer bases.

The second axis fixes the classical rule but varies structural
properties of the fixed points. Iwasaki (2024) \cite{Iwasaki2024}
partitions classical-rule Kaprekar fixed points across digit lengths
into five disjoint sets composed of blocks drawn from
$\{495, 6174, 36, 123456789, 27, 875421, 09\}$ via Diophantine
constraints. The relationship to our work is one of disjoint operator
families. Iwasaki's classification covers fixed points of the classical
rule $K_0(n) = \alpha(n) - \beta(n)$ at every digit length where they
exist; ours covers fixed points of \emph{non-classical} permutation-pair
rules at every digit length where they exist. The two classifications
operate on different operators and produce disjoint fixed-point
families. As a concrete check: $60714$ is not a classical-rule
Kaprekar number --- under $K_0$ at $d = 5$ it iterates to $74943$
in one step and enters the classical $d = 5$ four-cycle
$74943 \to 62964 \to 71973 \to 83952 \to 74943$. Conversely, none of
Iwasaki's seed numbers $\{495, 6174, 36, 123456789, 27, 875421, 09\}$
is a fixed point of the lifted-$60714$ rule at its native digit length
under our coefficient-preserving construction. The two frameworks
address structurally adjacent questions (classifications of
digit-permutation fixed points across $d$) on operationally disjoint
inputs, and we view them as complementary rather than competing.

The third axis is information-theoretic. Dahl (2026) \cite{Dahl2026}
studies classical-rule entropy contraction on coarse-grained digit-gap
coordinates.

All three axes fix the classical rule and vary structural or algebraic
properties of its action. Closest to our framework is the three-paper
series of Nuez (2021--2022) \cite{Nuez2021,Nuez2021b,Nuez2022}, who
develops a parametric formulation of the classical rule using
symmetric-digit differences and catalogues the transformation functions
acting on these parameters. Our parametric observations at the classical
rule --- in particular, the middle-digit cancellation at odd digit
lengths (§1.1) --- are special cases of Nuez's framework, and we credit
him accordingly. Our axis of generalization is distinct: we fix the base
at ten and vary the rule across the full space of ordered permutation
pairs at each digit length, reaching $33$ universal full-variable
fixed points at $d = 5$ that do not appear in any classical-rule
analysis.

The systematic classification of universal full-variable fixed points at
$d \leq 6$, the formulation of the cross-dimensional question in terms
of permutation-pair liftings, and the proof of Theorem 3 at every
$d \geq 5$ are, to our knowledge, all new.

\subsection{Organization}\label{organization}

§2 establishes the formal framework: permutation-pair rules, algebraic
rank, universal fixed points, native digit length, and effective rank at
$F$. §3 presents the exhaustive classifications at $d = 3, 4, 5, 6$.
§4 proves Theorem 2 (the $d = 5 \to d = 6$ cross-check). §5 develops
the coefficient-preserving lifting framework and proves Theorem 3. §6
develops the $\{7, 6, 4, 1\}$-thread and proves Theorem 4. §7 is a
short discussion of open questions. Appendices A--E contain the full
classification tables, the $60714$ ladder through $d = 20$, the
proofs of the support lemmas, observed transcendent fixed points at
$d = 7$ (non-exhaustive, honestly labeled), and the $6174$
$d = 8, 9$ audit details.

\section{Framework}\label{framework}

This section establishes the formal objects, notation, and basic
properties used throughout the paper. Everything here is definitional
and verifiable; no substantive claims are made in this section.

\subsection{Sorted-descending digit
sequences}\label{sorted-descending-digit-sequences}

Throughout, $d \geq 2$ is a fixed digit length (the digit length will
vary between sections but is fixed within each statement). For a
non-negative integer $n < 10^d$, the \textbf{padded digit string} of
$n$ at length $d$ is the sequence
$(n_{d-1}, n_{d-2}, \ldots, n_1, n_0)$ where $n_j$ is the digit at
place $10^j$, padding with leading zeros as needed. Thus
$n = \sum_{j = 0}^{d-1} 10^j \, n_j$.

The \textbf{sorted-descending form} of $n$ at length $d$ is the
tuple $x(n) = (x_0, x_1, \ldots, x_{d-1})$ obtained by sorting the
padded digit string in non-increasing order, so
$x_0 \geq x_1 \geq \cdots \geq x_{d-1}$. For example, at $d = 5$
with $n = 12345$, the padded digit string is $(1, 2, 3, 4, 5)$ and
the sorted-descending form is $(5, 4, 3, 2, 1)$.

An integer $n$ is a \textbf{repdigit} at length $d$ if all digits of
its padded form are equal, i.e., $x_0 = x_{d-1}$. An integer is a
\textbf{near-repdigit} at length $d$ if its sorted-descending form has
one digit appearing $d - 1$ or $d$ times (the $d$-case coincides
with repdigit). All classical Kaprekar analyses exclude repdigits; our
convention additionally excludes near-repdigits. This follows standard
usage in the literature \cite{Prichett1981,Dahl2026} and is necessary in
our setting because near-repdigits can produce degenerate trajectories
under coefficient-preserving liftings, discussed at §5.6. The
\textbf{admissible input set} at length $d$, denoted $A_d$, is the
set of non-repdigit, non-near-repdigit $d$-digit integers.

\subsection{Permutation-pair rules}\label{permutation-pair-rules-1}

Let $S_d$ be the symmetric group on $\{0, 1, \ldots, d-1\}$. Given
an ordered pair $(\pi, \sigma) \in S_d \times S_d$ with
$\pi \neq \sigma$, define the \textbf{permutation-pair rule}
$K_{\pi, \sigma}: \mathbb{Z}_{\geq 0} \to \mathbb{Z}_{\geq 0}$ by

\[K_{\pi, \sigma}(n) = \left| \pi \cdot x(n) - \sigma \cdot x(n) \right|,\]

where, for any permutation $\tau \in S_d$ and tuple
$x = (x_0, \ldots, x_{d-1})$,

\[\tau \cdot x \;=\; \sum_{i = 0}^{d-1} 10^{\tau_i} \, x_i.\]

Thus $\tau \cdot x$ is the integer obtained by placing the
sorted-descending digit $x_i$ at the $10^{\tau_i}$ place.

\textbf{The classical rule.} The classical Kaprekar rule $K_0$ at
digit length $d$ corresponds to $\pi = (d-1, d-2, \ldots, 1, 0)$
(descending) and $\sigma = (0, 1, \ldots, d-2, d-1)$ (ascending). At
$d = 4$, $K_0(n) = \alpha(n) - \beta(n)$ with
$\alpha(n) = \pi \cdot x(n) = \sum_i 10^{d-1-i} x_i$ arranging
sorted-descending digits in their own order, and
$\beta(n) = \sigma \cdot x(n) = \sum_i 10^i x_i$ arranging them in
reverse.

\textbf{Coefficient expansion.} Let $r(i) = \tau^{-1}(i)$ denote the
preimage of place $i$ under a permutation $\tau$ --- that is, the
sorted-descending position that contributes to the $10^i$ place of
$\tau \cdot x$. Then

\[\tau \cdot x \;=\; \sum_{i = 0}^{d-1} 10^i \, x_{\tau^{-1}(i)}.\]

Equivalently,

\[K_{\pi, \sigma}(n) \;=\; \left| \sum_{i = 0}^{d-1} c_i \, x_i \right|,\]

where the \textbf{coefficient vector} $c = (c_0, \ldots, c_{d-1})$ is
given by

\[c_i \;=\; 10^{\pi_i} - 10^{\sigma_i}.\]

\textbf{Sum-zero constraint.} Since $\{\pi_i\}_{i=0}^{d-1}$ and
$\{\sigma_i\}_{i=0}^{d-1}$ are each permutations of
$\{0, 1, \ldots, d-1\}$, we have
$\sum_i 10^{\pi_i} = \sum_i 10^{\sigma_i} = \frac{10^d - 1}{9}$.
Therefore $\sum_i c_i = 0$ identically. This will be important in §5.

\textbf{Three equivalent representations.} Throughout the paper we
freely switch between three representations of a permutation-pair rule,
and it is useful to see them side by side. We use the classical rule at
$d = 4$ as the running example.

Assign letters $a, b, c, d$ to the sorted-descending positions:
$a = x_0$ (the largest digit), $b = x_1$, $c = x_2$, $d = x_3$
(the smallest). A string of letters then denotes an integer: reading
left-to-right from the highest place to the lowest, each letter
specifies which sorted-descending position contributes that digit.

For the classical Kaprekar rule at $d = 4$, the three representations
are:

\begin{enumerate}
\def\labelenumi{(\roman{enumi})}
\item
  \textbf{Letter string:} $\;K_0(n) = \mathrm{abcd} - \mathrm{dcba}$.
\item
  \textbf{Permutation pair:} $\;\pi = (3, 2, 1, 0)$ and
  $\sigma = (0, 1, 2, 3)$.
\item
  \textbf{Coefficient vector:}
  $\;c = (c_0, c_1, c_2, c_3) = (999, \, 90, \, -90, \, -999)$.
\end{enumerate}

\textbf{Worked example.} Take $n = 6174$. Its sorted-descending form
is $x(n) = (x_0, x_1, x_2, x_3) = (7, 6, 4, 1)$. Using the
letter-string form, $\pi \cdot x = \mathrm{abcd} = 7641$ and
$\sigma \cdot x = \mathrm{dcba} = 1467$, so
$K_0(6174) = |7641 - 1467| = 6174$. Using the coefficient-vector form,
$K_0(6174) = 999 \cdot 7 + 90 \cdot 6 - 90 \cdot 4 - 999 \cdot 1 = 6993 + 540 - 360 - 999 = 6174$.
Both computations give the same answer, and both confirm that $6174$
is a fixed point of $K_0$ at $d = 4$.

The three representations encode the same information. The letter string
is the most legible: reading $\mathrm{abcd} - \mathrm{dcba}$ makes
immediately clear that the rule is ``digits in descending order minus
digits in ascending order,'' Kaprekar's original statement. The
permutation pair $(\pi, \sigma)$ is the formal definition used in
proofs. The coefficient vector $c = (c_0, \ldots, c_{d-1})$ with
$c_i = 10^{\pi_i} - 10^{\sigma_i}$ is the most useful for algebraic
manipulation, because it reduces the rule to a linear expression in the
sorted-descending digits.

For rules arising in §5 and §6 (the $60714$ lifting family and the
$6174$ cross-dimensional pattern), we will use whichever
representation is most convenient for the context, and Appendix B
provides a complete table of all three for $60714$'s ladder at each
$d \in \{5, 6, \ldots, 20\}$.

\subsection{Algebraic rank}\label{algebraic-rank}

The \textbf{algebraic rank} (or \textbf{surviving-variable count}) of a
rule $K_{\pi, \sigma}$ is

\[\mathrm{sv}(K_{\pi, \sigma}) \;=\; \bigl|\{\,i : c_i \neq 0\,\}\bigr| \;=\; \bigl|\{\,i : \pi_i \neq \sigma_i\,\}\bigr|.\]

The rule is \textbf{full-variable} at digit length $d$ if
$\mathrm{sv} = d$ --- equivalently, if the map
$i \mapsto (\pi_i, \sigma_i)$ has $\pi_i \neq \sigma_i$ for every
$i$.

\begin{proposition}[2.1 (derangement characterization of
full-variable)]
A rule $K_{\pi, \sigma}$ at digit length $d$
is full-variable if and only if $\sigma \circ \pi^{-1}$ is a
derangement of $\{0, 1, \ldots, d-1\}$.
\end{proposition}

\begin{proof}
$\mathrm{sv} = d$ means $\pi_i \neq \sigma_i$ for
all $i$. Equivalently, setting $j = \pi_i$ and
$\rho = \sigma \circ \pi^{-1}$, this says $\rho(j) \neq j$ for all
$j$ --- precisely that $\rho$ is a derangement.
\end{proof}

\textbf{Rule counts by rank.} The total number of ordered pairs
$(\pi, \sigma)$ with $\pi \neq \sigma$ is $d! \, (d! - 1)$. The
number of full-variable rules is $d! \, D_d$, where $D_d$ is the
number of derangements of $d$ elements. For small $d$:

\begin{longtable}[]{@{}ccc@{}}
\toprule\noalign{}
$d$ & $d! \, (d! - 1)$ (all rules) & $d! \, D_d$
(full-variable) \\
\midrule\noalign{}
\endhead
\bottomrule\noalign{}
\endlastfoot
2 & 2 & 2 \\
3 & 30 & 12 \\
4 & 552 & 216 \\
5 & 14\{,\}280 & 5\{,\}280 \\
6 & 517\{,\}680 & 190\{,\}800 \\
\end{longtable}

\subsection{Universality and fixed
points}\label{universality-and-fixed-points}

A \textbf{fixed point} of $K_{\pi, \sigma}$ at length $d$ is an
integer $F \in A_d$ with $K_{\pi, \sigma}(F) = F$. A fixed point
$F$ is \textbf{universal} for $K_{\pi, \sigma}$ at length $d$ if,
for every admissible input $n \in A_d$, the orbit
$n, K(n), K^2(n), \ldots$ reaches $F$ in finitely many steps. A rule
is \textbf{universal for $F$ at $d$} if it satisfies both
conditions.

\textbf{Example (classical).} The classical rule $K_0$ at $d = 4$ is
universal for $F = 6174$ \cite{Kaprekar1955}. It reaches $6174$
within at most seven iterations for every admissible input.

\textbf{Example (dimension-agnostic).} At $d = 3$, the classical rule
$K_0$ has coefficient vector $(99, \, 0, \, -99)$ --- the middle
coefficient vanishes by the forced borrow chain of §1.1 --- so
$\mathrm{sv}(K_0) = 2$. The classical rule at $d = 3$ is \emph{not}
full-variable. It is universal for $F = 495$, but $495$ is a fixed
point of a rank-$2$ rule, not a rank-$3$ rule.

\subsection{Native digit length}\label{native-digit-length}

For an integer $F$, its \textbf{native digit length} $d_F$ is the
smallest $d$ such that $F$ is a universal full-variable fixed point
of some rule at digit length $d$. If $F$ admits no universal
full-variable rule at any $d$, it has no native digit length in our
sense.

\textbf{Example.} At $d = 4$, the set of fixed points with native
digit length $4$ is $\{1746, 2538, 5382, 6174\}$ (see §3). Each is a
universal full-variable fixed point at $d = 4$ under some rule; for
$6174$, the rule is classical, but the other three arise from
non-classical permutation pairs.

\textbf{Dimension-agnostic fixed points.} Integers such as $45$
(universal at every $d \geq 2$ under low-rank rules derived from
classical reverse-pair structure), $495$ (universal at $d = 3$ under
$K_0$), and $450$ (universal at $d = 3$ under a cousin of the
classical rule) are \emph{not} in our full-variable classification
because their rules have $\mathrm{sv} < d$ at every $d$ where they
appear. These integers are fixed points of \emph{rank-reducing} rules,
where forced borrow chains or structural cancellations reduce the
effective input count below $d$. They persist across digit lengths
through this dimensional reduction.

Our full-variable classification excludes these by design: we study only
rules with $\mathrm{sv} = d$, where every sorted-descending position
contributes nontrivially. This restriction is principled rather than
definitional. Rank-reducing rules are those whose coefficient vector has
at least one zero --- equivalently, whose action depends on strictly
fewer than $d$ inputs. Such rules describe a $(d-k)$-dimensional
dynamical system embedded in a $d$-dimensional input space; their
fixed points and convergence behavior are determined by the
lower-dimensional projection, not by the digit-length $d$ at which the
rule happens to be expressed. Including them in the cross-dimensional
question would conflate persistence-by-projection (a structural
artifact) with persistence-by-genuine-extension (the phenomenon of
interest). The classical rule at $d = 5$ is itself such a rank-reducer
(sv $= 4$, not $5$), which is why classical Kaprekar fails at five
digits --- not because Kaprekar dynamics has no five-digit attractor,
but because the classical rule is rank-$4$ when viewed at $d = 5$.
The dimension-agnostic family ($45, 495, 450$, etc.) deserves
independent study but is outside the scope of this paper. See §1.1 for
the structural derivation of the middle-digit cancellation that produces
$495$'s rank-$2$ status.

\subsection{Effective rank at a fixed
point}\label{effective-rank-at-a-fixed-point}

For a fixed point $F$ with sorted-descending form
$(f_0, f_1, \ldots, f_{d-1})$ at length $d$, and for a rule
$K_{\pi, \sigma}$ with coefficient vector $c$, the \textbf{effective
rank at $F$} is

\[\mathrm{sv}_F(K_{\pi, \sigma}) \;=\; \bigl|\{\,i : c_i \neq 0 \text{ and } f_i \neq 0\,\}\bigr|.\]

Thus $\mathrm{sv}_F$ counts coefficients of the rule that \emph{land
on nonzero digits of $F$}. If $F$ has $z$ zero digits in its
sorted-descending form, then $\mathrm{sv}_F \leq d - z$ for any rule.

Effective rank plays a structural role in §4 and §5. Coefficients paired
with zero digits of $F$ vanish in the computation of $K(F)$, so the
fixed-point equation $K(F) = F$ is determined only by the coefficients
paired with nonzero digits. This is why coefficient-preserving lifting
(defined in §5) can introduce large new coefficients at positions
corresponding to $F$'s zero digits without disturbing $F$'s
fixed-point status.

\subsection{Basic properties}\label{basic-properties}

We collect two basic properties used without comment throughout the rest
of the paper.

\begin{proposition}[2.2 (sign-flip symmetry)]
If
$K_{\pi, \sigma}$ is universal for $F$ at length $d$, so is
$K_{\sigma, \pi}$, with the same basin. The coefficient vector of
$K_{\sigma, \pi}$ is the negation of that of $K_{\pi, \sigma}$.
\end{proposition}

\begin{proof}
The rule is
$K_{\pi, \sigma}(n) = |\pi \cdot x(n) - \sigma \cdot x(n)|$. Swapping
$\pi$ and $\sigma$ negates the argument of the absolute value, which
leaves the rule invariant.
\end{proof}

In consequence, universal rules come in sign-flipped pairs. Throughout
the paper, when we report ``the number of universal full-variable rules
for $F$,'' we count ordered pairs unless explicitly noted; the number
of geometrically distinct rules is half of this.

\begin{proposition}[2.3 (repdigit invariance)]
For any rule
$K_{\pi, \sigma}$ and any repdigit input $r$ (e.g., $11111$ at
$d = 5$), $K_{\pi, \sigma}(r) = 0$. Moreover,
$K_{\pi, \sigma}(0) = 0$. Thus $0$ is a fixed point of every rule.
\end{proposition}

\begin{proof}
If all digits of $r$ are equal to some value $v$,
then
$\pi \cdot x(r) = v \cdot \sum_i 10^{\pi_i} = v \cdot (10^d - 1)/9$,
and $\sigma \cdot x(r)$ has the same value. Their difference is zero.
\end{proof}

Proposition 2.3 is why we exclude repdigits from the admissible input
set: they provide no nontrivial dynamics and add a trivial fixed point
$0$ to every rule's fixed-point set. Near-repdigits, excluded for
similar reasons, can also produce trajectories that collapse to $0$
under coefficient-preserving liftings. The standard convention in
Kaprekar analysis is to restrict attention to ``generic'' inputs with at
least two distinct digit values, and we follow this.

\subsection{Summary of notation}\label{summary-of-notation}

For reference:

\begin{longtable}[]{@{}
  >{\raggedright\arraybackslash}p{(\columnwidth - 2\tabcolsep) * \real{0.5000}}
  >{\raggedright\arraybackslash}p{(\columnwidth - 2\tabcolsep) * \real{0.5000}}@{}}
\toprule\noalign{}
\begin{minipage}[b]{\linewidth}\raggedright
Symbol
\end{minipage} & \begin{minipage}[b]{\linewidth}\raggedright
Meaning
\end{minipage} \\
\midrule\noalign{}
\endhead
\bottomrule\noalign{}
\endlastfoot
$d$ & digit length \\
$n$ & a $d$-digit integer (possibly with leading zeros) \\
$x(n) = (x_0, \ldots, x_{d-1})$ & sorted-descending form of $n$,
with $x_0 \geq x_1 \geq \cdots$ \\
$A_d$ & admissible inputs at length $d$ (non-repdigit,
non-near-repdigit) \\
$\pi, \sigma$ & permutations in $S_d$ \\
$K_{\pi, \sigma}$ & permutation-pair rule \\
$K_0$ & classical rule (descending minus ascending) \\
$c_i = 10^{\pi_i} - 10^{\sigma_i}$ & coefficient at sorted-descending
position $i$ \\
$\mathrm{sv}(K)$ & algebraic rank (number of nonzero coefficients) \\
$\mathrm{sv}_F(K)$ & effective rank at $F$ (nonzero coefficients
paired with nonzero digits of $F$) \\
$F$ & a fixed point \\
$d_F$ & native digit length of $F$ \\
$(f_0, \ldots, f_{d-1})$ & sorted-descending form of $F$ at length
$d$ \\
$T_d$ & tail-two-zeros set: sorted-desc inputs at length $d$ with
$x_{d-1} = x_{d-2} = 0$ (§5.3) \\
$E_d$ & escape class: admissible inputs whose orbit reaches a
block-aligned multiset (§5.2) \\
$B_d$ & block-aligned multisets at length $d$ (Definition 5.2) \\
$d_0$ & base block size for the 60714 ladder: $5$ on the odd ladder,
$6$ on the even ladder \\
odd ladder & rules at odd $d \geq 5$ extending the $d = 5$ native
rule by zero-sum pair appending \\
even ladder & rules at even $d \geq 6$ extending the $d = 6$
split-lifted rule by zero-sum pair appending \\
$\mathrm{core}(x)$ & base-block contribution to $K^{(d)}(x)$:
$9900 x_0 + 9 x_1 + 90 x_2 - 9000 x_3 - 999 x_4$ (odd ladder) \\
$\delta_k$ & difference at the $k$-th appended pair:
$x_{2k+d_0} - x_{2k+d_0+1}$ (§C.4) \\
$\mathrm{sum\_locked\_spans}(K_F)$ & within-thread invariant:
$\sum_{i : f_i \neq 0} \lvert \pi_i - \sigma_i \rvert$ (§6.4) \\
coefficient-preserving lifting & a rule at $d + k$ whose coefficient
vector extends the $d$-coefficient vector by appending $k$ pairs
that sum to zero (§5.1) \\
dimension-transcendent & a fixed point with a coefficient-preserving
lifting at every $d \geq d_F$ \\
\end{longtable}

\section{\texorpdfstring{Classification of Universal Full-Variable Fixed
Points at
$d \leq 6$}{Classification of Universal Full-Variable Fixed Points at d \textbackslash leq 6}}\label{classification-of-universal-full-variable-fixed-points-at-d-leq-6}

This section presents the complete classification of universal
full-variable fixed points at digit lengths $d = 3, 4, 5, 6$. The
classification at each $d$ is obtained by exhaustive enumeration:
every full-variable rule is tested against every admissible input, and
the universal rules are identified. The counts are definitive; the full
list of fixed points at $d = 6$ appears in Appendix A.

\subsection{Method}\label{method}

At each digit length $d$, we enumerate the full-variable rules ---
ordered pairs $(\pi, \sigma) \in S_d \times S_d$ with
$\pi_i \neq \sigma_i$ for every $i$ (Proposition 2.1) --- and for
each such rule we check:

\begin{enumerate}
\def\labelenumi{\arabic{enumi}.}
\tightlist
\item
  Does the rule have at least one fixed point $F \in A_d$?
\item
  For each fixed point $F$, is the rule universal for $F$? That is,
  does every input $n \in A_d$ reach $F$ under iteration?
\end{enumerate}

A rule is \textbf{universal at $d$} if both conditions hold. The set
of \textbf{universal full-variable fixed points at $d$} is the set of
integers $F$ arising as the attractor of some universal full-variable
rule.

The enumerations are complete. At $d = 6$, the largest case we
present, $190{,}800$ full-variable rules are tested against
$4{,}905$ admissible digit multisets ($999{,}450$ admissible padded
six-digit strings), a search that runs in hours on commodity hardware
with the iteration-and-caching implementation described in Appendix A.

\subsection{\texorpdfstring{$d = 3$: no universal full-variable fixed
points}{d = 3: no universal full-variable fixed points}}\label{d-3-no-universal-full-variable-fixed-points}

At $d = 3$, there are $12$ full-variable rules (Proposition 2.1,
$3! \cdot D_3 = 6 \cdot 2 = 12$).

\begin{theorem}[3.1]
No full-variable rule at $d = 3$ is
universal for any fixed point.
\end{theorem}

\textbf{Proof sketch.} Direct enumeration. For each of the $12$ rules,
the iteration from any admissible $d = 3$ input enters a cycle of
length $\geq 2$ or a non-universal fixed point. No rule admits a fixed
point with basin covering all of $A_3$. Full computation in Appendix
A.1.

\begin{remark}[on $495$]
The classical attractor $495$ does not
appear here because the classical rule $K_0$ at $d = 3$ is not
full-variable: its middle coefficient vanishes by the forced borrow
chain (§1.1), giving $\mathrm{sv}(K_0) = 2$. $495$ is a fixed point
of a rank-$2$ rule, outside the scope of Theorem 3.1. It reappears in
Appendix §2.5 as a dimension-agnostic fixed point.
\end{remark}

\subsection{\texorpdfstring{$d = 4$: four universal full-variable
fixed
points}{d = 4: four universal full-variable fixed points}}\label{d-4-four-universal-full-variable-fixed-points}

At $d = 4$, there are $216$ full-variable rules (Proposition 2.1,
$4! \cdot D_4 = 24 \cdot 9 = 216$).

\begin{theorem}[3.2]
There are exactly four universal
full-variable fixed points at $d = 4$:
\end{theorem}

\[F_4 = \{\, 1746,\; 2538,\; 5382,\; 6174 \,\}.\]

\emph{These fixed points are reached by exactly $8$ universal
full-variable rules in total (each fixed point reached by exactly $2$
universal rules, which are sign-flips of each other per Proposition
2.2).}

\begin{proof}
Exhaustive enumeration of $216$ rules against $615$
admissible digit multisets ($9{,}630$ admissible padded four-digit
strings). For each universal rule, basin coverage of $A_4$ is verified
directly. Full computation in Appendix A.2.

\textbf{Structure of the classification.} All four universal fixed
points share the digit-sum property $\sum_i f_i \equiv 0 \pmod 9$. The
integers $6174$ and $1746$ are digit-permutations of each other
(both have multiset $\{1, 4, 6, 7\}$); $2538$ and $5382$ are
digit-permutations of each other (multiset $\{2, 3, 5, 8\}$). These
anagram clusters will recur in §3.4 and §6.

\begin{observation}[3.1]
The sum of the two fixed points in each
anagram cluster is $7920$:
\end{observation}

\[6174 + 1746 = 7920, \qquad 2538 + 5382 = 7920.\]

This identity is a specific instance of reverse-pair structure at
$d = 4$: if $F$ has digits summing to $18$, then
$F + \mathrm{reverse}(F)$ takes a predictable form. We do not develop
this further here; see \cite{Nuez2021} for parametric analysis of this
phenomenon.

\textbf{Connection to the classical rule.} $K_0$ at $d = 4$ is
full-variable (coefficient vector $(999, 90, -90, -999)$), and $K_0$
is one of the $2$ universal rules for $6174$. The classical Kaprekar
theorem is thus recovered within our classification as a specific entry
in Theorem 3.2.

\subsection{\texorpdfstring{$d = 5$: thirty-three universal
full-variable fixed
points}{d = 5: thirty-three universal full-variable fixed points}}\label{d-5-thirty-three-universal-full-variable-fixed-points}

At $d = 5$, there are $5{,}280$ full-variable rules (Proposition
2.1, $5! \cdot D_5 = 120 \cdot 44 = 5{,}280$).

\begin{theorem}[3.3]
There are exactly $33$ universal
full-variable fixed points at $d = 5$, reached collectively by $66$
universal full-variable rules. Each fixed point is reached by exactly
$2$ universal rules (sign-flip pairs).
\end{theorem}

\textbf{Proof.} Exhaustive enumeration of $5{,}280$ rules against
$1{,}902$ admissible digit multisets ($99{,}540$ admissible padded
five-digit strings). Full computation in Appendix A.3; the complete list
of $33$ fixed points appears in Table 3.1 below.

\textbf{Classical rule does not appear.} The classical rule $K_0$ at
$d = 5$ has coefficient vector
$(9999, \, 990, \, 0, \, -990, \, -9999)$ --- the middle coefficient
vanishes by the forced borrow chain at odd digit length (§1.1). $K_0$
at $d = 5$ has $\mathrm{sv} = 4$, not $5$; it is not in the
full-variable classification. The observation that ``the classical
Kaprekar routine fails at $d = 5$'' is, in our framing, the
observation that the classical rule at $d = 5$ is not even in the
full-variable space --- there is no contradiction to the existence of
$33$ universal attractors in that space.

\textbf{Reverse-pair structure fails at odd $d$.} The classical rule
is a specific ``reverse-pair'' rule ($\sigma_i = \pi_{d-1-i}$). At
even $d$, reverse-pair rules can be full-variable; at odd $d$, they
cannot --- the middle position always satisfies
$\sigma_{(d-1)/2} = \pi_{(d-1)/2}$, forcing
$\mathrm{sv} \leq d - 1$. Consequently, all $66$ universal
full-variable rules at $d = 5$ are \emph{non-reverse-pair} --- they
have no mirror symmetry. This is a structural distinction from the
classical case.

\textbf{Table 3.1.} \emph{The $33$ universal full-variable fixed
points at $d = 5$, grouped by digit multiset.}

\begin{longtable}[]{@{}
  >{\raggedright\arraybackslash}p{(\columnwidth - 4\tabcolsep) * \real{0.3077}}
  >{\raggedright\arraybackslash}p{(\columnwidth - 4\tabcolsep) * \real{0.3077}}
  >{\centering\arraybackslash}p{(\columnwidth - 4\tabcolsep) * \real{0.3846}}@{}}
\toprule\noalign{}
\begin{minipage}[b]{\linewidth}\raggedright
Multiset
\end{minipage} & \begin{minipage}[b]{\linewidth}\raggedright
Fixed points
\end{minipage} & \begin{minipage}[b]{\linewidth}\centering
Cluster size
\end{minipage} \\
\midrule\noalign{}
\endhead
\bottomrule\noalign{}
\endlastfoot
$\{0, 0, 0, 4, 5\}$ & $54$ & $1$ \\
$\{0, 3, 3, 5, 7\}$ & $3753$ & $1$ \\
$\{0, 1, 4, 6, 7\}$ & $60714$, $60417$ & $2$ \\
$\{1, 2, 3, 4, 8\}$ & $18342$, $21834$, $24183$, $41832$,
$42183$ & $5$ \\
$\{2, 3, 5, 8, 9\}$ & $28539$, $53928$, $58239$ & $3$ \\
$\{1, 5, 6, 7, 8\}$ & $16578$, $16758$, $17685$, $65781$,
$67581$ & $5$ \\
$\{3, 4, 5, 7, 8\}$ & $37584$, $37854$, $38754$, $43758$,
$43785$, $43875$ & $6$ \\
$\{1, 2, 4, 5, 6\}$ & $12456$, $14562$, $15642$, $16524$,
$21456$, $24156$, $24561$, $41562$, $45612$, $45621$ &
$10$ \\
\end{longtable}

\emph{Total: $33$ fixed points across $8$ multiset clusters. Each
fixed point is universal for exactly $2$ rules (sign-flip pair),
giving $66$ universal rules total.}

\begin{observation}[3.2]
The $\{0, 1, 4, 6, 7\}$ multiset
produces exactly two universal fixed points at $d = 5$: $60714$ and
$60417$. These are the multiset twins featured in §6. Despite
identical digit content, their native rules differ in a specific
structural invariant (developed in §6) that predicts their contrasting
cross-dimensional behavior.
\end{observation}

\begin{observation}[3.3]
The classical fixed points $6174$,
$1746$, $2538$, $5382$ (from $d = 4$) do not appear at $d = 5$
even as fixed points of any full-variable rule --- not merely as
non-universal fixed points, but as non-fixed points. Exhaustive
enumeration confirms: for each of these four integers padded to five
digits, no full-variable rule at $d = 5$ satisfies
$K(F_{\text{padded}}) = F$.
\end{observation}

\begin{observation}[3.4 (effective rank within the full-variable
classification)]
Two of the $33$ fixed points in Table 3.1 ---
namely $54$ and $3753$ --- have effective rank at $F$ strictly
less than their native algebraic rank:
\end{observation}

\begin{itemize}
\item
  \emph{$54$ has sorted-descending form $(5, 4, 0, 0, 0)$ at
  $d = 5$; three of the five digits are zero. Its native sv$=5$ rule
  is $\pi = (1, 2, 4, 3, 0), \sigma = (2, 0, 3, 4, 1)$ with
  coefficient vector $(-90, 99, 9000, -9000, -9)$. Three of these five
  coefficients are paired with zero digits in the fixed-point equation,
  giving $\mathrm{sv}_F(K_{54}) = 2 < 5 = \mathrm{sv}(K_{54})$.}
\item
  \emph{$3753$ has sorted-descending form $(7, 5, 3, 3, 0)$ at
  $d = 5$; one digit is zero. Its native rule has
  $\mathrm{sv}_F = 4 < 5$.}
\end{itemize}

\emph{These fixed points are genuinely in the full-variable
classification --- their native rules are algebraically rank-$5$ ---
but their effective rank at $F$ is reduced by absorption of
coefficients at zero-digit positions. This is the mechanism that
underlies coefficient-preserving lifting in §5, visible here at the root
of the classification: a rule can be algebraically full-variable while
being effectively lower-rank on a specific fixed point. The remaining
$31$ fps in Table 3.1 all have no zero digits and therefore have
$\mathrm{sv}_F = \mathrm{sv} = 5$.}

\subsection{\texorpdfstring{$d = 6$: five hundred and six universal
full-variable fixed
points}{d = 6: five hundred and six universal full-variable fixed points}}\label{d-6-five-hundred-and-six-universal-full-variable-fixed-points}

At $d = 6$, there are $190{,}800$ full-variable rules (Proposition
2.1, $6! \cdot D_6 = 720 \cdot 265$).

\begin{theorem}[3.4]
There are exactly $506$ universal
full-variable fixed points at $d = 6$, reached collectively by
$1{,}174$ universal full-variable rules.
\end{theorem}

\textbf{Proof.} Exhaustive enumeration of $190{,}800$ full-variable
rules against $4{,}905$ admissible digit multisets ($999{,}450$
admissible padded six-digit strings). Full enumeration data and the
complete list of $506$ fixed points are in Appendix A.4.

\begin{remark}[on enumeration soundness]
The reproducibility script
\texttt{classify\_at\_d.py} uses a candidate-filtering heuristic
followed by full-basin verification of each candidate, which scales
tractably to $d = 6$ but is not by itself a sound enumeration (a
candidate fp could in principle be missed by the filter). The $506$
count was independently cross-verified against three sources that all
agree: (i) a 3-day exhaustive enumeration on consumer hardware producing
\texttt{results\_db.json}, (ii) a no-shortcut sound enumeration on the
verifier container, and (iii) the supplementary file
\texttt{d6\_fps.txt}, whose stratification by zero-digit count
$(205 + 240 + 53 + 8 = 506)$ matches the table below. The three
sources independently arrive at the same fixed-point set and the same
count.
\end{remark}

\textbf{Stratification by zero digit count.} The $506$ fixed points
distribute as follows by the number of zero digits in each fixed point:

\begin{longtable}[]{@{}cc@{}}
\toprule\noalign{}
Zero-digit count & Fixed-point count \\
\midrule\noalign{}
\endhead
\bottomrule\noalign{}
\endlastfoot
$0$ & $205$ \\
$1$ & $240$ \\
$2$ & $53$ \\
$3$ & $8$ \\
\textbf{Total} & \textbf{$506$} \\
\end{longtable}

This stratification matters for §5 and §6: fixed points with more zero
digits admit more coefficient-preserving liftings to higher $d$, a
phenomenon visible already at the $d = 5 \to d = 6$ boundary and
developed in detail at §6.

\textbf{Note on the high-zero-count strata.} The $8$ fps with $3$
zero digits at $d = 6$ all share digit multiset
$\{0, 0, 0, 2, 2, 5\}$: they are
$\{252, 2520, 20025, 25200, 200025, 200250, 250200, 252000\}$. There
are no universal full-variable fps at $d = 6$ with $4$ or more zero
digits --- the trivial fixed point $F = 0$ is excluded throughout this
paper by the convention of §2 (Proposition 2.3). In particular, the
fixed point $549{,}945$ has digit multiset $\{4, 4, 5, 5, 9, 9\}$
with no zero digits and does not appear in this stratification.

\textbf{Connection to classical Kaprekar at $d = 6$.} The classical
rule $K_0$ at $d = 6$ has coefficient vector
$(999{,}999, \, 89{,}991, \, 8{,}991, \, -8{,}991, \, -89{,}991, \, -999{,}999)$
after simplification. It is full-variable at $d = 6$ (all coefficients
nonzero). \cite{Dahl2026} analyzes its basin: the rule is \emph{not}
universal at $d = 6$ --- it produces a 7-cycle that attracts
$93.55\%$ of inputs, with $6.25\%$ reaching $F = 631{,}764$ and
$0.20\%$ reaching $F = 549{,}945$. The universal fixed points at
$d = 6$ in our classification are attractors of \emph{other}
full-variable rules, not the classical one.

\begin{observation}[3.4]
The digit sum of every universal
full-variable fixed point at $d = 6$ is divisible by $9$. The
distribution of digit sums is strikingly bell-shaped about $27$ (the
midpoint of the admissible range $\{9, 18, 27, 36, 45\}$):
\end{observation}

\begin{longtable}[]{@{}cc@{}}
\toprule\noalign{}
Digit sum & Fixed-point count \\
\midrule\noalign{}
\endhead
\bottomrule\noalign{}
\endlastfoot
$9$ & $8$ \\
$18$ & $156$ \\
$27$ & $244$ \\
$36$ & $96$ \\
$45$ & $2$ \\
\textbf{Total} & \textbf{$506$} \\
\end{longtable}

This pattern reflects a general arithmetic fact about universal
Kaprekar-type attractors: $K(n) \equiv 0 \pmod 9$ whenever
$n \in A_d$, so fixed points satisfy $F \equiv 0 \pmod 9$. The
bell-shape reflects the combinatorics of digit-multisets satisfying this
modular constraint.

\subsection{Summary}\label{summary}

The classifications at $d = 3, 4, 5, 6$ are:

\begin{longtable}[]{@{}cccc@{}}
\toprule\noalign{}
$d$ & Full-variable rules & Universal fps & Universal rules \\
\midrule\noalign{}
\endhead
\bottomrule\noalign{}
\endlastfoot
$3$ & $12$ & $0$ & $0$ \\
$4$ & $216$ & $4$ & $8$ \\
$5$ & $5{,}280$ & $33$ & $66$ \\
$6$ & $190{,}800$ & $506$ & $1{,}174$ \\
\end{longtable}

Each classification is complete and reproducible; full enumeration
scripts and data are in Appendix A.

\textbf{Why the count grows so quickly.} The growth in universal
fixed-point count from $0$ at $d = 3$ to $506$ at $d = 6$
reflects two phenomena: (i) the space of full-variable rules grows
factorially in $d$, providing more opportunities for fixed points, and
(ii) integers with more digits have more admissible multisets, more
potential rearrangements, and so more potential fixed-point equations to
satisfy. The question of which of these $506$ fixed points at
$d = 6$ arise as coefficient-preserving liftings from lower $d$, and
which are genuinely new at $d = 6$, is the content of §4.

\section{\texorpdfstring{Cross-Dimensional Cross-Check:
$d = 5 \to d = 6$}{Cross-Dimensional Cross-Check: d = 5 \textbackslash to d = 6}}\label{cross-dimensional-cross-check-d-5-to-d-6}

§3 established the classifications at $d = 3, 4, 5, 6$ independently
at each digit length. This section addresses the cross-dimensional
question of §1.3: among the $33$ universal full-variable fixed points
at $d = 5$, which extend to $d = 6$ as universal full-variable fixed
points?

\subsection{The question, formally}\label{the-question-formally}

Let $F_5 = \{F^{(1)}, \ldots, F^{(33)}\}$ be the set of universal
full-variable fixed points at $d = 5$ (Theorem 3.3). For each
$F^{(k)} \in F_5$, consider the padded integer $F^{(k)}_{(6)}$ ---
that is, $F^{(k)}$ interpreted as a six-digit integer by prepending a
zero. Let $\overline{F^{(k)}}$ denote its sorted-descending form at
length $6$: this is the sorted-descending form at length $5$ with a
zero appended, since padding adds a zero that sorts to the smallest
position.

We ask two questions for each $F^{(k)}$:

\textbf{Question A (algebraic).} Does any full-variable rule
$K_{\pi, \sigma}$ at $d = 6$ satisfy $K(F^{(k)}_{(6)}) = F^{(k)}$?
That is, does the fixed-point equation admit any full-variable solution
at $d = 6$?

\textbf{Question B (dynamic).} If yes, is any such rule universal for
$F^{(k)}$ at $d = 6$ --- meaning every input in $A_6$ iterates to
$F^{(k)}$?

A positive answer to both questions constitutes
\textbf{coefficient-lifting success} from $d = 5$ to $d = 6$: the
fixed point $F^{(k)}$ is a universal full-variable fixed point at both
digit lengths. A negative answer to Question A is \textbf{algebraic
obstruction} (no rule even fixes $F^{(k)}$ at $d = 6$). A negative
answer to Question B despite a positive answer to A is \textbf{dynamic
obstruction} (rules fix $F^{(k)}$ but none is universal).

\subsection{The result}\label{the-result}

\begin{theorem}[4.1]
Among the $33$ universal full-variable
fixed points at $d = 5$, the cross-dimensional behavior at $d = 6$
is as follows:
\end{theorem}

\begin{itemize}
\tightlist
\item
  \textbf{$17$ fixed points face algebraic obstruction:} no
  full-variable rule at $d = 6$ fixes $F$.
\item
  \textbf{$15$ fixed points face dynamic obstruction:} full-variable
  rules at $d = 6$ fix $F$, but none is universal.
\item
  \textbf{$1$ fixed point admits a universal coefficient lifting:}
  $F = 60714$. Of the four full-variable rules at $d = 6$ satisfying
  the fixed-point equation $K(60714) = 60714$, exactly two are
  dynamically universal (basin $= 1$ over admissible inputs), and
  these two form a sign-flip pair.
\end{itemize}

\emph{Formally: the set of $d = 5$ universal full-variable fixed
points that are also $d = 6$ universal full-variable fixed points is
$\{60714\}$. All other $32$ $d = 5$ fixed points are
dimension-locked at $d = 5$.}

\textbf{Proof.} By exhaustive computation in two parts.

\emph{Algebraic part.} For each $F^{(k)} \in F_5$, we enumerate all
full-variable rules $K_{\pi, \sigma}$ at $d = 6$ and test whether
$K(F^{(k)}_{(6)}) = F^{(k)}$. This reduces to a solution-counting
problem: writing $\overline{F^{(k)}} = (f_0, \ldots, f_5)$, we seek
$(\pi, \sigma) \in S_6 \times S_6$ with $\pi_i \neq \sigma_i$ for
every $i$ and
$\sum_i (10^{\pi_i} - 10^{\sigma_i}) f_i = \pm F^{(k)}$. Since
$F^{(k)} \in F_5$ is nonzero, the equation is nontrivial. For each
$F^{(k)}$, the number of full-variable $(\pi, \sigma)$ satisfying
the equation is a specific non-negative integer. Computing this count
for all $33$ fixed points:

\begin{itemize}
\tightlist
\item
  $17$ fixed points admit zero full-variable solutions.
\item
  $16$ fixed points admit at least one full-variable solution. The
  number of full-variable rules satisfying the fixed-point equation per
  fixed point varies from $4$ (for $60714$ and $60417$) up to
  $528$ (for $F = 54$, with three zero digits when padded). Among
  these algebraic solutions, only some are \emph{dynamically} universal.
\end{itemize}

\emph{Dynamic part.} For each of the $16$ fixed points with at least
one algebraic solution, we test each fixed-point-satisfying rule for
universality: does every admissible input $n \in A_6$ iterate to
$F^{(k)}$ under this rule? This test is direct: iterate from each
admissible digit multiset at $d = 6$ ($4{,}905$ multisets) for up to
$300$ steps, and record whether $F^{(k)}$ is reached. The basin is
then the fraction of inputs reaching $F^{(k)}$.

For $15$ of the $16$ fixed points with algebraic solutions, every
fixed-point-satisfying rule has basin $< 1$. The best basins are
tabulated below.

For $1$ fixed point --- $F = 60714$ --- exactly two of the four
full-variable rules satisfying the fixed-point equation are universal
(basin $= 1.0$); the other two have basin substantially less than
$1$.

The computation is implemented as a two-stage enumeration: first a fast
enumeration over $(\pi, \sigma) \in S_6 \times S_6$ with derangement
filtering, producing the algebraic-solution set; then a basin test on
each surviving candidate. Total runtime on commodity hardware: a few
minutes for the algebraic part, several hours for the dynamic part. Full
enumeration logs and per-fixed-point data are in Appendix A.4.
\end{proof}

\subsection{Tabulation of the 32 dimension-locked fixed
points}\label{tabulation-of-the-32-dimension-locked-fixed-points}

\textbf{Table 4.1.} \emph{The $17$ fixed points at $d = 5$ with
algebraic obstruction at $d = 6$.}

\begin{longtable}[]{@{}cl@{}}
\toprule\noalign{}
Fixed point & Digit multiset at $d = 5$ \\
\midrule\noalign{}
\endhead
\bottomrule\noalign{}
\endlastfoot
$12456$ & $\{1, 2, 4, 5, 6\}$ \\
$14562$ & $\{1, 2, 4, 5, 6\}$ \\
$15642$ & $\{1, 2, 4, 5, 6\}$ \\
$16524$ & $\{1, 2, 4, 5, 6\}$ \\
$21456$ & $\{1, 2, 4, 5, 6\}$ \\
$24156$ & $\{1, 2, 4, 5, 6\}$ \\
$24561$ & $\{1, 2, 4, 5, 6\}$ \\
$28539$ & $\{2, 3, 5, 8, 9\}$ \\
$38754$ & $\{3, 4, 5, 7, 8\}$ \\
$41832$ & $\{1, 2, 3, 4, 8\}$ \\
$42183$ & $\{1, 2, 3, 4, 8\}$ \\
$43758$ & $\{3, 4, 5, 7, 8\}$ \\
$43785$ & $\{3, 4, 5, 7, 8\}$ \\
$43875$ & $\{3, 4, 5, 7, 8\}$ \\
$53928$ & $\{2, 3, 5, 8, 9\}$ \\
$65781$ & $\{1, 5, 6, 7, 8\}$ \\
$67581$ & $\{1, 5, 6, 7, 8\}$ \\
\end{longtable}

\emph{All $17$ fixed points share a structural feature: their
multisets at $d = 5$ contain at most one zero digit, so padding to
$d = 6$ gives at most two absorbing positions --- typically
insufficient for any full-variable sv$=6$ rule to satisfy
$K(F_{(6)}) = F_{(6)}$.}

\textbf{Table 4.2.} \emph{The $15$ fixed points at $d = 5$ with
dynamic obstruction at $d = 6$.}

For these $15$ fixed points, full-variable rules satisfying
$K(F_{(6)}) = F_{(6)}$ exist but are not universal. The best basin
each fixed point achieves at $d = 6$ over the $4{,}905$ admissible
multisets is given below.

\begin{longtable}[]{@{}ccc@{}}
\toprule\noalign{}
Fixed point & \# candidate sv=6 rules & Best basin \\
\midrule\noalign{}
\endhead
\bottomrule\noalign{}
\endlastfoot
$54$ & $528$ & $0.9631$ \\
$60417$ & $4$ & $0.9598$ \\
$21834$ & $4$ & $0.9472$ \\
$45621$ & $12$ & $0.7951$ \\
$24183$ & $4$ & $0.7182$ \\
$37584$ & $2$ & $0.5594$ \\
$45612$ & $12$ & $0.5111$ \\
$41562$ & $4$ & $0.5025$ \\
$3753$ & $16$ & $0.3199$ \\
$18342$ & $4$ & $0.2137$ \\
$16578$ & $2$ & $0.1953$ \\
$58239$ & $4$ & $0.1880$ \\
$16758$ & $2$ & $0.0520$ \\
$37854$ & $2$ & $0.0245$ \\
$17685$ & $2$ & $0.0049$ \\
\end{longtable}

\emph{The basins range from $0.0049$ to $0.9631$; no rule is
universal. For most fixed points, competing attractors (the rule's other
fixed points or short cycles) draw off substantial fractions of input
space. Notably, $54$ and $60417$ achieve basins very close to $1$
(respectively $0.9631$ and $0.9598$) but remain dynamically
obstructed --- they are ``near misses'' in a precise sense, and were
among the candidates investigated in detail before $60714$ was
identified as the unique transcendent fp at $d = 5$.}

\begin{observation}[4.1 (zero-digit count predicts algebraic solvability
at $d = 6$)]
Among the $33$ fixed points at $d = 5$, those
with one zero digit at $d = 5$ (so two zero digits when padded to
$d = 6$) are candidates for coefficient-preserving lifting with two
absorbing positions. Those with no zero digits have only one absorbing
position after padding, and empirically none admits a full-variable
sv=$6$ rule fixing them. This pattern --- more zero digits
$\Rightarrow$ more absorbing capacity $\Rightarrow$ more algebraic
solutions --- is developed further at §5 and §6.
\end{observation}

\subsection{\texorpdfstring{The unique success case:
$F = 60714$}{The unique success case: F = 60714}}\label{the-unique-success-case-f-60714}

We record the $d = 6$ universal lifting of $60714$ explicitly, as it
is the cornerstone of §5.

\begin{proposition}[4.1]
At $d = 6$, exactly $2$
full-variable rules are universal for $F = 60714$. These are sign-flip
pairs of each other. One canonical representative has permutations
\end{proposition}

\[\pi = (4, 1, 2, 3, 5, 0), \qquad \sigma = (2, 0, 1, 4, 3, 5),\]

\emph{with coefficient vector}

\[c = (9900,\; 9,\; 90,\; -9000,\; 99000,\; -99999).\]

\emph{In letter notation (§ Appendix B conventions), the rule is}

\[K(n) = \mathrm{eadcbf}(n) - \mathrm{fdeacb}(n).\]

\emph{The rule is universal at $d = 6$: every admissible input
$n \in A_6$ iterates to $60714$ under $K$, with basin $= 1.0$
verified over all $4{,}905$ admissible digit multisets ($999{,}450$
admissible padded six-digit strings).}

\textbf{Verification.} Apply $K$ to $F_{(6)} = 060714$.
Sorted-descending form: $(7, 6, 4, 1, 0, 0)$. Compute:

\[K(060714) = |eadcbf - fdeacb|_{a=7, b=6, c=4, d=1, e=0, f=0} = |071460 - 010746| = 60714. \quad \checkmark\]

Universality is verified by exhaustive iteration over $A_6$: every
admissible six-digit integer reaches $60714$ in at most some number of
steps, with maximum reaching time bounded for all inputs (direct fixing)
and up to a small number of steps for the remainder. See Appendix A.4
for the full basin-verification log. $\square$

\textbf{The structural shape of the lifting.} Compare $60714$'s native
rule at $d = 5$:

\[\pi_5 = (4, 1, 2, 3, 0), \qquad \sigma_5 = (2, 0, 1, 4, 3),\]

coefficient vector $(9900, 9, 90, -9000, -999)$.

The $d = 6$ rule does not simply append a position; it restructures.
Specifically: the $d = 5$ rule's fifth coefficient is $-999$; the
$d = 6$ rule's fifth and sixth coefficients are $99000$ and
$-99999$, which sum to $99000 - 99999 = -999$. This is the
\textbf{split lifting} structure developed at §5: the single native
coefficient $-999$ is ``split'' across two new positions at $d = 6$,
summing back to $-999$ to preserve the action of the rule on $F$'s
nonzero digits (since both new positions correspond to zero digits of
$F_{(6)}$, which absorb them in the $K(F) = F$ computation).

This lifting structure is explicit, and its generalization to
$d = 7, 8, 9, \ldots$ is the content of §5.

\subsection{Implications for the
paper}\label{implications-for-the-paper}

Theorem 4.1 is the central cross-dimensional result at the
$d = 5 \to d = 6$ boundary. It establishes:

\begin{enumerate}
\def\labelenumi{\arabic{enumi}.}
\tightlist
\item
  \textbf{Cross-dimensional universality is rare at this boundary.} Only
  $1$ of $33$ fixed points succeeds.
\item
  \textbf{The distinction is structural, not incidental.} $60714$'s
  native rule has a specific algebraic structure --- the split-lifting
  structure --- that admits extension to $d = 6$. The other $32$
  fixed points' native rules do not.
\item
  \textbf{The mechanism generalizes.} The coefficient-preserving lifting
  that works at $d = 5 \to d = 6$ can be iterated to
  $d = 6 \to d = 7 \to d = 8 \to \cdots$ for $60714$. This is the
  content of §5.
\item
  \textbf{The uniqueness of $60714$ is proven.} No other $d = 5$
  universal full-variable fixed point extends to $d = 6$. Future
  results about $60714$'s cross-dimensional behavior rest on this
  uniqueness, which is established here by exhaustion.
\end{enumerate}

\begin{remark}[on Conjecture 1 (informal statement)]
For a universal
full-variable fixed point $F$ at native digit length $d_F$, define
its \emph{dimension-locking spectrum} as the set of $d > d_F$ at which
$F$ is also a universal full-variable fixed point. Theorem 4.1
establishes that $32$ of $33$ fixed points at $d = 5$ have empty
dimension-locking spectrum extension to $d = 6$, and that one ---
$60714$ --- has $d = 6$ in its spectrum. Theorem 5.2 of the next
section extends this: $60714$'s dimension-locking spectrum contains
every $d \geq 5$. The other $32$ fixed points' spectra at $d > 6$
are addressed individually at the appendix level, but no systematic
characterization of which fixed points have unbounded dimension-locking
spectra is presently known beyond the case of $60714$.
\end{remark}

A formal conjecture --- that every full-variable fixed point with
unbounded dimension-locking spectrum admits a coefficient-preserving
lifting similar to $60714$'s --- is stated as Conjecture 7.1.

\section{\texorpdfstring{Dimension-Transcendence of
$F = 60714$}{Dimension-Transcendence of F = 60714}}\label{dimension-transcendence-of-f-60714}

This section develops the coefficient-preserving lifting framework and
proves the paper's main theorem (Theorem 5.2): $F = 60714$ is a fixed
point of an explicit coefficient-preserving lifting at every digit
length $d \geq 5$, with universal convergence on $A_d \setminus E_d$
at every $d \geq 5$, where the escape class $E_d$ is empty at
$d = 5, 6$ and at $d \geq 7$ is structurally continuous with the
classical Kaprekar escape phenomenon (the ten repdigits at $d = 3$,
the ten multiples of $1111$ at $d = 4$).

The section is structured as follows. §5.1 defines
coefficient-preserving liftings and establishes that such liftings
automatically satisfy the fixed-point equation (Proposition 5.1). §5.2
specializes to the zero-sum pair construction used for $60714$,
introduces the two ladders (odd and even), and states the main theorem.
§5.3 reduces the theorem to two lemmas: Lemma 5.1 (structural closure of
the tail-two-zeros set) and Lemma 5.2 (bounded reaching time). §5.4
proves Lemma 5.1. §5.5 proves Lemma 5.2 in two parts: finite-state
enumeration at $d \leq 16$, and a $d$-independent algebraic argument
(one-step $T_d$ closure at $d \geq 15$ on the odd ladder; bounded
reaching time at $d \geq 16$ on the even ladder). §5.6 addresses
near-repdigit exclusion. §5.7 completes the proof of the main theorem.

\subsection{Coefficient-preserving
liftings}\label{coefficient-preserving-liftings}

Throughout this section, $F$ denotes a universal full-variable fixed
point at some native digit length $d_F$, and $K_{F}^{(d_{F})}$
denotes a specific universal rule at $d_F$ with coefficient vector
$c^{(d_{F})} = (c_0, c_1, \ldots, c_{d_{F} - 1})$. Let
$f^{(d_{F})} = (f_0, \ldots, f_{d_{F} - 1})$ be the sorted-descending
form of $F$ at $d_F$, and let $Z = \{i : f_i = 0\}$ be the set of
zero-digit positions,
$N = \{i : f_i \neq 0\} = \{0, 1, \ldots, d_F - 1\} \setminus Z$ the
nonzero-digit positions.

For any $d > d_F$, write $F_{(d)}$ for $F$ padded to $d$ digits.
The sorted-descending form of $F_{(d)}$ is $f^{(d_{F})}$ with
$d - d_F$ zeros appended (since padding adds zeros, which sort to the
smallest positions). Thus
$f^{(d)} = (f_0, \ldots, f_{d_{F} - 1}, 0, 0, \ldots, 0)$ with
$d - d_F$ trailing zeros.

\begin{definition}[5.1 (coefficient-preserving lifting)]
A rule
$K^{(d)}$ at digit length $d > d_F$ with coefficient vector
$c^{(d)} = (\tilde c_0, \tilde c_1, \ldots, \tilde c_{d-1})$ is a
\textbf{coefficient-preserving lifting} of $K_{F}^{(d_{F})}$ at length
$d$ if
\end{definition}

\[\tilde c_i = c_i \quad \text{for every } i \in N.\]

\emph{(That is, the coefficients at $F$'s nonzero-digit positions are
preserved verbatim from the native rule.)}

\begin{proposition}[5.1 (fixed-point equation preservation)]
If
$K^{(d)}$ is a coefficient-preserving lifting of $K_{F}^{(d_{F})}$,
then $K^{(d)}(F_{(d)}) = F$.
\end{proposition}

\begin{proof}
By the coefficient expansion of §2.2,

\[K^{(d)}(F_{(d)}) = \left| \sum_{i = 0}^{d - 1} \tilde c_i \cdot f^{(d)}_i \right|.\]

For $i \in N$, $\tilde c_i = c_i$ and $f^{(d)}_i = f_i$. For
$i \in Z \cap \{0, \ldots, d_F - 1\}$, $f^{(d)}_i = 0$, so
$\tilde c_i \cdot 0 = 0$. For $i \in \{d_F, \ldots, d - 1\}$,
$f^{(d)}_i = 0$ (the appended-zero positions), so again
$\tilde c_i \cdot 0 = 0$. Therefore

\[K^{(d)}(F_{(d)}) = \left| \sum_{i \in N} c_i \cdot f_i \right| = \left| K^{(d_{F})}(F) \right| = F,\]

where the final equality uses that $K_{F}^{(d_{F})}$ fixes $F$ at
$d_F$.
\end{proof}

\begin{remark}[5.1]
Proposition 5.1 shows that the fixed-point equation
is automatic under coefficient-preserving lifting. All structural
content of the lifting --- both the nontrivial coefficients at
zero-digit positions and the sign-sum structure --- lives in the
coefficient positions corresponding to $F$'s zero digits, where it
does not affect $K(F) = F$. This is the fundamental observation
enabling the uniform construction across $d$.
\end{remark}

\begin{remark}[5.2]
Proposition 5.1 does \emph{not} assert universality
of the lifting --- it asserts only that the lifting satisfies
$K(F) = F$. Universality at $d > d_F$ requires that every admissible
input iterates to $F$, which is a substantively stronger claim. The
remainder of §5 is devoted to establishing this for the specific lifting
family constructed for $60714$.
\end{remark}

\begin{remark}[5.3 (discovered vs engineered)]
A natural objection is
that since the lifting is built so that appended coefficients vanish on
$F$'s zero-digit positions, $K(F) = F$ ``comes for free'' by
construction. This objection is correct as stated, and it is important
to be precise about what is and is not engineered. The fixed-point
equation $K^{(d)}(F) = F$ at $d > d_F$ is engineered: Proposition
5.1 shows it follows automatically from the coefficient-preserving
structure. \textbf{The non-trivial content lives elsewhere.} First, the
\emph{selection} of $F = 60714$ from the population of $33$
universal full-variable fixed points at $d = 5$ is not engineered ---
it is the unique survivor of the cross-dimensional test of Theorem 4.1,
with $32$ other candidates failing either algebraically (no rank-$6$
rule fixes them) or dynamically (the rules that fix them have basin
strictly less than $1$). Second, the \emph{dynamical} claim that the
lifted rule converges universally on $A_d \setminus E_d$ at every
$d \geq 5$ --- with $E_d = \emptyset$ at $d = 5, 6$ and an
exactly-characterized escape class continuous with the classical case at
$d \geq 7$ --- is a substantive theorem about the rule's iterated
behavior, not about the algebra of $K(F) = F$. Engineering the
algebraic identity is cheap; surviving the cross-dimensional selection
and yielding universal dynamics is not. The integer $60714$ is what
the framework selected, not what it was built to fit.
\end{remark}

\subsection{\texorpdfstring{The zero-sum pair construction for
$60714$}{The zero-sum pair construction for 60714}}\label{the-zero-sum-pair-construction-for-60714}

Let $K_{60714}^{(5)}$ denote the native rule at $d = 5$ for
$F = 60714$: permutations $\pi = (4, 1, 2, 3, 0)$,
$\sigma = (2, 0, 1, 4, 3)$, coefficient vector
$c^{(5)} = (9900, 9, 90, -9000, -999)$. The sorted-descending form of
$60714$ at $d = 5$ is $(7, 6, 4, 1, 0)$, so $N = \{0, 1, 2, 3\}$
and $Z = \{4\}$.

\begin{definition}[5.2 (zero-sum pair construction at
$d = d_F + 2k$)]
Given $K_{60714}^{(d)}$ at some $d$ on the
odd ladder ($d = 5, 7, 9, \ldots$) with coefficient vector $c^{(d)}$
and permutations $(\pi^{(d)}, \sigma^{(d)})$ using exponents
$\{0, 1, \ldots, d-1\}$, the \textbf{zero-sum pair lifting} to
$d + 2$ is the rule with permutations
\end{definition}

\[\pi^{(d+2)} = (\pi^{(d)}_0, \ldots, \pi^{(d)}_{d-1}, \; d + 1, \; d), \qquad \sigma^{(d+2)} = (\sigma^{(d)}_0, \ldots, \sigma^{(d)}_{d-1}, \; d, \; d + 1),\]

\emph{and coefficient vector}

\[c^{(d+2)} = (c_0^{(d)}, \ldots, c_{d-1}^{(d)}, \; 10^{d+1} - 10^d, \; 10^d - 10^{d+1}).\]

The last two appended coefficients are $+9 \cdot 10^d$ and
$-9 \cdot 10^d$, summing to zero. On the even ladder
($d = 6, 8, 10, \ldots$) the construction is the same with the signs
of the appended pair reversed; the two ladders differ only by this sign
convention.

\textbf{The two ladders for $60714$.}

\begin{itemize}
\item
  \textbf{Odd ladder root:} $K_{60714}^{(5)}$ as defined above.
  Coefficient vector $(9900, 9, 90, -9000, -999)$.
\item
  \textbf{Even ladder root:} the $d = 6$ \textbf{split lifting} of the
  native rule. Permutations $\pi = (4, 1, 2, 3, 5, 0)$,
  $\sigma = (2, 0, 1, 4, 3, 5)$, coefficient vector
  $c^{(6)} = (9900, 9, 90, -9000, 99000, -99999)$. This is a
  coefficient-preserving lifting at $d = 6$: the coefficients at
  nonzero-digit positions of $F_{(6)}$ are preserved
  ($9900, 9, 90, -9000$), while the native coefficient $-999$ at
  position $4$ is replaced by \emph{two} coefficients at positions
  $4$ and $5$ summing to $-999$: namely,
  $99000 + (-99999) = -999$. Both new coefficients are at zero-digit
  positions (the original position-$4$ zero and the new position-$5$
  zero from padding), so Proposition 5.1 applies.
\item
  \textbf{Iterative extension.} Starting from the odd ladder root, the
  odd ladder rules at $d = 7, 9, 11, \ldots$ are obtained by
  successive zero-sum pair liftings per Definition 5.2. Starting from
  the even ladder root, the even ladder rules at
  $d = 8, 10, 12, \ldots$ are obtained similarly.
\end{itemize}

Explicit rules at $d = 5, 6, \ldots, 20$ (and at $d = 100$) appear
in Appendix B.

\textbf{On the role of admissibility.} Throughout this section, $A_d$
denotes admissible inputs (non-repdigit and non-near-repdigit; see
§2.1). The exclusion is essential: repdigits trivially map to zero under
any sum-zero rule, and near-repdigits exhibit pathological trajectories
that escape the structural argument of Lemma 5.1 by projecting onto
inadmissible lower-dimensional inputs. The escape class $E_d$ defined
below is a \emph{strict subset} of $A_d$, characterizing the residual
class within the admissible domain. See §5.6 for full discussion of why
both restrictions are necessary.

\begin{definition}[5.2 (block-aligned multiset)]
Let $d_0 = 5$
if $d$ is odd, $d_0 = 6$ if $d$ is even. A sorted-descending
multiset $(x_0, x_1, \ldots, x_{d-1})$ at digit length $d$ is
\textbf{block-aligned} for the $60714$ ladder if:
\end{definition}

\emph{1. The first $d_0$ positions are all equal:
$x_0 = x_1 = \cdots = x_{d_0 - 1}$.}

\emph{2. Each appended pair has equal digits: $x_{d_0} = x_{d_0+1}$,
$x_{d_0+2} = x_{d_0+3}$, and so on.}

\emph{The set of block-aligned multisets at digit length $d$ is
denoted $B_d$.}

The block-aligned multisets are precisely the inputs on which the
constructed rule $K_{60714}^{(d)}$ evaluates to $0$ in a single
step: the $d_0$-block contributes
$(x_0)(c_0 + c_1 + \cdots + c_{d_0 - 1}) = 0$ because the base
coefficients sum to zero, and each appended zero-sum pair
$(c_k, -c_k)$ contributes $0$ when the corresponding two digits are
equal.

\begin{definition}[5.2.1 (escape class)]
The \textbf{escape
class} at digit length $d \geq 7$ is
\end{definition}
\[E_d \;=\; \{n \in A_d : K_{60714}^{(d)}{}^{\kappa}(n) \in B_d \text{ for some finite } \kappa\} \;\subseteq\; A_d.\]
\emph{By construction, $E_5 = E_6 = \emptyset$.}

\begin{theorem}[5.2 (Main theorem)]
The rules $K_{60714}^{(d)}$
constructed as above satisfy:
\end{theorem}

\emph{1. (\textbf{Algebraic.}) $K_{60714}^{(d)}(60714) = 60714$ at
every $d \geq 5$.}

\emph{2. (\textbf{Universal convergence on $A_d \setminus E_d$.}) For
every $d \geq 5$ and every admissible input
$n \in A_d \setminus E_d$, the orbit
$K_{60714}^{(d)}{}^{\kappa}(n) = 60714$ for some finite $\kappa$.
The escape class $E_d$ is empty at $d = 5, 6$ (so universal
convergence holds on all of $A_d$ there) and is non-empty at
$d \geq 7$, where it is structurally continuous with the classical
Kaprekar escape phenomenon at $d = 3, 4$ (Observation 7.6.0). Every
$n \in E_d$ has orbit reaching $0$ in finitely many steps; for
$d \leq 11$, exhaustive enumeration confirms collapse in at most $4$
iterations, but a uniform-in-$d$ bound remains open (Conjecture 7.6,
Basin Density Conjecture, Depth part).}

\begin{lemma}[5.5 (closed-form count of the step-1 escape class)]
Let $E_d^{(1)} = \{n \in A_d : K_{60714}^{(d)}(n) = 0\}$ be the
inputs that collapse in one step. Then $E_d^{(1)} = B_d \cap A_d$,
and
\end{lemma} \[|E_d^{(1)}| = \binom{10 + k - 1}{k} - 10\] \emph{where
$k = 1 + \lfloor (d - d_0)/2 \rfloor$ counts the number of blocks (one
base block plus appended pairs). The subtraction of $10$ removes the
repdigit case (a single digit replicated across all blocks).}

This closed form gives $|E_d^{(1)}|$ values
$45, 45, 210, 210, 705, 705, 1992, 1992, \ldots$ at
$d = 7, 8, 9, 10, 11, 12, 13, 14, \ldots$

The full escape class $E_d$ is the backward orbit of $B_d \cap A_d$
under $K_{60714}^{(d)}$. Empirical computation shows $|E_d|$ is
bounded by a small multiple of $|E_d^{(1)}|$ at each $d \leq 11$ and
that all orbits in $E_d$ reach $0$ in at most $4$ iterations. A
general bound on $|E_d|$ as $d \to \infty$ remains open.

The proof of Theorem 5.2 occupies the remainder of this section.

\subsection{Reduction to two lemmas}\label{reduction-to-two-lemmas}

The proof strategy has two parts.

First, we identify an \emph{absorbing set} $T_d \subseteq A_d$ with
two properties: (i) $T_d$ is closed under $K_{60714}^{(d)}$, and
(ii) on $T_d$, the rule $K_{60714}^{(d)}$ acts isomorphically to
$K_{60714}^{(d-2)}$ on a corresponding set at the lower digit length.
These two properties together mean that once an orbit enters $T_d$,
its convergence behavior reduces to the corresponding question at
$d - 2$.

Second, we show that every admissible input reaches $T_d$ in a bounded
number of iterations.

\begin{definition}[5.3 (tail-two-zeros set)]
At digit length
$d \geq 7$, the \textbf{tail-two-zeros set} $T_d \subseteq A_d$ is
the set of admissible inputs whose sorted-descending form ends in at
least two zeros: $x_{d-1} = x_{d-2} = 0$.
\end{definition}

\begin{lemma}[5.1 (Structural closure)]
For every $d \geq 7$ on
either ladder, and for every $n \in T_d$:
\end{lemma}

\begin{enumerate}
\def\labelenumi{\arabic{enumi}.}
\tightlist
\item
  $K_{60714}^{(d)}(n) \in T_d$ (the set is closed under the rule).
\item
  \emph{Writing $n$'s sorted-descending form as
  $(x_0, \ldots, x_{d-3}, 0, 0)$, we have}
  \[K_{60714}^{(d)}(n) = K_{60714}^{(d-2)}\bigl((x_0, \ldots, x_{d-3})\bigr)\]
  \emph{--- that is, the rule's action on $T_d$ is exactly the rule at
  $d - 2$ applied to the non-tail digits.}
\end{enumerate}

\begin{lemma}[5.2 (Bounded reaching time)]
For every $d \geq 7$
on either ladder, there exists a finite $T_d^*$ such that for every
$n \in A_d$, some iterate $K_{60714}^{(d)}{}^{\tau}(n) \in T_d$ with
$\tau \leq T_d^*$. In particular, orbits enter $T_d$ within a
bounded number of iterations.
\end{lemma}

Lemma 5.1 is structural and holds by direct computation on the zero-sum
pair construction. Lemma 5.2 is the substantive claim: it requires
ruling out orbits that cycle in $A_d \setminus T_d$ without ever
entering $T_d$.

Granting both lemmas, Theorem 5.2 follows by induction on $d$ along
each ladder:

\begin{itemize}
\tightlist
\item
  At $d = 5$ (odd ladder root) and $d = 6$ (even ladder root),
  universality is established by direct verification (these are the
  ladder-root cases from §3 and §4).
\item
  At $d \geq 7$ on a given ladder, every admissible input
  $n \in A_d$ reaches $T_d$ in $\leq T_d^*$ iterations by Lemma
  5.2. By Lemma 5.1, once in $T_d$, the orbit behaves as an orbit of
  $K_{60714}^{(d-2)}$ on the corresponding input at $d - 2$. By
  induction, this orbit reaches $60714$ in finitely many steps.
  Pulling back through the isomorphism of Lemma 5.1, the original orbit
  at $d$ reaches $60714$.
\end{itemize}

This completes the proof of Theorem 5.2 modulo the two lemmas, which are
proven in §5.4--§5.5.

\subsection{Proof of Lemma 5.1 (structural
closure)}\label{proof-of-lemma-5.1-structural-closure}

Fix $d \geq 7$ on the odd ladder (the even ladder argument is
identical with appropriate sign conventions). Let $n \in T_d$ with
sorted-descending form $(x_0, \ldots, x_{d-3}, 0, 0)$.

\textbf{Step 1: Computing $K(n)$.} The rule's coefficient vector at
$d$ is
$c^{(d)} = (c_0^{(d-2)}, \ldots, c_{d-3}^{(d-2)}, 10^{d-1} - 10^{d-2}, 10^{d-2} - 10^{d-1})$,
where the last two coefficients are the appended zero-sum pair.
Therefore

\[K^{(d)}(n) = \left| \sum_{i = 0}^{d - 3} c_i^{(d-2)} x_i + (10^{d-1} - 10^{d-2}) \cdot 0 + (10^{d-2} - 10^{d-1}) \cdot 0 \right| = \left| \sum_{i = 0}^{d - 3} c_i^{(d-2)} x_i \right|.\]

The last two terms vanish because $x_{d-2} = x_{d-1} = 0$. Therefore
$K^{(d)}(n)$ depends only on $(x_0, \ldots, x_{d-3})$, and its value
equals $K^{(d-2)}((x_0, \ldots, x_{d-3}))$. This proves (2).

\textbf{Step 2: Closure of $T_d$.} Let $m = K^{(d)}(n)$. By Step 1,
$m = K^{(d-2)}((x_0, \ldots, x_{d-3}))$. We must show that when $m$
is considered as a $d$-digit integer (by padding with two leading
zeros if necessary), its sorted-descending form at $d$ still ends in
two zeros.

Observe that $m \leq \max_{c \in \mathbb{Z}} |c| \cdot \max_x |x|$
over the coefficients and digits of $K^{(d-2)}$. Concretely, the value
of $K^{(d-2)}$ on any input at $d - 2$ is bounded above by
$10^{d-2}$ (the number has at most $d - 2$ digits, as it is the
output of a rule at $d - 2$). Hence $m < 10^{d - 2}$, which means
$m$ has at most $d - 2$ digits. When $m$ is considered as a
$d$-digit integer, its padded form has at least two leading zeros.

But a padded integer with at least two leading zeros has
sorted-descending form ending in at least two zeros (the leading zeros
sort to the smallest positions). Therefore $m \in T_d$, establishing
(1). $\square$

\begin{remark}[5.3]
Lemma 5.1 is purely structural: its proof uses only
the coefficient structure at zero-digit positions and the elementary
fact that a number with fewer than $d$ digits has a padded
representation with trailing sorted zeros. Nothing about the specific
value $60714$ or the specific native rule enters. Indeed, Lemma 5.1
holds verbatim for any zero-sum pair lifting of any
coefficient-preserving lifting family --- not just $60714$'s.
\end{remark}

\subsection{Proof of Lemma 5.2 (bounded reaching
time)}\label{proof-of-lemma-5.2-bounded-reaching-time}

Lemma 5.2 is the substantive content of Theorem 5.2. We prove it in two
parts: a finite-state verification at $d \in \{7, 8, \ldots, 16\}$
(§5.5.1), and a $d$-independent algebraic argument establishing
one-step $T_d$ closure at $d \geq 15$ (odd) and bounded (at most
two-step) $T_d$ closure at $d \geq 16$ (even) (§5.5.2). Together,
these cover all $d \geq 7$.

\subsubsection{\texorpdfstring{Finite-state verification at
$d \leq 16$}{Finite-state verification at d \textbackslash leq 16}}\label{finite-state-verification-at-d-leq-16}

\begin{proposition}[5.2 (finite-state verification)]
For each
$d \in \{7, 8, 9, 10, 11, 12, 13, 14, 15, 16\}$ and each ladder's rule
$K_{60714}^{(d)}$, every admissible input $n \in A_d$ reaches
$T_d$ within $\leq 8$ iterations.
\end{proposition}

\begin{proof}
Exhaustive enumeration. For each $d$ and each ladder's
rule, we iterate $K_{60714}^{(d)}$ starting from every admissible
digit multiset at $d$ --- that is, every multiset of $d$ digits from
$\{0, 1, \ldots, 9\}$ with at least two distinct digit values and not
a near-repdigit. At each iteration, we test whether the current multiset
has at least two trailing zeros in its sorted-descending form; if yes,
the orbit has entered $T_d$.

The set of admissible multisets at $d \leq 16$ is large but finite:

\begin{longtable}[]{@{}ccc@{}}
\toprule\noalign{}
$d$ & Admissible multisets & Max reaching time \\
\midrule\noalign{}
\endhead
\bottomrule\noalign{}
\endlastfoot
$7$ & $11{,}340$ & $8$ \\
$8$ & $24{,}210$ & $6$ \\
$9$ & $48{,}520$ & $3$ \\
$10$ & $92{,}278$ & $3$ \\
$11$ & $167{,}860$ & $2$ \\
$12$ & $293{,}830$ & $2$ \\
$13$ & $497{,}320$ & $2$ \\
$14$ & $817{,}090$ & $2$ \\
$15$ & $1{,}307{,}404$ & $1$ \\
$16$ & $2{,}042{,}875$ & $1$ \\
\end{longtable}

For every multiset at every $d$, every ladder, we verify
computationally that iteration enters $T_d$ within the stated reaching
time. There are no exceptions across the $\approx 25$ million
multisets tested. The computation is implemented in Python; the full
enumeration log appears in Appendix D.

The upper bound $\leq 8$ iterations holds at every $d \leq 16$. At
$d \geq 15$, reaching time is $\leq 1$ --- orbits enter $T_d$ in a
single step.
\end{proof}

\begin{remark}[5.4]
Counts in the table equal $\binom{d+9}{9} - 100$
(subtracting the $10$ repdigits and $90 \cdot d / d = 90$
near-repdigits, but here counted as digit multisets so the correction is
exactly $100$ for all $d \geq 4$). The enumeration is over digit
multisets, not over distinct integer inputs, because $K_{60714}^{(d)}$
depends only on the sorted-descending form of its input --- all integers
with the same multiset have identical orbits. This gives a
factor-of-$d!$ reduction in the verification complexity. The number of
non-repdigit non-near-repdigit multisets at $d = 16$ is just over
$2$ million, well within computational reach.
\end{remark}

\subsubsection{\texorpdfstring{$d$-independent algebraic argument at
large
$d$}{d-independent algebraic argument at large d}}\label{d-independent-algebraic-argument-at-large-d}

\begin{proposition}[5.3]
On the odd ladder at every
$d \geq 15$, every admissible input $n \in A_d$ satisfies
$K_{60714}^{(d)}(n) \in T_d$ (one-step $T_d$ closure). On the even
ladder at every $d \geq 16$, every admissible input reaches $T_d$ in
at most two iterations (bounded two-step $T_d$ closure). In both
cases, orbits enter $T_d$ within a bounded number of iterations, the
bound being uniform in $d$.
\end{proposition}

The proof uses two lemmas on the native rule's \emph{core contribution}
under sorted-descending inputs.

\begin{definition}[5.4 (native core contribution)]
At digit
length $d$ on either ladder, the \textbf{core contribution} of an
input $n$ with sorted-descending form $(x_0, \ldots, x_{d-1})$ is
\end{definition}

\[\mathrm{core}(n) = \sum_{i = 0}^{d_{F} - 1} c_i^{(d_{F})} \cdot x_i = 9900 \, x_0 + 9 \, x_1 + 90 \, x_2 - 9000 \, x_3 - 999 \, x_4\]

\emph{(for the odd ladder, taking $d_F = 5$).}

The core contribution is the sum, over the first $d_F$
sorted-descending positions, of the native coefficient times the digit.
It is the ``$d = 5$-style'' part of the full $K_{60714}^{(d)}$
computation, with the remainder coming from the appended zero-sum pair
contributions at positions $d_F, d_F + 1, \ldots, d - 1$.

\begin{lemma}[5.3 (Core non-negativity)]
For every
sorted-descending digit sequence $(x_0, \ldots, x_{d-1})$ with
$x_0 \geq x_1 \geq \cdots \geq x_{d-1} \geq 0$, the core contribution
$\mathrm{core}(n) = 9900 x_0 + 9 x_1 + 90 x_2 - 9000 x_3 - 999 x_4 \geq 0$.
\end{lemma}

\begin{proof}
By the sorted-descending constraint $x_0 \geq x_3$ and
$x_0 \geq x_4$. Bound:

\[\mathrm{core}(n) \geq 9900 x_0 + 0 + 0 - 9000 x_0 - 999 x_0 = (9900 - 9000 - 999) x_0 = -99 x_0 \geq -99 \cdot 9 = -891,\]

which does not immediately give non-negativity. A sharper argument is
needed, exploiting the relations $x_1, x_2 \leq x_0$ and
$x_3, x_4 \leq \min(x_0, x_1, x_2)$. The fully worked bound is in
Appendix C.
\end{proof}

\begin{lemma}[5.4 (Core upper bound)]
For every sorted-descending
sequence at $d \geq 7$, $\mathrm{core}(n) \leq 10^{d-2}$.
\end{lemma}

\begin{proof}
Direct:
$9900 x_0 + 9 x_1 + 90 x_2 - 9000 x_3 - 999 x_4 \leq 9900 \cdot 9 + 9 \cdot 9 + 90 \cdot 9 = 89{,}100 + 81 + 810 = 89{,}991 < 10^5 = 10^{d_{F}}$
for $d_F = 5$. For $d \geq 7$, $10^{d-2} \geq 10^5$, so the bound
holds.
\end{proof}

\textbf{Proof of Proposition 5.3 (outline).} The argument uses a
\emph{block-structure decomposition} of $K^{(d)}(n)$: each zero-sum
pair at positions $(3+2k, 4+2k)$ contributes a value of the form
$9 \cdot 10^{3+2k} \delta_k$ where
$\delta_k = x_{3+2k} - x_{4+2k} \in \{0, 1, \ldots, 9\}$ by
sorted-descending. Combined with Lemma 5.3 (core non-negativity), the
argument inside the absolute value is non-negative:

\[K^{(d)}(n) = \mathrm{core}(n) + \sum_{k=1}^{M} 9 \cdot 10^{3+2k} \, \delta_k, \qquad M = (d-5)/2.\]

A key telescoping observation (Lemma C.6, \emph{cliff-sequence bound})
gives $\sum_k \delta_k \leq 9$. Combined with the decimal-block
structure --- each $\delta_k = 0$ gives a $(0,0)$ block (2 zero
digits), each $\delta_k = 1$ gives a $(0,9)$ block (1 zero digit),
each $\delta_k \geq 2$ gives a block with no zero digits --- a
counting argument shows that $K^{(d)}(n)$ has at least $2M - 9$ zero
digits in the block region. For the odd ladder at $d \geq 15$ (so
$M \geq 5$), this gives at least one zero from the block region, which
combined with at least one zero from the core contribution (Lemma 5.4,
$\mathrm{core} < 10^5$) yields the required two zero digits, so
$K^{(d)}(n) \in T_d$ in one step.

For the even ladder at $d \geq 16$, the analogous bound holds with
$M' = (d-6)/2$. At $d = 18$ the bound is marginally sufficient but
does not yet give strict one-step closure (the split-lifted coefficients
at positions $4, 5$ can interact with the first appended block via
carries in $\approx 0.86\%$ of inputs); in these cases, direct
computation shows $T_d$ is reached in exactly two iterations. At
$d \geq 20$ on the even ladder, the block-count margin grows and
one-step closure holds.

The full algebraic computation --- the cliff-sequence bound, the
block-structure disjointness, and the zero-counting argument --- appears
as Lemmas C.3 and C.4 in Appendix C. Here we record only the resulting
bound: uniform reaching time $\leq 1$ on the odd ladder at
$d \geq 15$, and uniform reaching time $\leq 2$ on the even ladder
at $d \geq 16$. $\square$

Propositions 5.2 and 5.3 together establish Lemma 5.2.

\subsection{Near-repdigit exclusion}\label{near-repdigit-exclusion}

Near-repdigit inputs --- multisets where one digit appears $d - 1$ or
$d$ times --- can produce pathological trajectories under
coefficient-preserving liftings. Specifically, an input like
$(X, X, X, X, X, X, 0, 0)$ at $d = 8$ (six-of-a-kind plus two zeros)
lies in $T_d$ (two trailing zeros) but iterates to
$K^{(d)}((X, X, X, X, X, X)) = K^{(d-2)}$ applied to an input with no
zero-digit structure at $d - 2$. Such inputs can escape back out of
$T_d$ in pathological ways, and in some cases the orbit collapses to
zero instead of reaching $60714$.

We exclude near-repdigits from $A_d$ precisely to prevent such
pathologies. The resulting admissible set $A_d$ still contains an
overwhelming majority of $d$-digit integers (for $d \geq 7$, the
ratio $|A_d| / 10^d \to 1$ as $d$ grows, and at $d = 7$ already
exceeds $99.9\%$). The restriction is standard in Kaprekar analysis
and is compatible with the spirit of Kaprekar's original theorem, which
also excludes repdigits.

\textbf{Note on near-repdigits and $60714$.} In our finite-state
enumeration at $d \leq 16$, every admissible (non-near-repdigit) input
reaches $T_d$ within $\leq 8$ iterations. If we extended to
near-repdigit inputs, a small number would fail --- these are the
$45$-multiset escape class documented for $6174$ at §6 and a smaller
analogous class for $60714$. The exclusion of near-repdigits in
$A_d$ removes these pathological inputs from consideration.

\subsection{Completing the proof of Theorem
5.2}\label{completing-the-proof-of-theorem-5.2}

We now complete the proof of Theorem 5.2.

\textbf{Proof of Theorem 5.2.} Proceed by strong induction on $d$
along each ladder.

\textbf{Base cases.} At $d = 5$ (odd ladder root), $K_{60714}^{(5)}$
is universal for $60714$ by direct verification: exhaustive iteration
over $A_5$ shows every admissible input reaches $60714$ within a
bounded number of steps (see §3 and Appendix A for basin verification).
At $d = 6$ (even ladder root), $K_{60714}^{(6)}$ is universal for
$60714$ by the split-lifting construction of §4 and direct basin
verification at $d = 6$ (see Proposition 4.1).

\textbf{Inductive step.} Fix $d \geq 7$ on either ladder, and suppose
universality holds at $d - 2$ on the same ladder. Let $n \in A_d$.
By Lemma 5.2, some iterate $m = K^{(d)}{}^{\tau}(n) \in T_d$ with
$\tau \leq T_d^* < \infty$. By Lemma 5.1, subsequent iterations
satisfy

\[K^{(d)}(m) = K^{(d-2)}((\text{non-tail digits of }m)),\]

and the result stays in $T_d$ by closure. The non-tail digits of $m$
constitute an input at $d - 2$ that may be admissible, repdigit, or
near-repdigit; the three cases are handled in Lemma C.5 (§C.6). For Case
1 (admissible projection), the inductive hypothesis applies directly.
For Case 1b (repdigit projection), the orbit reaches $0$, placing
$n$ in the residual escape class. For Case 2 (near-repdigit
projection), Lemma C.10 shows the inductive hypothesis applies after at
most two iterations of $K^{(d)}$.

By the induction hypothesis, the orbit of the non-tail digits under
$K^{(d-2)}$ reaches $60714$ in finitely many steps. Pulling back
through the isomorphism of Lemma 5.1, the orbit of $m$ under
$K^{(d)}$ reaches $60714_{(d)}$ (i.e., $60714$ padded to $d$
digits --- equivalently, the integer $60714$ viewed as a $d$-digit
integer with leading zeros). Therefore the original orbit of $n$
reaches $60714$ in finitely many total steps. $\square$

\subsection{Summary}\label{summary-1}

The structure of the proof is:

\begin{enumerate}
\def\labelenumi{\arabic{enumi}.}
\item
  \textbf{Coefficient-preserving lifting} (Definition 5.1, Proposition
  5.1): a purely algebraic construction that makes the fixed-point
  equation $K(F) = F$ automatic at every $d > d_F$.
\item
  \textbf{Zero-sum pair lifting} (Definition 5.2): a specific instance
  of coefficient-preserving lifting for $60714$ that iterates along
  two ladders.
\item
  \textbf{Structural closure} (Lemma 5.1): the tail-two-zeros set
  $T_d$ is invariant, and the rule reduces to the lower-dimensional
  rule on it. Proven by direct computation using only the zero-sum pair
  structure.
\item
  \textbf{Bounded reaching time} (Lemma 5.2): every admissible orbit
  enters $T_d$ within finitely many iterations. Proven by:

  \begin{itemize}
  \tightlist
  \item
    Finite-state enumeration at $d \in \{7, \ldots, 16\}$ (Proposition
    5.2).
  \item
    A $d$-independent algebraic argument: one-step $T_d$ closure at
    $d \geq 15$ on the odd ladder, and bounded (at most two-step)
    $T_d$ closure at $d \geq 16$ on the even ladder, using core
    non-negativity and the cliff-sequence bound on sorted-descending
    digit differences (Proposition 5.3, Lemmas 5.3, 5.4, and C.6; full
    proofs in Appendix C).
  \end{itemize}
\item
  \textbf{Induction along each ladder} (§5.7): base cases at
  $d = 5, 6$, inductive step via Lemmas 5.1 and 5.2.
\end{enumerate}

Theorem 5.2 is thus established without any conjectural step. The
finite-state enumeration at low $d$ and the algebraic bound at high
$d$ together cover every $d \geq 5$, and each step of the reduction
is structurally transparent.

\section{\texorpdfstring{The
$\{7, 6, 4, 1\}$-Thread}{The \textbackslash\{7, 6, 4, 1\textbackslash\}-Thread}}\label{the-7-6-4-1-thread}

This section connects the central result of §5 --- that 60714 is
universal at every $d \geq 5$ --- to the classical Kaprekar
literature. The four-digit fixed point 6174, identified by Kaprekar in
1949, lives in the same digit multiset as 60714 and gives the section
its motivation: 60714 is the $d = 5$ generalization that Kaprekar's
classical recipe could not see, because the classical ``biggest-first
minus smallest-first'' rule degenerates at $d = 5$ to a sv$=4$
rank-deficient operator with no fixed point.

The unifying structure is what we call the \textbf{\{7, 6, 4,
1\}-thread} --- the digit multiset shared by 6174 (at $d = 4$), 60714
and 60417 (at $d = 5$), 1746 (at $d = 4$, classical-rule reverse
anagram of 6174), and a four-fp cluster at $d = 6$ (146070, 170460,
607140, and 60714 lifted). Within this thread, the four fps studied here
exhibit a clean ordering of cross-dimensional behaviors that is captured
empirically by an invariant we call $\mathrm{sum\_locked\_spans}$. The
ordering is \textbf{within-thread only}: the invariant does not predict
cross-dimensional behavior across different digit multisets, as we
demonstrate by comparison to the \{2, 3, 5, 8\}-thread.

One of this section's claims (Theorem 6.1, the $6174$
cross-dimensional pattern) is proven unconditionally at specific digit
lengths by exhaustive verification. The $\mathrm{sum\_locked\_spans}$
unification (§6.3) is an empirical observation, explicitly labeled as
such, that organizes and explains the within-thread data but is not
proven as a theorem.

\subsection{6174 as predecessor: where the classical recipe runs
out}\label{as-predecessor-where-the-classical-recipe-runs-out}

Kaprekar's 1949 result identifies 6174 as the unique attractor of the
classical four-digit Kaprekar routine: starting from any non-repdigit
four-digit integer, the iteration
$K_0(n) = (\text{digits of } n \text{ rearranged descending}) - (\text{digits rearranged ascending})$
reaches 6174 within at most seven steps. In our framework, this is the
rule with permutation pair $\pi = (3, 2, 1, 0)$,
$\sigma = (0, 1, 2, 3)$ and coefficient vector
$(999, 90, -90, -999)$.

At $d = 5$, the analogous rule $\mathrm{abcde} - \mathrm{edcba}$ has
coefficient vector $(9999, 990, 0, -990, -9999)$. The middle
coefficient is identically zero: the classical rule at $d = 5$ has
algebraic rank $4$, not $5$. The classical recipe degenerates at odd
digit lengths in this way, and the result is well-known: there is no
fixed point of the classical $d = 5$ rule. Every input enters a
four-cycle $(74943, 62964, 71973, 83952)$ that catches the entire
admissible set.

This is partially what 60714 addresses. The fixed point 60714 lives in
the digit multiset $\{7, 6, 4, 1, 0\}$ --- the same digits as 6174 (at
$d = 4$, padded with a zero at $d = 5$) --- but it is \emph{not} a
fixed point of the classical rule, which has no fixed points at
$d = 5$. Instead, 60714 is the universal attractor of a
\emph{different} permutation pair: the rule
$\mathrm{adcbe} - \mathrm{deacb}$, with $\pi = (4, 1, 2, 3, 0)$,
$\sigma = (2, 0, 1, 4, 3)$ and coefficient vector
$(9900, 9, 90, -9000, -999)$. This rule is full-variable (sv$=5$),
and 60714 is its unique universal fixed point at $d = 5$.

The thread that connects 6174 to 60714 is therefore not direct
succession --- there is no rule that fixes both --- but \textbf{shared
digit content}. Both fps live in the multiset $\{7, 6, 4, 1, 0\}$
(with 6174 padded with one zero at $d \geq 5$), and the
cross-dimensional behavior of each gives information about the
structural role of the multiset.

\subsection{Cross-dimensional pattern of
6174}\label{cross-dimensional-pattern-of-6174}

Even though the classical rule fails at $d = 5$, 6174 itself can be
lifted to higher $d$ via coefficient-preserving liftings of its native
$d = 4$ rule. Section §5 showed that 60714 admits
coefficient-preserving liftings at every $d \geq 5$ that are
universal. We now ask: what does the analogous lifting do for 6174?

\begin{theorem}[6.1 (cross-dimensional pattern of 6174)]
Let
$K_{6174}^{(d)}$ denote the natural coefficient-preserving lifting of
the classical Kaprekar rule $K_0$ from $d = 4$ to digit length
$d$, where the native coefficients $(999, 90, -90, -999)$ are
preserved at the four nonzero-digit positions of 6174's padded
sorted-descending form $(7, 6, 4, 1, 0, \ldots, 0)$, and the remaining
positions form a derangement of zero-sum pair structure. Then:
\end{theorem}

\begin{enumerate}
\def\labelenumi{\arabic{enumi}.}
\item
  \emph{$d = 4$ (native, strict universality).} $K_0$ is universal
  for 6174 at $d = 4$, with basin $1$ over the $9{,}630$
  admissible four-digit integers ($615$ admissible digit multisets).
  This is Kaprekar's 1949 result.
\item
  \emph{$d = 5$ (algebraic obstruction).} No sv$=5$ rule fixes 6174
  at $d = 5$. The trivial sv$=4$ extension of $K_0$ (coefficients
  $(999, 90, -90, -999, 0)$) achieves basin $0.99775$ over the
  $99{,}540$ admissible $d = 5$ multisets --- near-universal but not
  strict.
\item
  \emph{$d = 6$ (dynamic obstruction). Among the four sv$=6$ rules
  fixing 6174 at $d = 6$ (two coefficient-preserving liftings of
  $K_0$ plus their sign-flips), the best basin achieved is
  $0.968603$ over the $4{,}905$ admissible multisets at $d = 6$
  --- below the NEAR threshold of $0.99$. Thus 6174 is dynamically
  obstructed at $d = 6$: full-variable rules fix it, but none is even
  near-universal.}
\item
  \emph{$d = 7$ (NEAR-universal). Among the twelve sv$=7$
  coefficient-preserving liftings of $K_0$ at $d = 7$, no rule is
  strict-universal. The best basin is $0.989683$ over the $11{,}340$
  admissible multisets, achieved by the rule with
  $\pi = (3, 2, 1, 0, 6, 5, 4)$, $\sigma = (0, 1, 2, 3, 4, 6, 5)$
  and coefficient vector
  $(999, 90, -90, -999, 990000, -900000, -90000)$.}
\item
  \emph{$d = 8$ (NEAR-universal). Among the $216$
  coefficient-preserving sv$=8$ liftings of $K_0$, no rule is
  strict-universal. The best basin is $0.998141$. The $45$
  non-reaching inputs all have sorted-descending form
  $(X, X, X, X, Y, Y, Y, Y)$ for distinct digits $X > Y \geq 0$, and
  each collapses to $0$ rather than reaching 6174. Verified by
  exhaustive enumeration over the $24{,}210$ admissible multisets at
  $d = 8$.}
\item
  \emph{$d = 9$ (NEAR-universal with improved ratio). Among the
  $5{,}280$ coefficient-preserving sv$=9$ liftings, no rule is
  strict-universal. The best basin is $0.999073$. The escape class is
  the $d = 8$ structure extended by one trailing zero: $45$
  multisets of form $(X, X, X, X, Y, Y, Y, Y, 0)$ with
  $X > Y \geq 0$, each collapsing to $0$. Verified over the
  $48{,}520$ admissible multisets at $d = 9$.}
\end{enumerate}

\begin{proof}
Statement 1 is Kaprekar's original result. Statement 2
is verified by exhaustive enumeration over the
$5! \cdot D_5 = 5{,}280$ full-variable rules at $d = 5$: zero of
them fix 6174's padded form $(7, 6, 4, 1, 0)$, establishing algebraic
obstruction. Statements 3--6 are proven by direct enumeration at the
corresponding $d$. At each $d$, all coefficient-preserving liftings
are enumerated, the best basin is computed, and the non-reaching inputs
are characterized. Full enumeration logs and escape-set characterization
are in Appendix E.
\end{proof}

\textbf{The pattern: monotone NEAR-universality.} Statements 3--6
establish that 6174's basin under coefficient-preserving lifting
\textbf{monotonically improves} toward $1$ as $d$ grows:
$0.969 \to 0.989683 \to 0.998141 \to 0.999073$ at $d = 6, 7, 8, 9$.
The escape class is structurally fixed (the four-of-a-kind-paired
multisets at $d \geq 8$), but the admissible multiset count grows
combinatorially with $d$, so the ratio of escapes to admissible inputs
decays toward zero. This is \emph{near-universality with a fixed escape
class}, not the non-monotone ``fingerprint'' sometimes attributed to
6174.

\begin{remark}[on escape structure]
At $d \geq 8$, the $45$
non-reaching inputs have a specific combinatorial form: four-of-a-kind
of one digit $X$ paired with four-of-a-kind of another digit $Y$,
with $X > Y \geq 0$. There are exactly $\binom{10}{2} = 45$ such
pairs $(X, Y)$. These multisets are just outside the near-repdigit
threshold (four-of-a-kind rather than seven- or eight-of-a-kind), so
they remain admissible. Their trajectories under any
coefficient-preserving lifting collapse to zero rather than reaching
6174, because the rule acts on an effectively-rank-2 structure (only two
distinct digit values contribute), and the sum-zero coefficient
structure produces zero algebraically. The escape count $45$ is
$d$-independent: the same $45$ multisets reappear at $d = 8, 9$
(with appropriate trailing zeros) and presumably higher $d$. Since
$|A_d|$ grows with $d$ while the escape count stays fixed, the basin
ratio asymptotically approaches $1$ as $d \to \infty$.
\end{remark}

\subsection{60714: the d=5 generalization 6174 hinted
at}\label{the-d5-generalization-6174-hinted-at}

The contrast with 60714 is immediate. Where 6174 achieves only
near-universality at $d \geq 5$, \textbf{60714 achieves universal
convergence on $A_d \setminus E_d$ at every $d \geq 5$ --- with
$E_d = \emptyset$ at $d = 5, 6$ and an exactly-characterized escape
class continuous with the classical case at $d \geq 7$} (Theorem 5.2).
The two fixed points share digit multiset $\{7, 6, 4, 1, 0\}$ at every
$d \geq 5$ (with 6174 padded by additional zeros), but they differ in
three structural respects:

\begin{enumerate}
\def\labelenumi{\arabic{enumi}.}
\item
  \emph{Native digit length.} 6174 is native at $d = 4$ (no zero
  digits in its sorted-descending form $(7, 6, 4, 1)$); 60714 is
  native at $d = 5$ (one zero digit in $(7, 6, 4, 1, 0)$).
\item
  \emph{Native rule.} 6174's native rule is the classical Kaprekar rule
  $\mathrm{abcd} - \mathrm{dcba}$, with sum-zero coefficient vector
  $(999, 90, -90, -999)$. 60714's native rule is
  $\mathrm{adcbe} - \mathrm{deacb}$, with coefficient vector
  $(9900, 9, 90, -9000, -999)$ --- a non-classical rule that exploits
  the absorbing capacity of 60714's zero digit at sorted-descending
  position 4.
\item
  \emph{Cross-dimensional fate.} 6174's coefficient-preserving lifts
  achieve only near-universality at $d \geq 5$, with a fixed
  $45$-input escape class. 60714's coefficient-preserving lifts are
  universally convergent on $A_d \setminus E_d$ at every $d \geq 5$
  --- with $E_d = \emptyset$ at $d = 5, 6$ and an
  exactly-characterized escape class of block-aligned multisets at
  $d \geq 7$.
\end{enumerate}

The structural reason for the difference traces to the position of the
absorbing zero. 6174's lifts at $d \geq 5$ pad with zeros at positions
$\geq 4$, \emph{outside} the four nonzero-digit positions where the
native rule's coefficients are concentrated. The four-of-a-kind escape
class arises because the sum-zero coefficient structure of $K_0$ acts
trivially on inputs with only two distinct digit values. 60714's native
rule, by contrast, has its largest-magnitude coefficient $(-999)$
landing on its zero digit at sorted-descending position 4 --- that is,
\emph{within} the absorbing region. Coefficient-preserving liftings of
60714 inherit this structural placement, and the resulting rules have no
analogous escape class.

This is the structural connection between 6174 and 60714 that the \{7,
6, 4, 1\}-thread highlights. They share digit content but realize that
content through different rules with different absorbing-position
structure, leading to qualitatively different cross-dimensional fates.

\subsection{An empirical observation: sum\_locked\_spans within the
thread}\label{an-empirical-observation-sum_locked_spans-within-the-thread}

\emph{The remainder of §6 (§6.4 and §6.5) presents an empirical
observation about the four fixed points 6174, 1746, 60714, and 60417,
expressed through a structural invariant we call
$\mathrm{sum\_locked\_spans}$. \textbf{This material is not a
theorem.} The observation is verified on these specific four fixed
points within the $\{7, 6, 4, 1\}$-multiset thread; it does not
generalize across multisets, as a direct control test confirms (§6.5).
We present this material because it captures a structural pattern that
is suggestive and reproducible within the thread, but we are explicit
that it is empirical, not proven, and that its scope is bounded. A
formal conjecture extending the empirical pattern appears in §7.2.}

We now introduce the structural invariant that orders cross-dimensional
behavior across the four fps studied in this section.

\begin{definition}[6.1 (locked span and sum)]
For a fixed point
$F$ with sorted-descending form $(f_0, \ldots, f_{d-1})$ at native
digit length $d_F$, and for a rule $K_F$ at $d_F$ with
permutations $(\pi, \sigma)$, the \textbf{span} of position $i$ is
$s_i = |\pi_i - \sigma_i|$. The \textbf{locked spans} are those at
nonzero-digit positions of $F$: $\{s_i : f_i \neq 0\}$. The
\textbf{sum of locked spans} is
\end{definition}

\[\mathrm{sum\_locked\_spans}(K_F) = \sum_{i : f_i \neq 0} |\pi_i - \sigma_i|.\]

\textbf{Computation at four thread fixed points.} For each of 60714,
60417, 6174, and 1746, we compute $\mathrm{sum\_locked\_spans}$ of a
representative universal native rule.

\begin{longtable}[]{@{}
  >{\centering\arraybackslash}p{(\columnwidth - 10\tabcolsep) * \real{0.1667}}
  >{\centering\arraybackslash}p{(\columnwidth - 10\tabcolsep) * \real{0.1667}}
  >{\centering\arraybackslash}p{(\columnwidth - 10\tabcolsep) * \real{0.1667}}
  >{\centering\arraybackslash}p{(\columnwidth - 10\tabcolsep) * \real{0.1667}}
  >{\centering\arraybackslash}p{(\columnwidth - 10\tabcolsep) * \real{0.1667}}
  >{\centering\arraybackslash}p{(\columnwidth - 10\tabcolsep) * \real{0.1667}}@{}}
\toprule\noalign{}
\begin{minipage}[b]{\linewidth}\centering
Fixed point
\end{minipage} & \begin{minipage}[b]{\linewidth}\centering
Native $d_F$
\end{minipage} & \begin{minipage}[b]{\linewidth}\centering
Sorted-desc form
\end{minipage} & \begin{minipage}[b]{\linewidth}\centering
Native $\pi$
\end{minipage} & \begin{minipage}[b]{\linewidth}\centering
Native $\sigma$
\end{minipage} & \begin{minipage}[b]{\linewidth}\centering
$\mathrm{sum\_locked\_spans}$
\end{minipage} \\
\midrule\noalign{}
\endhead
\bottomrule\noalign{}
\endlastfoot
$1746$ & $4$ & $(7, 6, 4, 1)$ & $(0, 3, 2, 1)$ &
$(1, 2, 3, 0)$ & $4$ \\
$60714$ & $5$ & $(7, 6, 4, 1, 0)$ & $(4, 1, 2, 3, 0)$ &
$(2, 0, 1, 4, 3)$ & $5$ \\
$60417$ & $5$ & $(7, 6, 4, 1, 0)$ & $(4, 1, 0, 3, 2)$ &
$(0, 2, 1, 4, 3)$ & $7$ \\
$6174$ & $4$ & $(7, 6, 4, 1)$ & $(3, 2, 1, 0)$ &
$(0, 1, 2, 3)$ & $8$ \\
\end{longtable}

These are the natural sv$=d_F$ universal rules at each fp's native
digit length. The four thread fps share digit multiset
$\{7, 6, 4, 1\}$ (at their native digit lengths) but realize that
multiset through different permutation pairs, as reflected in their
distinct sum\_locked\_spans values.

\subsection{The unified empirical
observation}\label{the-unified-empirical-observation}

\begin{observation}[6.2 (within-thread, empirical, not proven)]
Within the $\{7, 6, 4, 1\}$-thread, sum\_locked\_spans
empirically orders cross-dimensional behavior at $d = 7$:
\end{observation}

\begin{longtable}[]{@{}
  >{\centering\arraybackslash}p{(\columnwidth - 6\tabcolsep) * \real{0.2500}}
  >{\centering\arraybackslash}p{(\columnwidth - 6\tabcolsep) * \real{0.2500}}
  >{\centering\arraybackslash}p{(\columnwidth - 6\tabcolsep) * \real{0.2500}}
  >{\centering\arraybackslash}p{(\columnwidth - 6\tabcolsep) * \real{0.2500}}@{}}
\toprule\noalign{}
\begin{minipage}[b]{\linewidth}\centering
Fixed point
\end{minipage} & \begin{minipage}[b]{\linewidth}\centering
$\mathrm{sum\_locked\_spans}$
\end{minipage} & \begin{minipage}[b]{\linewidth}\centering
Best basin at $d = 7$
\end{minipage} & \begin{minipage}[b]{\linewidth}\centering
\# NEAR rules
\end{minipage} \\
\midrule\noalign{}
\endhead
\bottomrule\noalign{}
\endlastfoot
$1746$ & $4$ & $0.996032$ & $12$ \\
$60714$ & $5$ & $0.992857$ (NEAR, basin over $A_7 \setminus E_7$
is $1.000000$) & (within-thread sls invariant) \\
$60417$ & $7$ & $0.992857$ & $2$ \\
$6174$ & $8$ & $0.989683$ & $0$ \\
\end{longtable}

\emph{The pattern within the thread: lower sum\_locked\_spans correlates
with higher cross-dimensional basin. 60714 (sls $= 5$) achieves the
highest basin among the four, and uniquely achieves strict universality
at $d = 5, 6$ via its native lifting family; at $d = 7$, all four
fps in the thread show NEAR-universal behavior. The invariant captures
something real about how the rule's coefficient structure interacts with
the absorbing-position placement of $F$.}

\textbf{Crucially, this is an empirical observation, not a theorem, and
it is restricted to the \{7, 6, 4, 1\}-thread.} We verified
within-thread ordering at $d = 5, 6, 7$ (above and elsewhere in this
paper), but the invariant does \emph{not} generalize to predict
cross-dimensional behavior across multisets. As a control, we tested the
$\{2, 3, 5, 8\}$-thread (the other multiset cluster at $d = 4$,
supporting fps $2538$ and $5382$):

\begin{longtable}[]{@{}
  >{\centering\arraybackslash}p{(\columnwidth - 6\tabcolsep) * \real{0.2500}}
  >{\centering\arraybackslash}p{(\columnwidth - 6\tabcolsep) * \real{0.2500}}
  >{\centering\arraybackslash}p{(\columnwidth - 6\tabcolsep) * \real{0.2500}}
  >{\centering\arraybackslash}p{(\columnwidth - 6\tabcolsep) * \real{0.2500}}@{}}
\toprule\noalign{}
\begin{minipage}[b]{\linewidth}\centering
Fixed point
\end{minipage} & \begin{minipage}[b]{\linewidth}\centering
Multiset
\end{minipage} & \begin{minipage}[b]{\linewidth}\centering
$\mathrm{sum\_locked\_spans}$
\end{minipage} & \begin{minipage}[b]{\linewidth}\centering
Best basin at $d = 7$
\end{minipage} \\
\midrule\noalign{}
\endhead
\bottomrule\noalign{}
\endlastfoot
$1746$ & $\{1, 4, 6, 7\}$ & $4$ & $0.996032$ \\
$5382$ & $\{2, 3, 5, 8\}$ & $4$ & $0.996032$ \\
$2538$ & $\{2, 3, 5, 8\}$ & $8$ & $0.996032$ \\
$6174$ & $\{1, 4, 6, 7\}$ & $8$ & $0.989683$ \\
\end{longtable}

Across the two multisets, sum\_locked\_spans does \emph{not} track best
basin: $5382$ (sls $= 4$, in the $\{2, 3, 5, 8\}$-thread) and
$2538$ (sls $= 8$, same thread) both reach $0.996032$ at
$d = 7$, identical to $1746$ (sls $= 4$, in the
$\{7, 6, 4, 1\}$-thread). The invariant orders within-thread behaviors
at $d = 7$ but does not generalize to a cross-multiset principle. A
formal conjecture in this direction appears in §7.

\subsection{Context and open
directions}\label{context-and-open-directions}

The $\{7, 6, 4, 1\}$-thread exposes three phenomena worth noting for
the broader research direction.

\textbf{1. The classical Kaprekar recipe degenerates at odd digit
lengths.} The rule $\mathrm{abc \cdots} - (\text{reverse})$ has
algebraic rank $d$ at even $d$ but rank $d - 1$ at odd $d$,
because the middle coefficient $10^{(d-1)/2} - 10^{(d-1)/2} = 0$
vanishes identically. Universal fixed points of this rule therefore can
only exist at even $d$ (where 6174 sits at $d = 4$, and
$631{,}764$ --- though not full-variable in our sense --- sits at
$d = 6$; cf.~§3). At odd $d$, universal full-variable fps must come
from non-classical rules. 60714 is the canonical such fp at $d = 5$.

\textbf{2. 6174 and 60714 share digit content but realize it through
different rules.} The two fps live in the same digit multiset
$\{7, 6, 4, 1, 0\}$ (with 6174 padded) but are fixed points of
structurally different permutation pairs. 6174's native rule has
sum-zero coefficient structure $(999, 90, -90, -999)$ that produces an
algebraic four-of-a-kind escape class. 60714's native rule places its
largest coefficient on the zero digit, eliminating the analogous escape
class. The result: 60714 is universally convergent on
$A_d \setminus E_d$ at every $d \geq 5$ with $E_d = \emptyset$ at
$d = 5, 6$, while 6174 has a $45$-input escape class persisting at
every $d \geq 8$ under its native lifting.

\textbf{3. 60714 is the uniquely complete case.} Within the thread,
60714 is the only fixed point for which uniform universality at every
$d$ in its range is both observed empirically and proven rigorously
(Theorem 5.2). 6174 is partially universal (NEAR everywhere from
$d = 5$ onward) with a proven pattern; 60417 is dimension-locked at
$d \geq 6$. The $\mathrm{sum\_locked\_spans}$ invariant provides an
empirical within-thread ordering of these three behaviors (plus 1746 as
a within-multiset control) but does not yet yield a uniform theorem
analogous to Theorem 5.2 for fps other than 60714. Extending the
lifting-theorem framework to 6174 --- proving the monotone
NEAR-convergence of basin to $1$ across all $d \geq 5$ --- is a
natural open direction (§7).

\subsection{Summary}\label{summary-2}

Theorem 6.1 establishes the cross-dimensional pattern of 6174 at
$d = 4, \ldots, 9$: strict universal at $d = 4$ (Kaprekar's 1949
result), algebraically obstructed at $d = 5$, dynamically obstructed
at $d = 6$ (basin $0.969$ below NEAR threshold), NEAR-universal at
$d = 7, 8, 9$ with monotonically improving basin and a precisely
characterized $45$-input escape class at $d \geq 8$. The
$\mathrm{sum\_locked\_spans}$ invariant (§6.4) organizes the
four-fixed-point thread (Observation 6.2) as a within-thread empirical
regularity: 1746 (sls $= 4$, NEAR), 60714 (sls $= 5$, strict
universal), 60417 (sls $= 7$, dimension-locked beyond $d = 5$), and
6174 (sls $= 8$, monotone NEAR). This unification is explicitly
labeled as empirical and within-thread; cross-multiset testing on the
$\{2, 3, 5, 8\}$-thread shows the invariant is not a global predictor.
60714 is the uniquely complete case --- the universal attractor that the
classical Kaprekar recipe at $d = 5$ could not see, but which exists
in the same digit multiset as 6174 once non-classical full-variable
rules are admitted.

\emph{End of §6.}

\section{Open Questions and
Conjectures}\label{open-questions-and-conjectures}

The results of §3--§6 establish a specific empirical and structural
picture at digit lengths $d \leq 9$ for the multiset thread
$\{7, 6, 4, 1\}$, with a full theorem (§5, Theorem 5.2) covering
$60714$ at every $d \geq 5$. Several natural questions extend beyond
what we have proven; we collect them here.

\subsection{Dimension-locking as a general
phenomenon}\label{dimension-locking-as-a-general-phenomenon}

Theorem 4.1 established that $32$ of $33$ universal full-variable
fixed points at $d = 5$ are dimension-locked at $d = 5$ --- no
rank-$6$ rule at $d = 6$ is universal for them. The natural
generalization is:

\begin{conjecture}[7.1 (effective rank bound at cross-dimensional
universality)]
Let $F$ be a universal full-variable fixed point
at native digit length $d_F$, and suppose $K$ is a full-variable
rule at some $d \neq d_F$ that is universal for $F$. Then
$\mathrm{sv}_F(K) < d$ --- that is, the effective rank of $K$ at
$F$ is strictly less than the algebraic rank.
\end{conjecture}

Informally: when a fixed point appears at a non-native digit length, it
always does so through a rule whose effective rank at $F$ is strictly
reduced --- some of the rule's coefficients are ``absorbed'' by $F$'s
zero digits and do not contribute to the computation $K(F) = F$.

\textbf{Status.} Verified exhaustively at $d = 6$ for all $33$
$d = 5$ fixed points (Theorem 4.1 and the associated enumeration). For
$60714$, Theorem 5.2's construction automatically satisfies the
conjecture at every $d \geq 6$: the coefficient-preserving lifting has
$\mathrm{sv}_F = 4 < d$ at every $d \geq 6$. For $6174$, the
liftings at $d = 7, 8, 9$ similarly satisfy
$\mathrm{sv}_F \leq d - 3 < d$. The conjecture has not been tested at
$d > 6$ for the other $32$ dimension-locked $d = 5$ fixed points
or for fps native at $d = 6$.

\textbf{Formal verification of Conjecture 7.1 at $d \geq 7$ is open.}
Verifying it at a specific $d \geq 7$ for a specific fixed point $F$
reduces to exhaustive enumeration over rank-$d$ rules satisfying
$K(F_{(d)}) = F$, checking that all such rules have
$\mathrm{sv}_F < d$. The enumeration is tractable at $d = 7$
(several million rules per fixed point) and becomes computationally
demanding at $d \geq 8$.

\subsection{The structural invariant and dimensional
fate}\label{the-structural-invariant-and-dimensional-fate}

Observation 6.2 records the empirical unification of $60714$,
$60417$, and $6174$ through the $\mathrm{sum\_locked\_spans}$
invariant. The natural formal statement is:

\begin{conjecture}[7.2 (structural invariant controls dimensional
fate)]
For a universal full-variable fixed point $F$ with
native digit length $d_F$ and native rule $K_F$, the
cross-dimensional behavior of $F$ at $d > d_F$ is controlled by
$\mathrm{sum\_locked\_spans}(K_F)$ and the number of zero digits of
$F$ at $d$. Specifically:
\end{conjecture}

\begin{itemize}
\tightlist
\item
  \emph{If $\mathrm{sum\_locked\_spans}(K_F)$ is small relative to the
  number of absorbing positions at $d$, a coefficient-preserving
  lifting exists that is universal for $F$ at $d$.}
\item
  \emph{If the invariant is large relative to the absorbing capacity,
  either no lifting exists (algebraic obstruction) or liftings exist but
  no one is universal (dynamic obstruction).}
\end{itemize}

\textbf{Status.} Purely conjectural. The empirical pattern holds for
$60714$, $60417$, $6174$; whether it holds for arbitrary universal
full-variable fixed points at $d_F = 4, 5, 6$ has not been
systematically tested. A positive resolution would give a quantitative
criterion for predicting cross-dimensional transcendence from the native
rule alone, obviating the need for direct enumeration at each candidate
$d$.

\subsection{Other multiset clusters}\label{other-multiset-clusters}

\textbf{Question 7.1.} \emph{Which digit multisets support universal
full-variable fixed points at three or more consecutive native digit
lengths?}

The $\{7, 6, 4, 1\}$ thread is the only such example in our
classification at $d \leq 6$. The other $d = 4$ multiset
$\{2, 3, 5, 8\}$ produces fps at $d = 4$ but not at $d = 5, 6$.
Several $d = 5$ multisets have analogues at $d = 4$ or $d = 6$ but
not at both.

Extending the search to $d \geq 7$ would require either a full
classification at $d = 7$ (the $d = 7$ enumeration of
$9{,}344{,}160$ full-variable rules is computationally demanding; see
Appendix D) or a more structural argument for which multisets can
support cross-dimensional threads.

\textbf{Question 7.2.} \emph{Is there an analogue of $60714$ at higher
native digit lengths --- a fixed point $F$ with native $d_F = 7$ or
$d_F = 8$ that is universal at every $d \geq d_F$ via a
coefficient-preserving lifting?}

Such fixed points would constitute a second (or third, etc.)
dimension-transcendent attractor family. Observed data at $d_F = 6$
and $d_F = 7$ (Appendix D) contain several candidates that transcend
to $d + 1$, but none has been proven to lift uniformly to every higher
$d$ as $60714$ does.

\subsection{Other bases}\label{other-bases}

All results in this paper concern base $10$. The analogous framework
in base $b \geq 2$ would replace $10^{\pi_i}$ with $b^{\pi_i}$ in
the coefficient expansion of §2.2 and otherwise proceed identically.

\textbf{Question 7.3.} \emph{For each base $b$, does there exist a
universal full-variable fixed point at some native digit length that is
universal at every $d \geq d_F$ via a coefficient-preserving lifting?}

At base $10$, the answer is yes, demonstrated by $60714$. For small
bases ($b = 2, 3, \ldots$), the analogous classification would produce
different specific fixed points, and the structural mechanism of
coefficient-preserving lifting would apply with appropriate
modifications. The related literature \cite{Kay2024,Peterson2008}
addresses the classical Kaprekar rule at other bases; our framework has
not been systematically extended.

\subsection{\texorpdfstring{The $6174$ pattern at
$d \geq 10$}{The 6174 pattern at d \textbackslash geq 10}}\label{the-6174-pattern-at-d-geq-10}

Theorem 6.1 establishes $6174$'s behavior at $d = 4, 5, \ldots, 9$.
The pattern at $d = 8, 9$ (near-universal with $45$-input escape
class) is expected to continue at $d \geq 10$: the escape class
structure is $d$-independent (same multiset shape at each $d$),
while the admissible input set grows, so the basin fraction asymptotes
to $1$.

\textbf{Question 7.4.} \emph{Is $6174$ strict-universal at any
$d \geq 10$?}

Empirical extrapolation suggests no: at each $d \geq 8$, the
$45$-input escape class persists, and no coefficient-preserving
lifting eliminates it. A proof (or disproof) of this pattern at some
specific large $d$ would complete the $6174$ pattern to an arbitrary
horizon. At the structural level, proving that the $45$-input escape
class is genuinely stable under all coefficient-preserving liftings is
an algebraic question on the coefficient vectors' action on multisets of
the form $(X, X, X, X, 0, \ldots, 0)$.

\subsection{The depth of the escape class --- the Basin Density
Conjecture}\label{the-depth-of-the-escape-class-the-basin-density-conjecture}

Before stating the Basin Density Conjecture, we record a structural
observation that places the escape class $E_d$ in continuity with the
classical Kaprekar phenomenon at $d = 3, 4$.

\begin{observation}[7.6.0 (continuity with the classical case)]
Classical Kaprekar dynamics already exhibits an escape class at
$d = 3$ and $d = 4$. At $d = 3$, the ten repdigits
$\{0, 111, 222, \ldots, 999\}$ each map to $0$ in one application of
$K_0$, rather than to the nontrivial fixed point $495$. At
$d = 4$, the ten multiples of $1111$ --- namely
$\{0, 1111, 2222, \ldots, 9999\}$ --- each map to $0$ in one
application of $K_0$, rather than to $6174$. In both cases, the
escape class is precisely the set of inputs whose sorted-descending and
sorted-ascending forms are identical, and these are exactly the inputs
that the literature has conventionally treated as ``inadmissible.'' The
escape class $E_d$ defined by Definition 5.2.1 for the lifted operator
$K^{(d)}$ at $d \geq 7$ is the same phenomenon at higher digit
lengths: a structurally characterized set of inputs whose orbit reaches
the block-aligned multiset set $B_d$ and then collapses to $0$,
rather than reaching the nontrivial fixed point $60714$. The set
$B_d \cap A_d$ that anchors the escape class plays a role precisely
analogous to that of the repdigits at $d = 3$ and the multiples of
$1111$ at $d = 4$.
\end{observation}

What is new at $d \geq 7$ is therefore not the existence of an escape
class --- that is a generic feature of digit-permutation dynamics under
any sum-zero rule, and the classical case at $d = 3, 4$ has always had
one. What is new is the \emph{explicit closed form} for $|E_d^{(1)}|$
given by Lemma 5.5, the proof of one-step collapse for inputs in
$B_d \cap A_d$ (Lemma 5.1), and the structural conjecture below that
the asymptotic basin density $|E_d|/|A_d| \to 0$ across the lifting
family. The classical $d = 3, 4$ escape classes are small and
structurally simple precisely because the digit space is small; the
lifted-operator escape class at $d \geq 7$ is larger but is governed
by an analogous block-alignment criterion. Conjecture 7.6 asks whether
the escape class, while growing combinatorially with $d$, has
vanishing density.

Lemma 5.5 gives a closed form for $|E_d^{(1)}|$ --- the number of
step-$1$ collapse inputs at digit length $d$. The full escape class
$E_d$ is the backward orbit of $B_d \cap A_d$, and at $d \leq 11$
exhaustive enumeration confirms that every orbit in $E_d$ reaches
$0$ in at most $4$ iterations. We name the asymptotic claim:

\begin{conjecture}[7.6 (Basin Density Conjecture)]
The Basin
Density Conjecture for the $60714$ lifting family has two parts:
\end{conjecture}

\begin{enumerate}
\def\labelenumi{\arabic{enumi}.}
\tightlist
\item
  \emph{(\textbf{Density}.) $|E_d| / |A_d| \to 0$ as $d \to \infty$,
  i.e., the basin of $60714$ at digit length $d$ has asymptotic
  density $1$ within admissible inputs.}
\item
  \emph{(\textbf{Depth}.) There exists a uniform constant $T^*$ such
  that every orbit in $E_d$ reaches $0$ in at most $T^*$
  iterations of $K^{(d)}$, independently of $d$.}
\end{enumerate}

Combinatorially, $|E_d^{(1)}|$ grows polynomially in $d$
(specifically as $O(d^{d_0 - 1})$ via Lemma 5.5) while $|A_d|$ grows
as $O(d^9)$ with much larger constant; the basin fraction
$1 - |E_d|/|A_d|$ should approach $1$, supporting the Density part.
For the Depth part, exhaustive enumeration at $d \leq 11$ shows
$T^* \leq 4$ uniformly across the verified range, but no proof of a
uniform bound for general $d$ is provided here.

\subsection{Beyond universal full-variable fixed
points}\label{beyond-universal-full-variable-fixed-points}

The full-variable constraint $\mathrm{sv} = d$ excludes
dimension-agnostic fixed points like $45$, $495$, $450$ (§2.5).
These fixed points exhibit their own cross-dimensional behavior through
rank-reducing rules, which we have not addressed.

\textbf{Question 7.5.} \emph{Is there a useful cross-dimensional theory
for fixed points of rank-reducing rules ($\mathrm{sv} < d$)?}

The dimension-agnostic family is the simplest case: these fixed points
persist across all digit lengths by construction. A more interesting
question is whether there are non-full-variable rules whose
cross-dimensional behavior mirrors $60714$'s full-variable
universality --- say, rules with $\mathrm{sv} = d - 1$ that produce
universal fixed points at every $d$ in some range. We leave this open.

\subsection{Concluding remarks}\label{concluding-remarks}

The central result of this paper is Theorem 5.2: $60714$ is a
universal full-variable fixed point at every digit length $d \geq 5$,
under an explicit coefficient-preserving lifting. This is the first
explicitly constructed cross-dimensional persistence result in the
generalized Kaprekar family. The supporting results --- the
classifications of §3, the uniqueness of §4, and the $6174$
cross-dimensional pattern of §6 --- place this theorem in a structural
context that raises more questions than it answers. Conjectures 7.1 and
7.2 and Questions 7.1--7.5 and the Basin Density Conjecture (7.6)
together outline a program for understanding cross-dimensional behavior
in the generalized family: which fixed points transcend, under what
structural conditions, and with what mechanism.

The methodology of the paper --- exhaustive enumeration at each digit
length, combined with structural proof on the cases the enumeration
distinguishes --- has been effective at $d \leq 6$ for the full rule
space and at $d \leq 9$ for targeted verification of specific fps.
Extending this methodology to $d \geq 7$ for the full rule space, or
to $d \geq 10$ for $6174$, is the natural next step.

\section{Acknowledgments}\label{acknowledgments}

The mathematical content of this paper --- the framing of the organizing
question, the classifications at $d = 3, 4, 5, 6$, the
cross-dimensional cross-check, the coefficient-preserving lifting
framework, the proof of Theorem 5.2 (including the structural arguments
of Lemmas 5.1, 5.2, 5.3, 5.4 and the support lemmas of Appendix C), the
$\{7, 6, 4, 1\}$-thread analysis, and the open-question program of §7
--- was developed by the author, who is solely responsible for the
paper's content, including any remaining errors.

The development of the paper used substantial computer assistance, in
the spirit of the increasing role of computational tools in
number-theoretic research. Verification scripts were written and
executed for every numerical claim in the paper --- classifications,
basin counts, ladder-extension data, escape-class enumerations.
Structural arguments (most notably Lemma C.10's case enumeration and the
$d \geq 15$ algebraic argument in Appendix C.4) were tested against
worked examples and edge cases, a process that identified and corrected
several defects during development --- including a parenthetical
algebraic error in an earlier version of Lemma C.9, an enumeration-count
error in an earlier version of Lemma C.10, and an off-by-one error in a
closed-form formula in supplementary material. Prose was edited
iteratively. The cross-reference of the escape-class size sequence to
OEIS A000582 was identified through assisted literature search.

These computer-assisted activities used AI tools (Anthropic's Claude,
model versions Claude Opus 4.6 and 4.7, accessed during 2025 and 2026)
alongside conventional scripting in Python. The full development record,
the verification scripts, and the iterative review history are publicly
available at \url{https://github.com/clayelmore/Kaprekar-60714}.

Theorems, lemmas, propositions, conjectures, proofs, and computational
claims are the author's; computational tools were used as instruments
for verification, drafting, and review, not as authors.

\appendix

\section{\texorpdfstring{Classification Tables at
$d = 3, 4, 5, 6$}{Appendix A. Classification Tables at d = 3, 4, 5, 6}}\label{appendix-a.-classification-tables-at-d-3-4-5-6}

This appendix provides the complete enumerated classifications at digit
lengths $d = 3, 4, 5, 6$. Each section lists all universal
full-variable fixed points at the given $d$ along with representative
data for one universal rule per fixed point. Sign-flipped duplicates are
not listed separately (see Proposition 2.2).

The classifications were computed by exhaustive enumeration in Python.
Runtime on commodity hardware: negligible at $d = 3, 4$;
\textasciitilde5 seconds at $d = 5$; \textasciitilde5--10 minutes at
$d = 6$. Source code for the enumerations is available in the
supplementary materials.

\subsection{\texorpdfstring{Classification at
$d = 3$}{Classification at d = 3}}\label{classification-at-d-3}

\textbf{Full-variable rules enumerated:}
$3! \cdot D_3 = 6 \cdot 2 = 12$. \textbf{Admissible inputs
enumerated:} $720$ (non-repdigit, non-near-repdigit $3$-digit
integers). \textbf{Universal full-variable fixed points:} $0$.

No full-variable rule at $d = 3$ is universal for any fixed point. For
every rule, iteration from at least some admissible input enters a cycle
or a non-universal fixed point.

\textbf{Note on the dimension-agnostic fixed point $495$.} The
classical rule $K_0$ at $d = 3$ is universal for $495$, but has
algebraic rank $2$ (the middle-position coefficient vanishes; see
§1.1). $495$ is thus an sv$=2$ fixed point, not an sv$=3$ fixed
point, and is outside the scope of the full-variable classification.

\subsection{\texorpdfstring{Classification at
$d = 4$}{Classification at d = 4}}\label{classification-at-d-4}

\textbf{Full-variable rules enumerated:}
$4! \cdot D_4 = 24 \cdot 9 = 216$. \textbf{Admissible inputs
enumerated:} $615$ digit multisets ($9{,}630$ padded four-digit
strings). \textbf{Universal full-variable fixed points:} $4$.
\textbf{Universal full-variable rules (total, counting sign-flip
pairs):} $8$.

\begin{longtable}[]{@{}
  >{\centering\arraybackslash}p{(\columnwidth - 8\tabcolsep) * \real{0.2083}}
  >{\centering\arraybackslash}p{(\columnwidth - 8\tabcolsep) * \real{0.2083}}
  >{\centering\arraybackslash}p{(\columnwidth - 8\tabcolsep) * \real{0.2083}}
  >{\centering\arraybackslash}p{(\columnwidth - 8\tabcolsep) * \real{0.2083}}
  >{\raggedright\arraybackslash}p{(\columnwidth - 8\tabcolsep) * \real{0.1667}}@{}}
\toprule\noalign{}
\begin{minipage}[b]{\linewidth}\centering
Fixed point
\end{minipage} & \begin{minipage}[b]{\linewidth}\centering
Digit multiset
\end{minipage} & \begin{minipage}[b]{\linewidth}\centering
Sample $\pi$
\end{minipage} & \begin{minipage}[b]{\linewidth}\centering
Sample $\sigma$
\end{minipage} & \begin{minipage}[b]{\linewidth}\raggedright
Sample coefficient vector
\end{minipage} \\
\midrule\noalign{}
\endhead
\bottomrule\noalign{}
\endlastfoot
$1746$ & $\{1, 4, 6, 7\}$ & $(0, 3, 2, 1)$ & $(1, 2, 3, 0)$ &
$(-9, 900, -900, 9)$ \\
$2538$ & $\{2, 3, 5, 8\}$ & $(1, 0, 3, 2)$ & $(2, 3, 0, 1)$ &
$(-90, -999, 999, 90)$ \\
$5382$ & $\{2, 3, 5, 8\}$ & $(2, 1, 0, 3)$ & $(3, 0, 1, 2)$ &
$(-900, 9, -9, 900)$ \\
$6174$ & $\{1, 4, 6, 7\}$ & $(0, 1, 2, 3)$ & $(3, 2, 1, 0)$ &
$(-999, -90, 90, 999)$ \\
\end{longtable}

\textbf{Note on $6174$.} One of the two universal rules for $6174$
is the classical Kaprekar rule $K_0$ (descending minus ascending),
recovering Kaprekar's original 1949 result within our classification.

\textbf{Note on anagram clusters.} The four fps fall into two anagram
pairs: $\{1746, 6174\}$ (multiset $\{1, 4, 6, 7\}$) and
$\{2538, 5382\}$ (multiset $\{2, 3, 5, 8\}$). In each cluster, the
two fps sum to $7920$, a specific instance of reverse-pair identity at
$d = 4$.

\subsection{\texorpdfstring{Classification at
$d = 5$}{Classification at d = 5}}\label{classification-at-d-5}

\textbf{Full-variable rules enumerated:}
$5! \cdot D_5 = 120 \cdot 44 = 5{,}280$. \textbf{Admissible inputs
enumerated:} $99{,}540$. \textbf{Universal full-variable fixed
points:} $33$. \textbf{Universal full-variable rules (total, counting
sign-flip pairs):} $66$.

The $33$ fixed points grouped by digit multiset:

\begin{longtable}[]{@{}
  >{\raggedright\arraybackslash}p{(\columnwidth - 4\tabcolsep) * \real{0.3077}}
  >{\raggedright\arraybackslash}p{(\columnwidth - 4\tabcolsep) * \real{0.3077}}
  >{\centering\arraybackslash}p{(\columnwidth - 4\tabcolsep) * \real{0.3846}}@{}}
\toprule\noalign{}
\begin{minipage}[b]{\linewidth}\raggedright
Multiset
\end{minipage} & \begin{minipage}[b]{\linewidth}\raggedright
Fixed points
\end{minipage} & \begin{minipage}[b]{\linewidth}\centering
Cluster size
\end{minipage} \\
\midrule\noalign{}
\endhead
\bottomrule\noalign{}
\endlastfoot
$\{0, 0, 0, 4, 5\}$ & $54$ & $1$ \\
$\{0, 3, 3, 5, 7\}$ & $3753$ & $1$ \\
$\{0, 1, 4, 6, 7\}$ & $60714$, $60417$ & $2$ \\
$\{2, 3, 5, 8, 9\}$ & $28539$, $53928$, $58239$ & $3$ \\
$\{1, 2, 3, 4, 8\}$ & $18342$, $21834$, $24183$, $41832$,
$42183$ & $5$ \\
$\{1, 5, 6, 7, 8\}$ & $16578$, $16758$, $17685$, $65781$,
$67581$ & $5$ \\
$\{3, 4, 5, 7, 8\}$ & $37584$, $37854$, $38754$, $43758$,
$43785$, $43875$ & $6$ \\
$\{1, 2, 4, 5, 6\}$ & $12456$, $14562$, $15642$, $16524$,
$21456$, $24156$, $24561$, $41562$, $45612$, $45621$ &
$10$ \\
\end{longtable}

The complete list of $33$ fixed points in ascending order:

\[54,\ 3753,\ 12456,\ 14562,\ 15642,\ 16524,\ 16578,\ 16758,\ 17685,\ 18342,\]
\[21456,\ 21834,\ 24156,\ 24183,\ 24561,\ 28539,\ 37584,\ 37854,\ 38754,\ 41562,\]
\[41832,\ 42183,\ 43758,\ 43785,\ 43875,\ 45612,\ 45621,\ 53928,\ 58239,\ 60417,\]
\[60714,\ 65781,\ 67581.\]

\textbf{Representative rules at the $\{0, 1, 4, 6, 7\}$ cluster} (the
$\{7, 6, 4, 1\}$-thread of §6):

\begin{longtable}[]{@{}
  >{\centering\arraybackslash}p{(\columnwidth - 6\tabcolsep) * \real{0.2632}}
  >{\centering\arraybackslash}p{(\columnwidth - 6\tabcolsep) * \real{0.2632}}
  >{\centering\arraybackslash}p{(\columnwidth - 6\tabcolsep) * \real{0.2632}}
  >{\raggedright\arraybackslash}p{(\columnwidth - 6\tabcolsep) * \real{0.2105}}@{}}
\toprule\noalign{}
\begin{minipage}[b]{\linewidth}\centering
Fixed point
\end{minipage} & \begin{minipage}[b]{\linewidth}\centering
$\pi$
\end{minipage} & \begin{minipage}[b]{\linewidth}\centering
$\sigma$
\end{minipage} & \begin{minipage}[b]{\linewidth}\raggedright
Coefficient vector
\end{minipage} \\
\midrule\noalign{}
\endhead
\bottomrule\noalign{}
\endlastfoot
$60714$ & $(4, 1, 2, 3, 0)$ & $(2, 0, 1, 4, 3)$ &
$(9900,\; 9,\; 90,\; -9000,\; -999)$ \\
$60417$ & $(4, 1, 0, 3, 2)$ & $(0, 2, 1, 4, 3)$ &
$(9999,\; -90,\; -9,\; -9000,\; -900)$ \\
\end{longtable}

\textbf{Note on $54$.} The fixed point $54$ appears in the $d = 5$
full-variable classification under the rule $\pi = (1, 2, 4, 3, 0)$,
$\sigma = (2, 0, 3, 4, 1)$ with coefficient vector
$(-90, 99, 9000, -9000, -9)$. This rule is algebraically full-variable
at $d = 5$ ($\mathrm{sv} = 5$), but its effective rank at $F = 54$
is only $2$ --- three of the five coefficients land on zero digits of
$F$'s sorted-descending form $(5, 4, 0, 0, 0)$ and do not contribute
to $K(F) = F$. Thus
$\mathrm{sv}_F(K_{54}) = 2 < 5 = \mathrm{sv}(K_{54})$. This is
structurally analogous to the $495$ phenomenon at $d = 3$ (§1.1),
occurring at higher digit length within a genuinely full-variable rule.

\textbf{Note on $3753$.} Similarly, $3753$ has sorted-descending
form $(7, 5, 3, 3, 0)$ at $d = 5$ with one zero digit. Its effective
rank under its native rule is $\mathrm{sv}_F = 4 < 5$.

Both $54$ and $3753$ are dimension-locked at $d = 5 \to d = 6$:
the cross-check of §4 finds no universal full-variable rule for either
at $d = 6$.

\subsection{\texorpdfstring{Classification at
$d = 6$}{Classification at d = 6}}\label{classification-at-d-6}

\textbf{Full-variable rules enumerated:}
$6! \cdot D_6 = 720 \cdot 265 = 190{,}800$. \textbf{Admissible inputs
enumerated:} $4{,}905$ digit multisets ($999{,}450$ padded six-digit
strings). \textbf{Universal full-variable fixed points:} $506$.
\textbf{Universal full-variable rules (total, counting sign-flip
pairs):} $1{,}174$.

\subsubsection{Summary statistics}\label{summary-statistics}

The $506$ fixed points distribute by zero-digit count:

\begin{longtable}[]{@{}cc@{}}
\toprule\noalign{}
Zero-digit count & Fixed-point count \\
\midrule\noalign{}
\endhead
\bottomrule\noalign{}
\endlastfoot
$0$ & $205$ \\
$1$ & $240$ \\
$2$ & $53$ \\
$3$ & $8$ \\
\end{longtable}

And by digit sum (every universal fp at $d = 6$ has digit sum
divisible by $9$):

\begin{longtable}[]{@{}cc@{}}
\toprule\noalign{}
Digit sum & Fixed-point count \\
\midrule\noalign{}
\endhead
\bottomrule\noalign{}
\endlastfoot
$9$ & $8$ \\
$18$ & $156$ \\
$27$ & $244$ \\
$36$ & $96$ \\
$45$ & $2$ \\
\end{longtable}

\subsubsection{Distinguished fixed
points}\label{distinguished-fixed-points}

The following $d = 6$ universal full-variable fixed points are singled
out elsewhere in the paper:

\begin{longtable}[]{@{}
  >{\centering\arraybackslash}p{(\columnwidth - 4\tabcolsep) * \real{0.3571}}
  >{\centering\arraybackslash}p{(\columnwidth - 4\tabcolsep) * \real{0.3571}}
  >{\raggedright\arraybackslash}p{(\columnwidth - 4\tabcolsep) * \real{0.2857}}@{}}
\toprule\noalign{}
\begin{minipage}[b]{\linewidth}\centering
Fixed point
\end{minipage} & \begin{minipage}[b]{\linewidth}\centering
Digit multiset
\end{minipage} & \begin{minipage}[b]{\linewidth}\raggedright
Note
\end{minipage} \\
\midrule\noalign{}
\endhead
\bottomrule\noalign{}
\endlastfoot
$549945$ & $\{4, 4, 5, 5, 9, 9\}$ & zero-zero fp; algebraically
obstructed at $d = 6 \to d = 7$ \\
$60714$ & $\{0, 0, 1, 4, 6, 7\}$ & central result; cross-dimensional
fixed point at every $d \geq 5$, strict-universal at $d = 5, 6$
(Theorem 5.2) \\
$146070$ & $\{0, 0, 1, 4, 6, 7\}$ & in the
$\{7, 6, 4, 1\}$-thread; native at $d = 6$; transcendent to
$d = 7$ \\
$170460$ & $\{0, 0, 1, 4, 6, 7\}$ & in the
$\{7, 6, 4, 1\}$-thread; native at $d = 6$ \\
$607140$ & $\{0, 0, 1, 4, 6, 7\}$ & in the
$\{7, 6, 4, 1\}$-thread; native at $d = 6$; transcendent to
$d = 7$ \\
$631764$ & $\{1, 3, 4, 6, 6, 7\}$ & attractor at $d = 6$ under
classical rule with basin $0.0625$ \cite{Dahl2026} \\
\end{longtable}

\subsubsection{Fixed points with two or more zero
digits}\label{fixed-points-with-two-or-more-zero-digits}

The complete list of $506$ universal full-variable fixed points at
$d = 6$ is lengthy; we provide it in the supplementary materials as a
plain-text file (\texttt{d6\_fps.txt}). The enumeration is fully
reproducible from the source code.

Of particular relevance to §5 (the $60714$ theorem) and §6 (the
$\{7, 6, 4, 1\}$-thread) are the fps with at least two zero digits,
since these are the strata where dimension-transcendent behavior
empirically concentrates (cf.~Conjecture 7.2).

\textbf{The $8$ fps with $3$ zero digits.} All eight share digit
multiset $\{5, 2, 2, 0, 0, 0\}$ and are distinct integer arrangements
of the same digits:

\[252,\ \ 2520,\ \ 20025,\ \ 25200,\ \ 200025,\ \ 200250,\ \ 250200,\ \ 252000.\]

That all eight 3-zero fps belong to a single multiset is a notable
structural fact --- at higher zero-count strata, the multiset diversity
drops sharply.

\textbf{The $53$ fps with $2$ zero digits}, grouped by digit
multiset:

\begin{longtable}[]{@{}
  >{\raggedright\arraybackslash}p{(\columnwidth - 4\tabcolsep) * \real{0.3077}}
  >{\centering\arraybackslash}p{(\columnwidth - 4\tabcolsep) * \real{0.3846}}
  >{\raggedright\arraybackslash}p{(\columnwidth - 4\tabcolsep) * \real{0.3077}}@{}}
\toprule\noalign{}
\begin{minipage}[b]{\linewidth}\raggedright
Multiset
\end{minipage} & \begin{minipage}[b]{\linewidth}\centering
Cluster size
\end{minipage} & \begin{minipage}[b]{\linewidth}\raggedright
Sample fps
\end{minipage} \\
\midrule\noalign{}
\endhead
\bottomrule\noalign{}
\endlastfoot
$\{9, 9, 8, 1, 0, 0\}$ & $32$ &
$8919, 8991, 9189, 10899, 80919, 80991, \ldots$ \\
$\{5, 5, 4, 4, 0, 0\}$ & $7$ &
$4545, 44505, 54450, 445005, 445050, 504450, 544500$ \\
$\{9, 7, 1, 1, 0, 0\}$ & $5$ &
$7191, 17019, 700191, 701901, 719100$ \\
$\{7, 6, 4, 1, 0, 0\}$ & $4$ & $60714, 146070, 170460, 607140$ \\
$\{7, 7, 3, 1, 0, 0\}$ & $2$ & $37017, 707130$ \\
$\{9, 5, 2, 2, 0, 0\}$ & $1$ & $9225$ \\
$\{7, 6, 3, 2, 0, 0\}$ & $1$ & $67023$ \\
$\{6, 5, 5, 2, 0, 0\}$ & $1$ & $52605$ \\
\end{longtable}

The complete sorted list of all $53$ fps with $2$ zero digits:

\[4545, 7191, 8919, 8991, 9189, 9225, 10899, 17019, 37017, 44505, 52605, 54450,\]
\[60714, 67023, 80919, 80991, 81099, 89019, 90819, 90918, 91809, 98091, 100899,\]
\[108099, 109809, 146070, 170460, 180909, 189009, 190089, 190809, 445005,\]
\[445050, 504450, 544500, 607140, 700191, 701901, 707130, 719100, 809091,\]
\[809109, 810909, 890091, 900891, 901890, 908091, 908109, 908190, 908901,\]
\[909081, 910089, 910809.\]

\textbf{The $\{7, 6, 4, 1, 0, 0\}$ multiset cluster.} The four fps in
this cluster --- $60714, 146070, 170460, 607140$ --- share the
multiset structure of the central object of §5 (60714 padded to
$d = 6$). They are precisely the universal full-variable fps in this
multiset at $d = 6$, and the cross-dimensional behavior of all four is
documented at §6 and in Appendix D.

\textbf{Note on the unique $4$-zero fp.} §A.4.1 records one universal
full-variable fp at $d = 6$ with $4$ zero digits. Its identity is
recorded in \texttt{d6\_fps.txt} in the supplementary materials but is
not analyzed individually in this paper.

\textbf{Fixed points with $0$ or $1$ zero digit.} The remaining
$445$ fps ($205$ with $0$ zeros, $240$ with $1$ zero) are
listed in the supplementary file and are not analyzed individually here.
Cross-dimensional behavior at $d = 6 \to d = 7$ for fps with fewer
than $2$ zeros is empirically rare (cf.~Observation D.1).

\subsection{Cross-reference to §4 and
§6}\label{cross-reference-to-4-and-6}

The $33$ fixed points at $d = 5$ of §A.3 are the subject of §4's
cross-check: the exhaustive computation of which are dimension-locked at
$d = 5 \to d = 6$. The result ($17$ algebraic obstruction, $15$
dynamic obstruction, $1$ universal lifting to $d = 6$) is proven at
that section.

The $8$ fixed points with $3$ zero digits at $d = 6$ of §A.4.1 are
the subject of §6's transcendent-fp analysis at $d = 6 \to d = 7$
(empirical observations recorded in Appendix D).

\emph{End of Appendix A.}

\section{\texorpdfstring{Coefficient-preserving ladder for
$F = 60714$}{Appendix B. Coefficient-preserving ladder for F = 60714}}\label{appendix-b.-coefficient-preserving-ladder-for-f-60714}

For each digit length $d \geq 5$, the following tables give the
permutation-pair rule $K^{(d)}$ under which $F = 60714$ is a
universal fixed point, as constructed in §5.

\textbf{Conventions.} At digit length $d$, sorted-descending digit
positions are labeled $x_0, x_1, \ldots, x_{d-1}$ with
$x_0 \geq x_1 \geq \cdots \geq x_{d-1}$. We assign letters
$a, b, c, d, e, f, g, \ldots$ to these positions in order, so
$a = x_0$, $b = x_1$, etc. The rule
$K = \pi \cdot x - \sigma \cdot x$ is written as two letter strings:
reading each string left-to-right, the $k$-th letter indicates which
sorted-descending position contributes the digit at place $10^{d-1-k}$
of the resulting integer. The coefficient vector is
$c_i = 10^{\pi_i} - 10^{\sigma_i}$ for $i = 0, 1, \ldots, d-1$.

\textbf{Fixed-point verification.} At every $d$, the assignment
$a = 7, b = 6, c = 4, d = 1$, and every subsequent letter $= 0$
corresponds to $F = 60714$'s sorted-descending form (padded to $d$
digits). Under this assignment, $\pi \cdot x$ evaluates to $71460$
and $\sigma \cdot x$ evaluates to $10746$, so
$K(F) = |71460 - 10746| = 60714$ at every $d$.

\subsection{\texorpdfstring{Ladder rules through
$d = 20$}{Ladder rules through d = 20}}\label{ladder-rules-through-d-20}

\begin{longtable}[]{@{}
  >{\centering\arraybackslash}p{(\columnwidth - 12\tabcolsep) * \real{0.1562}}
  >{\raggedright\arraybackslash}p{(\columnwidth - 12\tabcolsep) * \real{0.1250}}
  >{\raggedright\arraybackslash}p{(\columnwidth - 12\tabcolsep) * \real{0.1250}}
  >{\raggedright\arraybackslash}p{(\columnwidth - 12\tabcolsep) * \real{0.1250}}
  >{\centering\arraybackslash}p{(\columnwidth - 12\tabcolsep) * \real{0.1562}}
  >{\centering\arraybackslash}p{(\columnwidth - 12\tabcolsep) * \real{0.1562}}
  >{\centering\arraybackslash}p{(\columnwidth - 12\tabcolsep) * \real{0.1562}}@{}}
\toprule\noalign{}
\begin{minipage}[b]{\linewidth}\centering
$d$
\end{minipage} & \begin{minipage}[b]{\linewidth}\raggedright
Ladder
\end{minipage} & \begin{minipage}[b]{\linewidth}\raggedright
$\pi$ (letters)
\end{minipage} & \begin{minipage}[b]{\linewidth}\raggedright
$\sigma$ (letters)
\end{minipage} & \begin{minipage}[b]{\linewidth}\centering
$\pi \cdot F$
\end{minipage} & \begin{minipage}[b]{\linewidth}\centering
$\sigma \cdot F$
\end{minipage} & \begin{minipage}[b]{\linewidth}\centering
$K(F)$
\end{minipage} \\
\midrule\noalign{}
\endhead
\bottomrule\noalign{}
\endlastfoot
5 & odd (native) & \texttt{adcbe} & \texttt{deacb} & 71460 & 10746 &
60714 \\
6 & even (split) & \texttt{eadcbf} & \texttt{fdeacb} & 71460 & 10746 &
60714 \\
7 & odd & \texttt{fgadcbe} & \texttt{gfdeacb} & 71460 & 10746 & 60714 \\
8 & even & \texttt{hgeadcbf} & \texttt{ghfdeacb} & 71460 & 10746 &
60714 \\
9 & odd & \texttt{hifgadcbe} & \texttt{ihgfdeacb} & 71460 & 10746 &
60714 \\
10 & even & \texttt{jihgeadcbf} & \texttt{ijghfdeacb} & 71460 & 10746 &
60714 \\
11 & odd & \texttt{jkhifgadcbe} & \texttt{kjihgfdeacb} & 71460 & 10746 &
60714 \\
12 & even & \texttt{lkjihgeadcbf} & \texttt{klijghfdeacb} & 71460 &
10746 & 60714 \\
13 & odd & \texttt{lmjkhifgadcbe} & \texttt{mlkjihgfdeacb} & 71460 &
10746 & 60714 \\
14 & even & \texttt{nmlkjihgeadcbf} & \texttt{mnklijghfdeacb} & 71460 &
10746 & 60714 \\
15 & odd & \texttt{nolmjkhifgadcbe} & \texttt{onmlkjihgfdeacb} & 71460 &
10746 & 60714 \\
16 & even & \texttt{ponmlkjihgeadcbf} & \texttt{opmnklijghfdeacb} &
71460 & 10746 & 60714 \\
17 & odd & \texttt{pqnolmjkhifgadcbe} & \texttt{qponmlkjihgfdeacb} &
71460 & 10746 & 60714 \\
18 & even & \texttt{rqponmlkjihgeadcbf} & \texttt{qropmnklijghfdeacb} &
71460 & 10746 & 60714 \\
19 & odd & \texttt{rspqnolmjkhifgadcbe} & \texttt{srqponmlkjihgfdeacb} &
71460 & 10746 & 60714 \\
20 & even & \texttt{tsrqponmlkjihgeadcbf} &
\texttt{stqropmnklijghfdeacb} & 71460 & 10746 & 60714 \\
\end{longtable}

\subsection{\texorpdfstring{Coefficient vectors through
$d = 12$}{Coefficient vectors through d = 12}}\label{coefficient-vectors-through-d-12}

The coefficient vectors are the same data as the letter strings in Table
B.1, presented numerically. The structure is uniform along each ladder:
every rule begins with the four locked coefficients
$(9900, 9, 90, -9000)$ at $F$'s nonzero-digit positions, and every
higher-$d$ rule extends the previous one by appending one zero-sum
pair.

The first five (odd ladder) or six (even ladder) coefficients are:

\begin{longtable}[]{@{}cl@{}}
\toprule\noalign{}
Ladder & Locked + root coefficients \\
\midrule\noalign{}
\endhead
\bottomrule\noalign{}
\endlastfoot
odd, $d = 5$ & $(9900,\ 9,\ 90,\ -9000,\ -999)$ \\
even, $d = 6$ & $(9900,\ 9,\ 90,\ -9000,\ 99000,\ -99999)$ \\
\end{longtable}

For each subsequent step along a ladder, the rule at $d$ extends the
rule at $d - 2$ by appending the zero-sum pair
$(\pm 9 \cdot 10^{d-2},\ \mp 9 \cdot 10^{d-2})$. Concretely:

\begin{longtable}[]{@{}
  >{\centering\arraybackslash}p{(\columnwidth - 4\tabcolsep) * \real{0.3571}}
  >{\centering\arraybackslash}p{(\columnwidth - 4\tabcolsep) * \real{0.3571}}
  >{\raggedright\arraybackslash}p{(\columnwidth - 4\tabcolsep) * \real{0.2857}}@{}}
\toprule\noalign{}
\begin{minipage}[b]{\linewidth}\centering
$d$
\end{minipage} & \begin{minipage}[b]{\linewidth}\centering
Ladder
\end{minipage} & \begin{minipage}[b]{\linewidth}\raggedright
Appended pair $(c_{d-2}, c_{d-1})$
\end{minipage} \\
\midrule\noalign{}
\endhead
\bottomrule\noalign{}
\endlastfoot
7 & odd & $(9 \cdot 10^5,\ -9 \cdot 10^5) = (900\,000,\ -900\,000)$ \\
8 & even &
$(-9 \cdot 10^6,\ 9 \cdot 10^6) = (-9\,000\,000,\ 9\,000\,000)$ \\
9 & odd &
$(9 \cdot 10^7,\ -9 \cdot 10^7) = (90\,000\,000,\ -90\,000\,000)$ \\
10 & even &
$(-9 \cdot 10^8,\ 9 \cdot 10^8) = (-900\,000\,000,\ 900\,000\,000)$ \\
11 & odd &
$(9 \cdot 10^9,\ -9 \cdot 10^9) = (9 \cdot 10^9,\ -9 \cdot 10^9)$ \\
12 & even &
$(-9 \cdot 10^{10},\ 9 \cdot 10^{10}) = (-9 \cdot 10^{10},\ 9 \cdot 10^{10})$ \\
\end{longtable}

The full coefficient vector at any $d$ is the locked-plus-root prefix
followed by the chain of appended pairs from $d - 2$ down to the
ladder root. The construction is stated formally in §B.3.

\subsection{The construction, stated
formally}\label{the-construction-stated-formally}

The tables above are the output of a simple recursive recipe that
generalizes to every $d \geq 5$:

\textbf{Base cases.}\\
Odd ladder root: at $d = 5$, $\pi = (4, 1, 2, 3, 0)$,
$\sigma = (2, 0, 1, 4, 3)$, coefficient vector
$(9900, 9, 90, -9000, -999)$.\\
Even ladder root: at $d = 6$, $\pi = (4, 1, 2, 3, 5, 0)$,
$\sigma = (2, 0, 1, 4, 3, 5)$, coefficient vector
$(9900, 9, 90, -9000, 99000, -99999)$.

\textbf{Inductive step.} Given the rule at $d$ on a ladder, the rule
at $d + 2$ on the same ladder is obtained by appending two new
positions:

\[\pi_{d} = \begin{cases} d + 1 & \text{(odd ladder)} \\ d & \text{(even ladder)} \end{cases}, \quad \pi_{d+1} = \begin{cases} d & \text{(odd ladder)} \\ d + 1 & \text{(even ladder)} \end{cases},\]

\[\sigma_{d} = \begin{cases} d & \text{(odd ladder)} \\ d + 1 & \text{(even ladder)} \end{cases}, \quad \sigma_{d+1} = \begin{cases} d + 1 & \text{(odd ladder)} \\ d & \text{(even ladder)} \end{cases}.\]

In both cases, the two appended coefficients
$c_d = 10^{\pi_d} - 10^{\sigma_d}$ and
$c_{d+1} = 10^{\pi_{d+1}} - 10^{\sigma_{d+1}}$ satisfy
$c_d + c_{d+1} = 0$ --- they form a \textbf{zero-sum pair}. The only
difference between ladders is the sign convention on the new pair.

\subsection{\texorpdfstring{Demonstration at
$d = 100$}{Demonstration at d = 100}}\label{demonstration-at-d-100}

To show the recipe is not merely finite-table-bound but scales
mechanically, we give the rule at $d = 100$ on the even ladder.

\textbf{Structure of the coefficient vector at $d = 100$.}

\begin{itemize}
\tightlist
\item
  Positions $0$ through $3$: $(9900,\ 9,\ 90,\ -9000)$. These are
  the coefficients at $F$'s four nonzero-digit positions, preserved
  verbatim from the native rule at $d = 5$. (Locked nonzero
  coefficients.)
\item
  Positions $4, 5$: $(99000,\ -99999)$. These are the even-ladder
  ``split'' coefficients fixed at the even root ($d = 6$).
\item
  Positions $6, 7, 8, \ldots, 99$: 47 zero-sum pairs, one appended at
  each even-ladder step from $d = 8$ to $d = 100$.
\end{itemize}

The $k$-th appended pair (for $k = 1, 2, \ldots, 47$) consists of

\[c_{2k + 4} = -10^{2k + 4} \cdot 9 + (\text{higher-order terms}) = -\,\underbrace{\overbrace{9 \cdot 10^{2k + 4}}^{\text{negative}}}_{\text{(see table)}}, \qquad c_{2k + 5} = +9 \cdot 10^{2k + 4}.\]

Concretely, the first four appended pairs at $d = 10, 12, 14, 16$ are:

\begin{longtable}[]{@{}ccl@{}}
\toprule\noalign{}
Pair index $k$ & Positions & $(c_i, c_{i+1})$ \\
\midrule\noalign{}
\endhead
\bottomrule\noalign{}
\endlastfoot
1 & (6, 7) & $(-9 \times 10^6,\ +9 \times 10^6)$ \\
2 & (8, 9) & $(-9 \times 10^8,\ +9 \times 10^8)$ \\
3 & (10, 11) & $(-9 \times 10^{10},\ +9 \times 10^{10})$ \\
4 & (12, 13) & $(-9 \times 10^{12},\ +9 \times 10^{12})$ \\
$\vdots$ & $\vdots$ & $\vdots$ \\
47 & (98, 99) & $(-9 \times 10^{98},\ +9 \times 10^{98})$ \\
\end{longtable}

\textbf{Fixed-point verification at $d = 100$.} The digits of
$F = 60714$ padded to $100$ digits, sorted in descending order, are
$a = 7$, $b = 6$, $c = 4$, $d = 1$, and 96 trailing zeros. Under
this assignment, every appended coefficient $c_i$ for $i \geq 4$
multiplies a zero digit and contributes nothing to $K(F)$. The
computation reduces to

\[K(F) = | (9900)(7) + (9)(6) + (90)(4) + (-9000)(1) | = |\,69300 + 54 + 360 - 9000\,| = 60714.\]

This is the same arithmetic at every $d \geq 5$. \textbf{The
coefficient-preserving lifting makes the fixed-point equation trivial at
every $d$} because the structural content of the lifting lives
entirely in the zero-digit positions, which absorb the new coefficients
without affecting $K(F)$.

The nontrivial content of Theorem 3 is not that $K(F) = F$ at every
$d$ --- that is automatic from the construction. The nontrivial
content is that every admissible orbit outside the escape class $E_d$
at every $d \geq 5$ reaches $F$, which is the work done by Lemmas
5.1 and 5.2.

\section{Support Lemmas for the Proof of Theorem
5.2}\label{appendix-c.-support-lemmas-for-the-proof-of-theorem-5.2}

This appendix provides the technical support lemmas referenced in §5.
The organization mirrors the proof of Theorem 5.2: §C.1 fixes setup,
§C.2 proves core non-negativity (Lemma 5.3), §C.3 proves core upper
bound (Lemma 5.4), §C.4 proves the reaching-time bound on the odd
ladder, §C.5 addresses the even ladder, and §C.6 establishes the
admissible-projection lemma.

\subsection{Setup}\label{setup}

Throughout, $(x_0, x_1, \ldots, x_{d-1})$ denotes a sorted-descending
digit sequence with $x_0 \geq x_1 \geq \cdots \geq x_{d-1}$ and each
$x_i \in \{0, 1, \ldots, 9\}$. The \textbf{native core contribution}
(Definition 5.4) is

\[\mathrm{core}(x) \;=\; 9900 \, x_0 + 9 \, x_1 + 90 \, x_2 - 9000 \, x_3 - 999 \, x_4.\]

\subsection{Core non-negativity (Proof of Lemma
5.3)}\label{core-non-negativity-proof-of-lemma-5.3}

\begin{lemma}[5.3 (restated)]
For every sorted-descending
sequence $(x_0, x_1, x_2, x_3, x_4)$, $\mathrm{core}(x) \geq 0$.
\end{lemma}

\begin{proof}
We proceed by case analysis on the values of $x_3$ and
$x_4$.

\textbf{Case 1: $x_3 = x_4 = 0$.} Then
$\mathrm{core}(x) = 9900 x_0 + 9 x_1 + 90 x_2 \geq 0$, with equality
iff $x_0 = x_1 = x_2 = 0$.

\textbf{Case 2: $x_3 \geq 1$, $x_4 = 0$.} By sorted-descending,
$x_0 \geq x_3 \geq 1$. Writing
$\mathrm{core}(x) = 9000(x_0 - x_3) + 900 x_0 + 9 x_1 + 90 x_2$, each
term is non-negative, so $\mathrm{core}(x) \geq 0$.

\textbf{Case 3: $x_3 \geq 1$ and $x_4 \geq 1$.} We use
$x_2 \geq x_3$ and $x_1 \geq x_4$:

\[\mathrm{core}(x) = 9900 x_0 + 9 x_1 + 90 x_2 - 9000 x_3 - 999 x_4.\]

Replace $90 x_2 \geq 90 x_3$ (valid since $x_2 \geq x_3$):

\[\mathrm{core}(x) \geq 9900 x_0 + 9 x_1 + 90 x_3 - 9000 x_3 - 999 x_4 = 9900 x_0 + 9 x_1 - 8910 x_3 - 999 x_4.\]

Replace $9 x_1 \geq 9 x_4$:

\[\mathrm{core}(x) \geq 9900 x_0 - 8910 x_3 - 999 x_4 + 9 x_4 = 9900 x_0 - 8910 x_3 - 990 x_4.\]

Replace $x_4 \leq x_3$:

\[\mathrm{core}(x) \geq 9900 x_0 - 8910 x_3 - 990 x_3 = 9900 x_0 - 9900 x_3 = 9900 (x_0 - x_3) \geq 0,\]

using $x_0 \geq x_3$ in the last step.

All three cases cover the admissible possibilities (Case 4 with
$x_3 = 0$ and $x_4 \geq 1$ is impossible by sorted-descending). In
each, $\mathrm{core}(x) \geq 0$.
\end{proof}

\begin{remark}[C.1]
The minimum $\mathrm{core}(x) = 0$ is achieved at
$x = (0, 0, 0, 0, 0)$ (repdigit zero). Direct enumeration over all
$2{,}002$ sorted-descending sequences in $\{0, \ldots, 9\}^5$
confirms the minimum is exactly $0$.
\end{remark}

\subsection{Core upper bound (Proof of Lemma
5.4)}\label{core-upper-bound-proof-of-lemma-5.4}

\begin{lemma}[5.4 (restated)]
For every sorted-descending
sequence $(x_0, x_1, x_2, x_3, x_4)$,
$\mathrm{core}(x) \leq 89{,}991 < 10^5$.
\end{lemma}

\begin{proof}
The negative terms $-9000 x_3 - 999 x_4$ contribute at
most $0$ (when $x_3 = x_4 = 0$), so:

\[\mathrm{core}(x) \leq 9900 x_0 + 9 x_1 + 90 x_2 \leq 9900 \cdot 9 + 9 \cdot 9 + 90 \cdot 9 = 89{,}100 + 81 + 810 = 89{,}991.\]

The maximum is achieved at
$(x_0, x_1, x_2, x_3, x_4) = (9, 9, 9, 0, 0)$.
\end{proof}

\begin{corollary}[C.2 (computational verification)]
Direct
enumeration over all $2{,}002$ sorted-descending sequences
$(x_0, \ldots, x_4) \in \{0, \ldots, 9\}^5$ confirms
$\mathrm{core}(x) \in [0, 89{,}991]$ with minimum $0$ at
$(0, 0, 0, 0, 0)$ and maximum $89{,}991$ at $(9, 9, 9, 0, 0)$.
Runtime: milliseconds.
\end{corollary}

\subsection{\texorpdfstring{The $T_d$ reaching-time bound (odd
ladder)}{The T\_d reaching-time bound (odd ladder)}}\label{the-t_d-reaching-time-bound-odd-ladder}

The main technical claim underlying Proposition 5.3 rests on a
\textbf{block-structure decomposition} of the lifted rule's output.

\subsection{Block-structure
decomposition}\label{block-structure-decomposition}

At odd digit length $d \geq 7$, the odd-ladder rule
$K^{(d)}_{60714}$ has coefficient vector

\[c^{(d)} = \underbrace{(9900, 9, 90, -9000, -999)}_{\text{native, positions } 0..4} \; \| \; \underbrace{(+p_1, -p_1)}_{\text{pair }1, \text{positions } 5, 6} \; \| \; \cdots \; \| \; \underbrace{(+p_M, -p_M)}_{\text{pair }M, \text{positions } d-2, d-1},\]

where $M = (d-5)/2$ and $p_k = 9 \cdot 10^{3 + 2k}$ is the
pair-$k$ magnitude. Pair $k$ occupies positions
$(3 + 2k, 4 + 2k)$.

\textbf{Pair contribution.} For input $(x_0, \ldots, x_{d-1})$, each
pair $k$ contributes

\[p_k \cdot x_{3 + 2k} - p_k \cdot x_{4 + 2k} = p_k \, \delta_k, \qquad \text{where } \delta_k := x_{3 + 2k} - x_{4 + 2k} \in \{0, 1, \ldots, 9\}.\]

By the sorted-descending constraint, each $\delta_k \geq 0$.

\begin{lemma}[C.6 (cliff-sequence bound)]
For any
sorted-descending sequence $(x_0, \ldots, x_{d-1})$ with each
$x_i \in \{0, \ldots, 9\}$,
\end{lemma}

\[\sum_{k = 1}^{M} \delta_k \;=\; \sum_{k = 1}^{M} (x_{3 + 2k} - x_{4 + 2k}) \;\leq\; 9.\]

\begin{proof}
Each $\delta_k$ is a non-negative adjacent-position
difference within the sorted-descending sequence. The sum is a subset of
the telescoping total
$\sum_{i=3}^{d-2} (x_i - x_{i+1}) = x_3 - x_{d-1}$. Since all
differences are non-negative and $x_3 - x_{d-1} \leq 9$, the selected
subsum is bounded by the full telescoping:

\[\sum_{k = 1}^{M} \delta_k \leq x_3 - x_{d-1} \leq 9.\]
\end{proof}

\subsection{\texorpdfstring{Non-negativity of
$K^{(d)}(x)$}{Non-negativity of K\^{}\{(d)\}(x)}}\label{non-negativity-of-kdx}

\begin{lemma}[C.7]
For every sorted-descending
$(x_0, \ldots, x_{d-1})$ and every odd $d \geq 7$,
\end{lemma}

\[K^{(d)}(x) \;=\; \mathrm{core}(x) + \sum_{k=1}^{M} p_k \, \delta_k.\]

\emph{In particular, $K^{(d)}(x) \geq 0$, so the absolute value
operation in the rule definition is vacuous on sorted-descending
inputs.}

\begin{proof}
Direct substitution into the expanded rule gives

\[K^{(d)}(x) \;=\; \left| \sum_{i = 0}^{d-1} c_i \, x_i \right| \;=\; \left| \mathrm{core}(x) + \sum_{k=1}^{M} p_k \, \delta_k \right|.\]

By Lemma 5.3, $\mathrm{core}(x) \geq 0$. Each
$p_k \, \delta_k \geq 0$ (since $p_k > 0$ and $\delta_k \geq 0$).
Hence the argument is non-negative and the absolute-value bars can be
dropped.
\end{proof}

\subsection{The reaching-time theorem}\label{the-reaching-time-theorem}

\begin{lemma}[C.3 (one-step $T_d$ closure on the odd ladder at
$d \geq 15$)]
Let $d \geq 15$ be odd. For every admissible
input $n \in A_d$ with sorted-descending form
$(x_0, \ldots, x_{d-1})$,
\end{lemma}

\[K^{(d)}_{60714}(n) \in T_d.\]

\emph{Equivalently, the padded $d$-digit representation of
$K^{(d)}(n)$ has at least two zero digits.}

\begin{proof}
By Lemma C.7,

\[K^{(d)}(x) \;=\; \mathrm{core}(x) + \sum_{k=1}^{M} 9 \cdot 10^{3 + 2k} \, \delta_k.\]

\textbf{Decimal-block structure.} Each term
$9 \cdot 10^{3 + 2k} \delta_k$ occupies decimal positions
$[3 + 2k, 4 + 2k]$ as a two-digit value. Writing
$9 \delta_k = 10 u_k + v_k$ for $u_k, v_k \in \{0, \ldots, 9\}$:

\[9 \cdot 10^{3 + 2k} \delta_k = u_k \cdot 10^{4 + 2k} + v_k \cdot 10^{3 + 2k}.\]

The digit pairs $(u_k, v_k)$ for each $\delta_k$:

\begin{longtable}[]{@{}cccc@{}}
\toprule\noalign{}
$\delta_k$ & $9 \delta_k$ & $(u_k, v_k)$ & Zeros \\
\midrule\noalign{}
\endhead
\bottomrule\noalign{}
\endlastfoot
$0$ & $0$ & $(0, 0)$ & $2$ \\
$1$ & $9$ & $(0, 9)$ & $1$ \\
$2$ & $18$ & $(1, 8)$ & $0$ \\
$3$ & $27$ & $(2, 7)$ & $0$ \\
$4$ & $36$ & $(3, 6)$ & $0$ \\
$5$ & $45$ & $(4, 5)$ & $0$ \\
$6$ & $54$ & $(5, 4)$ & $0$ \\
$7$ & $63$ & $(6, 3)$ & $0$ \\
$8$ & $72$ & $(7, 2)$ & $0$ \\
$9$ & $81$ & $(8, 1)$ & $0$ \\
\end{longtable}

\textbf{Disjoint blocks.} Pair $k$'s block occupies decimals
$[3 + 2k, 4 + 2k]$. Since the block value
$9 \delta_k \leq 81 < 100$, it fits entirely within two decimal
positions and does not propagate carries to higher blocks. Similarly,
$\mathrm{core}(x) \leq 89{,}991 < 10^5$ (Lemma 5.4) fits within
decimals $[0, 4]$ with at most one carry into decimal $5$. That
carry combines additively with pair $1$'s block at positions
$[5, 6]$: the combined value is at most
$u_1 \cdot 10 + v_1 + 1 = 9 \delta_1 + 1 \leq 82$, still a two-digit
value, so no further carry to decimal $7$. Blocks for $k \geq 2$
retain their $(u_k, v_k)$ form exactly.

\textbf{Counting zeros in the block region.} Define

\[Z_0 := \bigl|\{k : \delta_k = 0\}\bigr|, \qquad Z_1 := \bigl|\{k : \delta_k = 1\}\bigr|, \qquad Z_{2+} := \bigl|\{k : \delta_k \geq 2\}\bigr|.\]

Then $Z_0 + Z_1 + Z_{2+} = M$, and the number of zero digits
contributed by blocks is at least $2 Z_0 + Z_1$ (each $\delta_k = 0$
block contributes $2$ zeros; each $\delta_k = 1$ block contributes
$1$ zero from $u_k = 0$).

From Lemma C.6, $\sum_k \delta_k \leq 9$. Since $\delta_k \geq 2$
for $k \in Z_{2+}$:

\[2 Z_{2+} + Z_1 \leq \sum_k \delta_k \leq 9.\]

Therefore:

\[2 Z_0 + Z_1 \;=\; 2(M - Z_1 - Z_{2+}) + Z_1 \;=\; 2M - Z_1 - 2 Z_{2+} \;\geq\; 2M - 9.\]

\textbf{At odd $d \geq 17$: $M \geq 6$, so $2 Z_0 + Z_1 \geq 3$.}
This alone gives at least $3$ zero digits in the block region, more
than enough for $K^{(d)}(x) \in T_d$.

\textbf{At $d = 15$ (odd): $M = 5$, so $2 Z_0 + Z_1 \geq 1$.} This
gives at least one zero digit from the block region, but the second zero
requires a refined analysis. We complete the argument by partitioning
according to the value of $x_4$.

\begin{lemma}[C.7 (sharpened cliff-sum bound for $d = 15$)]
For
sorted-descending $(x_0, \ldots, x_{14})$ with
$x_i \in \{0, \ldots, 9\}$,
\end{lemma}

\[\sum_{k = 1}^{5} \delta_k \;\leq\; x_4.\]

\emph{Proof.} Each $\delta_k = x_{3 + 2k} - x_{4 + 2k}$ for
$k = 1, \ldots, 5$ is a non-negative adjacent-position difference
among positions $5, 6, \ldots, 14$. The full telescoping over these
positions is

\[\sum_{i = 5}^{13} (x_i - x_{i+1}) = x_5 - x_{14}.\]

Our sum $\sum_k \delta_k$ picks out the five differences
$(x_5 - x_6), (x_7 - x_8), (x_9 - x_{10}), (x_{11} - x_{12}), (x_{13} - x_{14})$,
a strict subset of these nine. Since every omitted difference is
non-negative,

\[\sum_{k=1}^{5} \delta_k \;\leq\; x_5 - x_{14} \;\leq\; x_5 \;\leq\; x_4. \qquad \square\]

\begin{lemma}[C.8 (core-bound when $x_4 \in \{8, 9\}$)]
If
$x_4 \geq 8$ in sorted-descending $(x_0, \ldots, x_4)$, then
$\mathrm{core}(x) < 10^4$.
\end{lemma}

\emph{Proof.} By sorted-descending, $x_4 \geq 8$ forces
$x_0, x_1, x_2, x_3 \in \{8, 9\}$. We compute:

\[\mathrm{core}(x) = 9900 x_0 + 9 x_1 + 90 x_2 - 9000 x_3 - 999 x_4 = 9000(x_0 - x_3) + 900 x_0 + 9 x_1 + 90 x_2 - 999 x_4.\]

Since $x_0, x_3 \in \{8, 9\}$ with $x_0 \geq x_3$, we have
$x_0 - x_3 \in \{0, 1\}$.

\emph{Case $x_0 = x_3$:} The first term vanishes. Then

\[\mathrm{core}(x) = 900 x_0 + 9 x_1 + 90 x_2 - 999 x_4.\]

The maximum over $x_0, x_1, x_2 \in \{8, 9\}$ and $x_4 \in \{8, 9\}$
with $x_4 \leq x_3 = x_0$ is achieved at $x_0 = x_1 = x_2 = 9$,
$x_4 = 8$:
$900 \cdot 9 + 9 \cdot 9 + 90 \cdot 9 - 999 \cdot 8 = 8100 + 81 + 810 - 7992 = 999$.
So $\mathrm{core}(x) \leq 999 < 10^3$.

\emph{Case $x_0 = x_3 + 1$:} Then $x_0 = 9$ and $x_3 = x_4 = 8$
(forcing $x_4 = 8$ since $x_4 \leq x_3$). We get

\[\mathrm{core}(x) = 9000 + 900 \cdot 9 + 9 x_1 + 90 x_2 - 999 \cdot 8 = 9108 + 9 x_1 + 90 x_2\]

for $x_1, x_2 \in \{8, 9\}$ with $x_1 \geq x_2 \geq 8$. The minimum
is $9108 + 72 + 720 = 9900$ and the maximum is
$9108 + 81 + 810 = 9999$. So $\mathrm{core}(x) \in [9900, 9999]$,
and in particular $\mathrm{core}(x) < 10^4$.

In both cases, $\mathrm{core}(x) < 10^4$.
\end{proof}

We now complete the $d = 15$ closure argument.

\textbf{Three cases for the second zero digit.}

\emph{Case A: $x_4 = 0$.} By sorted-descending,
$x_5 = x_6 = \cdots = x_{14} = 0$. All five $\delta_k = 0$, so every
block contributes $(0, 0)$ at decimals
$[5, 6], [7, 8], \ldots, [13, 14]$. The block region produces $10$
zero digits at decimals $5$ through $14$. Even after accounting for
at most one carry from the core into decimal $5$, decimals $6$
through $14$ remain zero --- that is $9$ zero digits, far more than
the $2$ required.

\emph{Case B1: $1 \leq x_4 \leq 7$.} By Lemma C.7,
$\sum_k \delta_k \leq x_4 \leq 7$. Since $\delta_k \geq 2$ for
$k \in Z_{2+}$,

\[2 Z_{2+} + Z_1 \leq \sum_k \delta_k \leq 7,\]

which gives

\[2 Z_0 + Z_1 = 2M - Z_1 - 2 Z_{2+} \geq 10 - 7 = 3.\]

So the block region contributes at least $3$ zero digits.

\emph{Case B2: $x_4 \in \{8, 9\}$.} By Lemma C.8,
$\mathrm{core}(x) < 10^4$, so the padded core occupies only decimals
$0$ through $3$, and decimal $4$ is a zero digit in the core's
contribution to $K^{(15)}(x)$. Combined with at least one zero digit
from the block region (the $2 Z_0 + Z_1 \geq 1$ bound), this gives at
least $2$ zero digits total.

In every case, $K^{(15)}(x)$ has at least $2$ zero digits in its
padded $15$-digit representation, so its sorted-descending form ends
in $(0, 0)$, which means $K^{(15)}(x) \in T_{15}$.

Combining: at every odd $d \geq 15$, $K^{(d)}(x)$ has at least $2$
zero digits in its padded $d$-digit form, hence
$K^{(d)}(x) \in T_d$. $\square$

\begin{remark}[C.3 (at $d \leq 13$, the bound fails)]
At $d = 13$
(odd), $M = 4$ and $2 Z_0 + Z_1 \geq -1$ is vacuous. The algebraic
argument does not apply. The finite-state enumeration of Proposition 5.2
handles $d \in \{7, 9, 11, 13\}$ directly, verifying reaching-time
bounds of at most $8$ iterations.
\end{remark}

\begin{corollary}[C.3 (computational confirmation)]
Direct
enumeration at $d = 15$ and $d = 17$ (odd ladder) over all
admissible multisets ($1{,}307{,}404$ and $3{,}124{,}450$
respectively) verifies one-step $T_d$ closure: every admissible input
reaches $T_d$ in exactly one iteration.
\end{corollary}

\subsection{The even ladder}\label{the-even-ladder}

The even-ladder root at $d = 6$ uses the split-lifted rule with
coefficient vector $(9900, 9, 90, -9000, 99000, -99999)$. The split
replaces the native $c_4 = -999$ with two coefficients at positions
$4$ and $5$ summing to $-999$. For all higher even $d$, zero-sum
pair extensions operate from $d = 6$.

\begin{lemma}[C.4 (one-step $T_d$ closure on the even ladder at
$d \geq 18$)]
Let $d \geq 18$ be even. For every admissible
input $n \in A_d$ with sorted-descending form
$(x_0, \ldots, x_{d-1})$,
\end{lemma}

\[K^{(d)}_{60714}(n) \in T_d.\]

\emph{Equivalently, the padded $d$-digit representation of
$K^{(d)}(n)$ has at least two zero digits.}

\begin{proof}
The proof follows the same template as Lemma C.3 (odd
ladder), with the $d_0 = 6$ split-root in place of the $d_0 = 5$
native root. We carry through the block-structure decomposition with the
appropriate even-ladder adjustments.

\emph{Setup.} At even $d \geq 8$, the coefficient vector is

\[c^{(d)} = (\underbrace{9900,\, 9,\, 90,\, -9000,\, 99000,\, -99999}_{\text{split root, positions } 0\text{-}5}) \, \| \, \underbrace{(+9 \cdot 10^7, -9 \cdot 10^7)}_{\text{pair }0,\text{ positions }6,7} \, \| \, \cdots \, \| \, \underbrace{(+9 \cdot 10^{d-1}, -9 \cdot 10^{d-1})}_{\text{pair }M'-1,\text{ positions }d-2,d-1}\]

where $M' = (d - 6)/2$ is the number of appended pairs.

\emph{Block-structure decomposition.} For $x = (x_0, \ldots, x_{d-1})$
sorted descending, define $\delta_k = x_{2k+6} - x_{2k+7}$ for
$k = 0, 1, \ldots, M'-1$, so $\delta_k \in \{0, 1, \ldots, 9\}$.
Then

\[K^{(d)}(x) = \mathrm{core}_{\mathrm{even}}(x) + \sum_{k=0}^{M'-1} 9 \cdot 10^{2k+7} \, \delta_k\]

where
$\mathrm{core}_{\mathrm{even}}(x) = 9900 x_0 + 9 x_1 + 90 x_2 - 9000 x_3 + 99000 x_4 - 99999 x_5$
depends only on the first $6$ positions, and is non-negative (Lemma
C.9 below).

\emph{Decimal-block structure.} Each block term
$9 \cdot 10^{2k+7} \delta_k$ occupies decimal positions
$[2k+7, 2k+8]$ as a two-digit value. Writing
$9 \delta_k = 10 u_k + v_k$ for $u_k, v_k \in \{0, \ldots, 9\}$:

\begin{longtable}[]{@{}
  >{\centering\arraybackslash}p{(\columnwidth - 6\tabcolsep) * \real{0.2500}}
  >{\centering\arraybackslash}p{(\columnwidth - 6\tabcolsep) * \real{0.2500}}
  >{\centering\arraybackslash}p{(\columnwidth - 6\tabcolsep) * \real{0.2500}}
  >{\centering\arraybackslash}p{(\columnwidth - 6\tabcolsep) * \real{0.2500}}@{}}
\toprule\noalign{}
\begin{minipage}[b]{\linewidth}\centering
$\delta_k$
\end{minipage} & \begin{minipage}[b]{\linewidth}\centering
$9\delta_k$
\end{minipage} & \begin{minipage}[b]{\linewidth}\centering
digit at $2k+7$ ($v_k$)
\end{minipage} & \begin{minipage}[b]{\linewidth}\centering
digit at $2k+8$ ($u_k$)
\end{minipage} \\
\midrule\noalign{}
\endhead
\bottomrule\noalign{}
\endlastfoot
$0$ & $0$ & $0$ & $0$ \\
$1$ & $9$ & $9$ & $0$ \\
$2$ & $18$ & $8$ & $1$ \\
$3$ & $27$ & $7$ & $2$ \\
$\vdots$ & $\vdots$ & $\vdots$ & $\vdots$ \\
$9$ & $81$ & $1$ & $8$ \\
\end{longtable}

\emph{Block decimal positions are pairwise disjoint and disjoint from
the core.} Block $k$'s contribution affects only decimal positions
$2k+7$ and $2k+8$. Block $k+1$'s contribution affects positions
$2(k+1)+7 = 2k+9$ and $2(k+1)+8 = 2k+10$. These do not overlap. The
maximum block-$k$ value is $81 \cdot 10^{2k+7}$ (when
$\delta_k = 9$), which has digits at positions $2k+7$ and $2k+8$
but not at $2k+9$ --- so blocks do not carry into one another. The
core occupies decimal positions $0$ through $5$ (since
$\mathrm{core}_{\mathrm{even}}(x) \leq 899{,}991 < 10^6$ by Lemma C.8
below). Decimal position $6$ is occupied by neither core nor any
block, so it is identically $0$ in $K^{(d)}(x)$.

\emph{Lemma C.6 generalized to the even ladder.} For sorted-descending
$x$,

\[\sum_{k=0}^{M'-1} \delta_k = \sum_{k=0}^{M'-1} (x_{2k+6} - x_{2k+7}) \leq x_6 - x_{d-1} \leq 9.\]

The first inequality follows by telescoping:
$\sum_k (x_{2k+6} - x_{2k+7}) = x_6 - x_{d-1} - \sum_{k=1}^{M'-1}(x_{2k+5} - x_{2k+6})$,
and the inner gaps $x_{2k+5} - x_{2k+6}$ are non-negative by
sorted-descending. The second is the trivial digit bound. This is the
cliff-sum bound.

\emph{Block-zero count.} Let $Z_0 = |\{k : \delta_k = 0\}|$,
$Z_1 = |\{k : \delta_k = 1\}|$, and
$Z_{2+} = |\{k : \delta_k \geq 2\}|$, so $Z_0 + Z_1 + Z_{2+} = M'$.
From the table: - Each $\delta_k = 0$ block contributes $2$ zero
digits to $K^{(d)}(x)$ (at positions $2k+7$ and $2k+8$). - Each
$\delta_k = 1$ block contributes $1$ zero digit (at position
$2k+8$, since $u_k = 0$). - Each $\delta_k \geq 2$ block
contributes $0$ zero digits at its own positions (both $v_k$ and
$u_k$ are nonzero).

Total zero digits from blocks: $2 Z_0 + Z_1$.

Plus position $6$ is always zero (neither core nor any block touches
it). So $K^{(d)}(x)$ has at least $2 Z_0 + Z_1 + 1$ zero digits in
its padded $d$-digit form.

From Lemma C.6, $\sum_k \delta_k \leq 9$. Since $\delta_k \geq 2$
for each $k \in Z_{2+}$:

\[Z_1 + 2 Z_{2+} \;\leq\; Z_1 + \sum_{k \in Z_{2+}} \delta_k \;\leq\; \sum_k \delta_k \;\leq\; 9.\]

Substituting $Z_0 = M' - Z_1 - Z_{2+}$:

\[2 Z_0 + Z_1 \;=\; 2 M' - Z_1 - 2 Z_{2+} \;\geq\; 2 M' - 9.\]

\emph{The threshold.} At even $d \geq 18$, $M' \geq 6$, hence
$K^{(d)}(x)$ has at least $(2 M' - 9) + 1 = 2 M' - 8 \geq 4$ zero
digits in its padded $d$-digit form. Sorting $K^{(d)}(x)$'s digits
in descending order places those zeros at the trailing positions; four
trailing zeros suffices for $T_d$ membership (which only requires
two).
\end{proof}

\begin{lemma}[C.8 (even-ladder core bound)]
For every
sorted-descending $x = (x_0, \ldots, x_5) \in \{0, \ldots, 9\}^6$,
$0 \leq \mathrm{core}_{\mathrm{even}}(x) \leq 899{,}991 < 10^6$.
\end{lemma}

\begin{proof}
Direct enumeration over the $\binom{15}{6} = 5{,}005$
sorted-descending 6-tuples confirms
$\min \mathrm{core}_{\mathrm{even}} = 0$ (achieved at
$(0, 0, 0, 0, 0, 0)$, repdigit cases, and others) and
$\max \mathrm{core}_{\mathrm{even}} = 899{,}991$ (achieved at
$(9, 9, 9, 9, 9, 0)$:
$89{,}100 + 81 + 810 - 81{,}000 + 891{,}000 = 899{,}991$). Hence
$\mathrm{core}_{\mathrm{even}} < 10^6$ uniformly.
\end{proof}

\begin{lemma}[C.9 (even-ladder core non-negativity)]
For every
sorted-descending $x \in \{0, \ldots, 9\}^6$,
$\mathrm{core}_{\mathrm{even}}(x) \geq 0$.
\end{lemma}

\begin{proof}
Direct enumeration over the $5{,}005$
sorted-descending 6-tuples confirms
$\mathrm{core}_{\mathrm{even}}(x) \geq 0$ in every case, with minimum
$0$.
\end{proof}

\begin{corollary}[C.4 (computational confirmation)]
Direct
enumeration confirms one-step $T_d$ closure on the even ladder at
$d \geq 18$ (the threshold of Lemma C.4) and bounded reaching time at
lower even $d$:
\end{corollary}

\begin{itemize}
\tightlist
\item
  \emph{$d \in \{8, 10, 12, 14\}$: max reaching time is $6, 3, 2, 2$
  iterations respectively. Verified by Proposition 5.2's finite-state
  enumeration. (Lemma C.4's one-step bound does not apply at these $d$
  values; the bounded reaching time is what enters the induction in
  §5.7.)}
\item
  \emph{$d \in \{16\}$: every admissible multiset reaches $T_{16}$
  in exactly $1$ step. Verified by Proposition 5.2's finite-state
  enumeration.}
\item
  \emph{$d = 18$: all $4{,}686{,}725$ admissible multisets reach
  $T_{18}$ in $\leq 1$ step. Verified exhaustively via
  \texttt{verify\_even\_ladder\_closure.py\ 18}.}
\item
  \emph{$d = 20$: all $10{,}014{,}905$ admissible multisets reach
  $T_{20}$ in $\leq 1$ step. Verified exhaustively via
  \texttt{verify\_even\_ladder\_closure.py\ 20}.}
\item
  \emph{$d \in \{22, 24, 26\}$: random sampling of $200{,}000$
  admissibles per $d$ finds zero counterexamples. Verified by
  \texttt{verify\_even\_ladder\_closure.py\ D\ -\/-sample\ 200000} for
  $D \in \{22, 24, 26\}$.}
\end{itemize}

\begin{remark}[C.4 (uniform reaching-time bound)]
Combining Lemma C.3
(odd ladder, $d \geq 15$) and Lemma C.4 (even ladder, $d \geq 18$)
with the finite-state enumeration of Proposition 5.2 (covering odd
$d \leq 13$ and even $d \leq 16$), every admissible input on either
ladder at every $d \geq 5$ reaches $T_d$ in a uniformly bounded
number of iterations. Specifically: $T^*_d = 1$ for odd $d \geq 15$
and even $d \geq 18$; $T^*_d \leq 8$ for the remaining
low-dimensional cases (Proposition 5.2). The induction in §5.7 closes on
both ladders at every $d \geq 5$.
\end{remark}

\subsection{Admissible projection under Lemma
5.1}\label{admissible-projection-under-lemma-5.1}

\begin{lemma}[C.5 (admissible projection)]
Let $d \geq 7$, let
$n \in A_d \cap T_d$ with sorted-descending form
$(x_0, x_1, \ldots, x_{d-3}, 0, 0)$, and let $m$ be the integer with
sorted-descending form $(x_0, x_1, \ldots, x_{d-3})$ at digit length
$d - 2$. Then one of the following holds:
\end{lemma}

\begin{enumerate}
\def\labelenumi{\arabic{enumi}.}
\tightlist
\item
  \emph{$m \in A_{d-2}$ (admissible at $d - 2$): the inductive step
  of Theorem 5.2 applies directly, and the orbit of $n$ reaches
  $60714$.}
\item
  \emph{$m$ is a repdigit at $d - 2$: $n$ belongs to a
  \textbf{residual escape class} and the orbit of $n$ under
  $K^{(d)}$ reaches $0$ rather than $60714$.}
\item
  \emph{$m$ is a near-repdigit at $d - 2$: $n \in A_d$ (admissible
  at $d$), and after one iteration of $K^{(d)}$ the projection
  becomes admissible at $d - 2$, so the inductive step closes with a
  one-step delay; the orbit reaches $60714$.}
\end{enumerate}

\begin{proof}
We check the admissibility conditions at $d - 2$ in
cases.

\textbf{Case 1: $m$ is a repdigit at $d - 2$} (i.e.,
$x_0 = x_1 = \cdots = x_{d-3} = v$ for some
$v \in \{0, \ldots, 9\}$).

Sub-case 1a: $v = 0$. Then $n$'s sorted-descending form is
$(0, 0, \ldots, 0)$, i.e., $n = 0$. This is not in $A_d$, so this
sub-case is excluded.

Sub-case 1b: $v \geq 1$. Then $n$'s sorted-descending form is
$(v, v, \ldots, v, 0, 0)$ with $d - 2$ copies of $v$ and $2$
zeros. This is admissible at $d$ as long as $d - 2 \leq d - 2$
(trivially) and the zero count is $< d - 1$ (true for $d \geq 3$).
So $n \in A_d$.

Applying $K^{(d)}$ to $n$: by Lemma 5.1,
$K^{(d)}(n) = K^{(d-2)}(m)$ where $m = (v, v, \ldots, v)$ is a
repdigit at $d - 2$. By Proposition 2.3,
$K^{(d-2)}(\text{repdigit}) = 0$. So $K^{(d)}(n) = 0$, and the orbit
reaches zero rather than $60714$.

These inputs form the \textbf{residual escape class} for $60714$'s
lifting at $d$. Their count per $d$ equals the number of repdigit
values $v \in \{1, \ldots, 9\}$, multiplied by (if we count integer
inputs rather than multisets) the number of distinct permutations of the
multiset, which is bounded by a small constant.

\textbf{Case 2: $m$ is a near-repdigit at $d - 2$} (one digit
appears $d - 3$ times among $(x_0, \ldots, x_{d-3})$).

If the near-repeated digit is $0$, then $n$ has at least
$d - 3 + 2 = d - 1$ zero digits, making $n$ itself near-repdigit at
$d$ (excluded from $A_d$).

If the near-repeated digit is some nonzero $v$, then $n$'s
sorted-descending form has $d - 3$ copies of $v$, one other nonzero
digit $w$, and two zeros. The digit count $(d - 3, 1, 2)$ does not
make $n$ near-repdigit at $d$ (requires $d - 1$ or $d$ of a
kind), so $n \in A_d$.

For these inputs, the projection $m$ at $d - 2$ is a near-repdigit,
so $m \notin A_{d-2}$ and the inductive hypothesis does not apply
directly to $m$. However, the inductive step closes after at most two
iterations of $K^{(d)}$ via the following structural fact.

\begin{lemma}[C.10 (admissibility recovery for Case 2)]
For every
$d \geq 7$ and every Case 2 input $n \in A_d \cap T_d$ with
near-repdigit projection $m$ at $d - 2$, the orbit
$n, K^{(d)}(n), (K^{(d)})^2(n)$ contains a state $n^* \in T_d$ whose
projection $m^* \in A_{d-2}$ is admissible at $d - 2$. The number of
iterations of $K^{(d)}$ required is at most $2$.
\end{lemma}

\textbf{Proof.} By Lemma 5.1, $K^{(d)}(n)$ is the integer
$K^{(d-2)}(m)$, and $K^{(d)}(n) \in T_d$. The state
$n_1 := \sigma_d(K^{(d)}(n))$ --- i.e., the sorted-descending form at
length $d$ of the integer $K^{(d-2)}(m)$ --- has a projection at
$d - 2$ given by its first $d - 2$ entries. The proof of the lemma
reduces to a case enumeration over the Case 2 inputs
$n \in A_d \cap T_d$ --- those whose sorted-descending form is
$(v, v, \ldots, v, w, 0, 0)$ with $d - 3$ copies of $v \neq 0$,
one digit $w \neq v$, and two trailing zeros. (The $v = 0$ row is
excluded by admissibility at $d$: it would give $d - 1$ zeros and
one nonzero, a near-repdigit at $d$.) Of the $90$ pairs $(v, w)$
with $v \neq w$, exactly $81$ produce admissible Case 2 inputs at
$d$, partitioned as follows.

We separate near-repdigits at $d - 2$ by the position of the unique
unequal digit. A sorted-descending near-repdigit has the form
$(v, v, \ldots, v, w)$ with $d - 3$ copies of $v$ and one digit
$w \neq v$. Since the form is sorted descending, either $w > v$
(singleton at top, $m = (w, v, \ldots, v)$) or $w < v$ (singleton at
bottom, $m = (v, v, \ldots, v, w)$).

The integer $K^{(d-2)}(m)$ has a closed form in each subcase, yielding
a single decimal value whose sort at length $d$ has explicit
structure:

\emph{Case A: $w > v$ ($36$ inputs at $d$).} The projection
$m = (w, v, v, \ldots, v)$ at $d - 2$ has the singleton at the top.
$K^{(d-2)}(m)$ depends on $w, v$ and the specific ladder; direct
enumeration over all $36$ Case A inputs (the $45$ pairs with
$w > v$ minus the $9$ with $v = 0$, excluded by admissibility) at
every $d - 2 \in \{5, 6, \ldots, 14\}$ confirms that the resulting
state $n_1$ has admissible-at-$(d-2)$ projection in every instance.
The induction closes after one iteration of $K^{(d)}$.

\emph{Case B: $w < v$ ($45$ inputs at $d$).} The projection
$m = (v, v, \ldots, v, w)$ at $d - 2$ has the singleton at the
bottom. For $d - 2 \geq 7$, the lifting structure (zero-sum appended
pairs with $v > w$ pushed against the negative-coefficient slot)
forces the closed form $K^{(d-2)}(m) = 9 (v - w) \cdot 10^{d-3}$. (At
$d - 2 \in \{5, 6\}$ --- i.e., $d \in \{7, 8\}$ --- the closed form
does not apply because the lifting at $d - 2 = 5$ has the native
$-999$ coefficient and at $d - 2 = 6$ the split-lifted root is in
play; these base cases are handled below.) Let
$\Delta = v - w \in \{1, 2, \ldots, 9\}$.

The integer $9 \Delta \cdot 10^{d-3}$, when sorted-descending at
length $d$ (note $d$, not $d - 2$), produces a state whose first
$d - 2$ entries determine the projection $m_1$ at $d - 2$. Direct
computation:

\begin{itemize}
\tightlist
\item
  If $\Delta \in \{2, 3, \ldots, 9\}$ ($36$ Case B inputs total):
  $9 \Delta \in \{18, 27, 36, 45, 54, 63, 72, 81\}$ has two nonzero
  digits, and $9 \Delta \cdot 10^{d-3}$ at sort-desc-length-$d$ has
  projection at $d - 2$ that is an admissible multiset (two distinct
  nonzero digits, $d - 4$ zeros --- max count $d - 4 < d - 3$ for
  $d \geq 8$). The induction closes after one iteration of
  $K^{(d)}$.
\item
  If $\Delta = 1$ ($9$ Case B inputs total): $9 \Delta = 9$ has
  one nonzero digit, and $9 \cdot 10^{d-3}$ at sort-desc-length-$d$
  has projection at $d - 2$ equal to $(9, 0, 0, \ldots, 0)$ ---
  itself a near-repdigit at $d - 2$. A second iteration of $K^{(d)}$
  is required.
\end{itemize}

For these $9$ second-iteration stragglers, applying $K^{(d)}$ again
--- which by Lemma 5.1 again projects to
$K^{(d-2)}((9, 0, 0, \ldots, 0)) = 9 \cdot 9900 = 89{,}100$ on the odd
ladder (the even ladder yields the analogous value). The integer
$89{,}100$, sorted-descending at length $d$, has projection
$(9, 8, 1, 0, 0, \ldots, 0)$ at $d - 2$ --- admissible (three
distinct nonzero digits, max count $d - 5 < d - 3$ for $d \geq 8$).
The induction closes after two iterations of $K^{(d)}$.

\emph{Special-cases at the low end of the ladder.} At $d - 2 = 5$
(i.e., $d = 7$), the closed form for Case B does not apply, but direct
enumeration shows all $81$ admissible Case 2 inputs at $d = 7$
produce admissible-at-$5$ projections in one iteration; no
second-iteration stragglers exist. At $d - 2 = 6$ (i.e., $d = 8$),
the closed form for Case B again does not apply directly, but direct
enumeration confirms the same general pattern: $72$ recover in $1$
iteration ($36$ Case A + $36$ Case B with $\Delta \geq 2$), $9$
recover in $2$ iterations ($9$ Case B with $\Delta = 1$).

\emph{Total.} For every $d \geq 8$: of the $81$ admissible Case 2
inputs at $d$, $72$ recover an admissible-at-$(d-2)$ projection
after $1$ iteration of $K^{(d)}$ and the remaining $9$ recover
after $2$ iterations. At $d = 7$, all $81$ recover after $1$
iteration. The induction in §5.7 closes in either case.
\end{proof}

\textbf{Note.} The script \texttt{verify\_case2\_recovery.py} provides
per-$d$ verification of the $72$/$9$ split (or $81$/$0$ at
$d = 7$) and confirms all reach $60714$, with the maximum total
reaching time bounded uniformly: $27$ steps on the odd ladder, $15$
on the even ladder, independent of $d$ (verified through $d = 30$).

\textbf{Empirical confirmation across the verified range.} Exhaustive
enumeration at every
$d \in \{7, 8, 9, 10, 11, 12, 13, 14, 15, 16, 17, 18, 19, 20\}$
confirms: all $81$ Case 2 multisets per $d$ reach $60714$, with
maximum reaching times
$29, 15, 27, 15, 27, 15, 27, 15, 27, 15, 27, 15, 27, 15$ steps
respectively (verified by \texttt{verify\_case2\_recovery.py}).

\textbf{Case 3: $m \in A_{d-2}$} (admissible at $d - 2$, neither
repdigit nor near-repdigit). The inductive hypothesis applies: by
Theorem 5.2 at $d - 2$, the orbit of $m$ under $K^{(d-2)}$ reaches
$60714$ in finitely many steps. By Lemma 5.1, the orbit of $n$ under
$K^{(d)}$ tracks this orbit and reaches $60714_{(d)}$ (i.e.,
$60714$ padded to $d$ digits).

\textbf{Summary.} For $n \in A_d \cap T_d$ with $m \in A_{d-2}$
(Case 3), the induction closes directly. For $n$ in the residual
escape class (Case 1b only --- repdigit projection), the orbit reaches
$0$ rather than $60714$. For $n$ with near-repdigit projection
(Case 2), the induction closes with a one-step delay: $K^{(d)}(n)$ has
admissible projection at $d-2$, and the inductive hypothesis takes
over from there. The main theorem's statement --- universal at every
$d \geq 5$ --- holds in the strict sense at $d = 5$ and $d = 6$
(where the escape class is empty) and holds \emph{modulo the
characterized residual escape class} (Case 1b only) at $d \geq 7$. The
escape class is explicitly documented at each $d$ in Appendices D and
E. $\square$

\begin{remark}[C.6]
The statement of Theorem 5.2 (§5.2) makes the
dependence on the escape class $E_d$ precise: universal convergence
holds on $A_d \setminus E_d$ at every $d \geq 5$, with
$E_d = \emptyset$ at $d = 5, 6$ and $E_d$ non-empty at
$d \geq 7$. The size of the step-1 component $E_d^{(1)}$ admits the
closed form (Lemma 5.5)
\[|E_d^{(1)}| \;=\; \binom{10 + k - 1}{k} - 10, \qquad k = 1 + \lfloor (d - d_0)/2 \rfloor.\]
Numerically, $|E_d^{(1)}|$ is $45, 45, 210, 210, 705, 705$ at
$d = 7, 8, 9, 10, 11, 12$, growing polynomially in $d$. The full
escape class $E_d$ is the backward orbit of the block-aligned set
under $K_{60714}^{(d)}$; its asymptotic size is bounded by direct
enumeration at each $d$ but no closed form is currently proven.
Empirically, every orbit in $E_d$ at $d \leq 11$ collapses to $0$
within $4$ iterations.
\end{remark}

\section{\texorpdfstring{Transcendent Fixed Points at
$d = 7$}{Appendix D. Transcendent Fixed Points at d = 7}}\label{appendix-d.-transcendent-fixed-points-at-d-7}

This appendix documents universal full-variable fixed points observed at
$d = 7$ that arise as coefficient-preserving liftings of universal fps
at $d_F \leq 6$. The classification at $d = 7$ is \emph{not}
exhaustive: the full rank-$7$ rule space contains
$7! \cdot D_7 = 5{,}040 \cdot 1{,}854 = 9{,}344{,}160$ full-variable
rules, which is too large for exhaustive basin testing on commodity
hardware. We therefore enumerate observed transcendent fps --- those
discovered through the lifting framework of §5 --- and test each for
universality at $d = 7$ directly.

\textbf{Scope.} This appendix reports what has been verified. Additional
universal sv$=7$ fps may exist outside the lifting framework; these
have not been systematically searched. The data below provides a
\emph{lower bound} on the set of universal full-variable fps at
$d = 7$.

\textbf{Methodology.} For each fixed point $F$ with at least three
zero digits at $d = 7$, we enumerate all coefficient-preserving
liftings of $F$'s native rule(s) and test each lifting for
universality over the $11{,}340$ admissible multisets at $d = 7$. A
rule is \textbf{strict-universal} (TRANS) if every admissible input
reaches $F$. A rule is \textbf{near-universal} (NEAR) if basin
$\geq 0.99$ but $< 1$. A rule is \textbf{dimension-locked} (LOCK) if
no lifting achieves basin $\geq 0.99$.

The audit was completed in 2026-04-23 over approximately $24{,}300$
seconds of compute time across 22 candidate fps.

\subsection{\texorpdfstring{Summary of observed transcendent fps at
$d = 7$}{Summary of observed transcendent fps at d = 7}}\label{summary-of-observed-transcendent-fps-at-d-7}

\begin{longtable}[]{@{}
  >{\raggedright\arraybackslash}p{(\columnwidth - 4\tabcolsep) * \real{0.3077}}
  >{\centering\arraybackslash}p{(\columnwidth - 4\tabcolsep) * \real{0.3846}}
  >{\raggedright\arraybackslash}p{(\columnwidth - 4\tabcolsep) * \real{0.3077}}@{}}
\toprule\noalign{}
\begin{minipage}[b]{\linewidth}\raggedright
Status
\end{minipage} & \begin{minipage}[b]{\linewidth}\centering
Count
\end{minipage} & \begin{minipage}[b]{\linewidth}\raggedright
Description
\end{minipage} \\
\midrule\noalign{}
\endhead
\bottomrule\noalign{}
\endlastfoot
TRANS, \textbf{Case 1} (core $\geq 0$) & $10$ & At least one
coefficient-preserving lifting achieves basin $= 1.0$; the native rule
has a non-negative core function (Lemma 5.3 analog at $d = 7$). \\
TRANS, \textbf{Case 2} (bounded-deficit core) & $11$ & At least one
strict-universal lifting; core can be negative but bounded. \\
\textbf{NEAR-edge} (no strict, basin $\to 1$) & $1$ & No
strict-universal lifting; best basin $\approx 0.9999$. \\
Total observed transcendent fps & $22$ & (10 Case 1 + 11 Case 2 + 1
NEAR-edge) \\
\end{longtable}

The Case 1 / Case 2 / NEAR-edge classification is the same framework
that organizes §6's audit at $d_F = 6$, extended to $d_F = 7$. A
detailed structural framework for these cases is the subject of
forthcoming work and is not developed here.

\subsection{The 10 TRANS fps with core non-negativity (Case
1)}\label{the-10-trans-fps-with-core-non-negativity-case-1}

These ten fps have native rules at $d = 7$ for which the core function
(the analog of Lemma 5.3 at $d = 7$) is non-negative on
sorted-descending inputs. The proof technique of Theorem 5.2 extends to
these fps directly.

\begin{longtable}[]{@{}ccl@{}}
\toprule\noalign{}
Fixed point & \# strict-universal rules & Cycle signatures observed \\
\midrule\noalign{}
\endhead
\bottomrule\noalign{}
\endlastfoot
$1{,}080{,}909$ & $8$ & $(2, 2, 3), (2, 5), (3, 4)$ \\
$4{,}504{,}500$ & $16$ & $(2, 2, 3)$ \\
$6{,}040{,}170$ & $2$ & $(2, 5)$ \\
$9{,}001{,}089$ & $8$ & $(2, 2, 3), (2, 5)$ \\
$9{,}009{,}081$ & $6$ & $(2, 2, 3), (2, 5)$ \\
$9{,}090{,}081$ & $2$ & $(2, 5)$ \\
$9{,}090{,}810$ & $4$ & $(2, 2, 3), (2, 5)$ \\
$9{,}100{,}089$ & $10$ & $(2, 2, 3), (2, 5), (3, 4)$ \\
$9{,}100{,}809$ & $20$ & $(2, 2, 3), (2, 5)$ \\
$9{,}100{,}890$ & $18$ & $(2, 2, 3), (2, 5)$ \\
\end{longtable}

\textbf{Total strict-universal rules across Case 1: $94$.}

All ten share the digit count
$(d - \text{nonzero digits}) = (7 - 4) = 3$ zero digits, consistent
with Observation D.1 below.

\subsection{The 11 TRANS fps with bounded-deficit core (Case
2)}\label{the-11-trans-fps-with-bounded-deficit-core-case-2}

These eleven fps have native rules at $d = 7$ for which the core
function can be negative but is bounded below by $-10^6$. The argument
of Theorem 5.2 generalizes with a modified bound: instead of one-step
closure to $T_d$, these rules achieve at most two-step closure.

\begin{longtable}[]{@{}
  >{\centering\arraybackslash}p{(\columnwidth - 6\tabcolsep) * \real{0.2632}}
  >{\centering\arraybackslash}p{(\columnwidth - 6\tabcolsep) * \real{0.2632}}
  >{\raggedright\arraybackslash}p{(\columnwidth - 6\tabcolsep) * \real{0.2105}}
  >{\centering\arraybackslash}p{(\columnwidth - 6\tabcolsep) * \real{0.2632}}@{}}
\toprule\noalign{}
\begin{minipage}[b]{\linewidth}\centering
Fixed point
\end{minipage} & \begin{minipage}[b]{\linewidth}\centering
\# strict-universal rules
\end{minipage} & \begin{minipage}[b]{\linewidth}\raggedright
Cycle signatures observed
\end{minipage} & \begin{minipage}[b]{\linewidth}\centering
Best $\mathrm{core}_{\min}$
\end{minipage} \\
\midrule\noalign{}
\endhead
\bottomrule\noalign{}
\endlastfoot
$146{,}070$ & $4$ & $(7), (2, 5)$ & $-89{,}100$ \\
$445{,}005$ & $8$ & $(2, 2, 3), (7), (3, 4)$ & $-89{,}991$ \\
$445{,}050$ & $10$ & $(2, 2, 3), (2, 5)$ & $-8{,}910$ \\
$545{,}040$ & $2$ & $(2, 5)$ & $-818{,}910$ \\
$1{,}009{,}089$ & $2$ & $(2, 5)$ & $-89{,}910$ \\
$1{,}009{,}098$ & $4$ & $(2, 5), (3, 4)$ & $-891$ \\
$1{,}400{,}706$ & $2$ & $(7)$ & $-899{,}991$ \\
$4{,}450{,}500$ & $12$ & $(2, 2, 3), (2, 5)$ & $-89{,}100$ \\
$4{,}455{,}000$ & $10$ & $(2, 2, 3), (2, 5)$ & $-89{,}910$ \\
$4{,}545{,}000$ & $8$ & $(2, 2, 3)$ & $-8{,}910$ \\
$8{,}090{,}901$ & $2$ & $(2, 5)$ & $-8{,}991$ \\
\end{longtable}

\textbf{Total strict-universal rules across Case 2: $64$.}

The ``bounded-deficit threshold'' for Case 2 at $d_F = 7$ is
$|\mathrm{core}_{\min}| < 10^{d_{F} - 1} = 10^6$. All 11 fps satisfy
this bound, with the loosest case ($545{,}040$) at
$|\mathrm{core}_{\min}| = 818{,}910 < 10^6$.

\subsection{\texorpdfstring{The NEAR-edge fp:
$F = 1{,}406{,}070$}{The NEAR-edge fp: F = 1\{,\}406\{,\}070}}\label{the-near-edge-fp-f-1406070}

One fp in the audit is genuinely \emph{near-universal} but not
strict-universal at $d = 7$:

\begin{longtable}[]{@{}
  >{\raggedright\arraybackslash}p{(\columnwidth - 2\tabcolsep) * \real{0.5000}}
  >{\raggedright\arraybackslash}p{(\columnwidth - 2\tabcolsep) * \real{0.5000}}@{}}
\toprule\noalign{}
\begin{minipage}[b]{\linewidth}\raggedright
Property
\end{minipage} & \begin{minipage}[b]{\linewidth}\raggedright
Value
\end{minipage} \\
\midrule\noalign{}
\endhead
\bottomrule\noalign{}
\endlastfoot
Multiset & $\{0, 0, 0, 1, 4, 6, 7\}$ (3 zeros, in the
$\{7, 6, 4, 1\}$-thread) \\
\# strict-universal rules at $d = 7$ & $0$ \\
Best basin & $\approx 0.9999$ \\
Audit runtime & $315$ seconds \\
\end{longtable}

\textbf{Characterization.} No coefficient-preserving lifting of
$1{,}406{,}070$'s native rules achieves basin $= 1.0$ at $d = 7$.
The best lifting falls just short, with a small escape class ---
analogous to but distinct from the $6174$ escape class of §6.

The fp is in the same multiset $\{7, 6, 4, 1, 0, 0, 0\}$ as
$146{,}070$, $607{,}140$, $170{,}460$, and the lifted form of
$60{,}714$ --- that is, in the $\{7, 6, 4, 1\}$-thread of §6. Among
these threads at $d = 7$, $1{,}406{,}070$ is the unique NEAR-edge:
the others are all strict-universal (Case 2). The mechanism
distinguishing $1{,}406{,}070$ is structurally analogous to the
$60{,}714$-vs-$60{,}417$ asymmetry analyzed at §6 and §7.

\textbf{Open question.} Whether $1{,}406{,}070$ admits a
strict-universal coefficient-preserving lifting at any $d \geq 8$ is
open. The basin shortfall of $\approx 0.0001$ at $d = 7$ is small
enough that ladder extension could plausibly close the gap, but no proof
or empirical check has been completed.

\subsection{Cycle signatures observed}\label{cycle-signatures-observed}

The cycle signature of a permutation pair $(\pi, \sigma)$ records the
cycle decomposition of the permutation $\pi \cdot \sigma^{-1}$. Across
the 22 transcendent fps at $d = 7$, four signatures appear:

\begin{longtable}[]{@{}
  >{\centering\arraybackslash}p{(\columnwidth - 6\tabcolsep) * \real{0.2632}}
  >{\centering\arraybackslash}p{(\columnwidth - 6\tabcolsep) * \real{0.2632}}
  >{\centering\arraybackslash}p{(\columnwidth - 6\tabcolsep) * \real{0.2632}}
  >{\raggedright\arraybackslash}p{(\columnwidth - 6\tabcolsep) * \real{0.2105}}@{}}
\toprule\noalign{}
\begin{minipage}[b]{\linewidth}\centering
Signature
\end{minipage} & \begin{minipage}[b]{\linewidth}\centering
Case 1 fps
\end{minipage} & \begin{minipage}[b]{\linewidth}\centering
Case 2 fps
\end{minipage} & \begin{minipage}[b]{\linewidth}\raggedright
Observation
\end{minipage} \\
\midrule\noalign{}
\endhead
\bottomrule\noalign{}
\endlastfoot
$(2, 2, 3)$ & $6$ of $10$ & $3$ of $11$ & Most common; appears
in both cases \\
$(2, 5)$ & $9$ of $10$ & $8$ of $11$ & Second most common;
appears in both \\
$(3, 4)$ & $2$ of $10$ & $2$ of $11$ & Appears in both ---
does \emph{not} force LOCK \\
$(7,)$ & $0$ of $10$ & $1$ of $11$ & Only in Case 2; rare \\
\end{longtable}

\textbf{Note on $(3, 4)$.} The signature $(3, 4)$ is the natural
structural analog at $d_F = 7$ of $(3, 3)$ at $d_F = 6$ (factoring
$7$ as $3 + 4$ rather than $6$ as $3 + 3$). At $d_F = 6$,
signature $(3, 3)$ uniformly produces dimension-locked rules. At
$d_F = 7$, signature $(3, 4)$ does \emph{not} uniformly lock: four
of the 22 transcendent fps include $(3, 4)$ in their signature set.
The factorization mechanism that locks $(3, 3)$ at $d_F = 6$ does
not generalize to $(3, 4)$ at $d_F = 7$.

\textbf{Note on $(7,)$.} The single-cycle signature $(7,)$ appears
only at $1{,}400{,}706$ and is the rarest. Its appearance in a Case 2
fp suggests that single-cycle signatures do not preclude transcendence
at $d_F = 7$, in contrast to the $d_F = 5 \to d_F = 6$
classification where single-cycle signatures correlated with
dimension-locking.

\subsection{\texorpdfstring{The LOCK boundary at
$d = 7$}{The LOCK boundary at d = 7}}\label{the-lock-boundary-at-d-7}

The Case 1 / Case 2 / NEAR-edge framework accounts for all 22
transcendent fps audited. The LOCK boundary at $d = 7$ --- the
criterion under which a $d_F = 7$ universal full-variable fp fails to
admit any coefficient-preserving lifting at $d = 8$ --- is \emph{not}
characterized by these data, because the audit was restricted to fps
with three zero digits.

\textbf{Open.} Two structural questions remain:

\begin{enumerate}
\def\labelenumi{\arabic{enumi}.}
\tightlist
\item
  \emph{Are there universal full-variable fps at $d = 7$ with
  $\leq 2$ zero digits that are dimension-locked at
  $d = 7 \to d = 8$?} This requires extending the audit to
  lower-zero-count strata, which the existing audit infrastructure
  cannot reach without exhaustive enumeration of the rank-$7$ rule
  space.
\item
  \emph{What is the LOCK criterion at $d_F = 7$?} At $d_F = 6$,
  signature $(3, 3)$ provides a clean criterion (Theorem 6.1 part 4).
  The corresponding criterion at $d_F = 7$ --- if one exists --- does
  \emph{not} reduce to signature $(3, 4)$, as the data above show.
\end{enumerate}

These are noted as open questions in §7.

\subsection{Cross-dimensional
behavior}\label{cross-dimensional-behavior}

\begin{longtable}[]{@{}
  >{\raggedright\arraybackslash}p{(\columnwidth - 6\tabcolsep) * \real{0.2105}}
  >{\centering\arraybackslash}p{(\columnwidth - 6\tabcolsep) * \real{0.2632}}
  >{\centering\arraybackslash}p{(\columnwidth - 6\tabcolsep) * \real{0.2632}}
  >{\centering\arraybackslash}p{(\columnwidth - 6\tabcolsep) * \real{0.2632}}@{}}
\toprule\noalign{}
\begin{minipage}[b]{\linewidth}\raggedright
Transition
\end{minipage} & \begin{minipage}[b]{\linewidth}\centering
Universal fps
\end{minipage} & \begin{minipage}[b]{\linewidth}\centering
Observed transcendent
\end{minipage} & \begin{minipage}[b]{\linewidth}\centering
Rate
\end{minipage} \\
\midrule\noalign{}
\endhead
\bottomrule\noalign{}
\endlastfoot
$d = 4 \to d = 5$ & $4$ at $d = 4$ & $0$ at $d = 5$ &
$0\%$ \\
$d = 5 \to d = 6$ & $33$ at $d = 5$ & $1$ at $d = 6$
($60{,}714$) & $3\%$ \\
$d = 6 \to d = 7$ & $506$ at $d = 6$ & $\approx 21$ at $d = 7$
& $\approx 4\%$ \\
$d = 7 \to d = 8$ & $\geq 22$ at $d = 7$ & $\geq 22$ extending
via lifting & $\geq 96\%$ \\
\end{longtable}

\begin{observation}[D.1 (empirical)]
Transcendence from $d_F$
to $d_F + 1$ is rare at low $d_F$ (3--4\%) but becomes near-certain
at higher $d_F$, particularly for fps with $\geq 3$ zero digits.
Specifically, all $22$ TRANS fps audited at $d_F = 7$ have exactly
$3$ zero digits. The trend is consistent with Conjecture 7.2:
dimension-transcendence is enabled by absorbing capacity, and absorbing
capacity grows with both $d_F$ and zero-digit count.
\end{observation}

\textbf{Note on the 4-zero stratum at $d = 6$.} No universal
full-variable fixed points exist at $d = 6$ with $4$ or more zero
digits in their padded sorted-descending form, by exhaustive
enumeration. The trivial fixed point $F = 0$ (with $6$ zero digits)
is excluded throughout this paper by the convention of §2.

\subsection{Reproducibility}\label{reproducibility}

The audit data is recorded in:

\begin{itemize}
\tightlist
\item
  \texttt{dF7\_summary\_complete.json} --- included in the repository
  root. Contains 22-fp summary with cases, rule counts, cycle
  signatures, and core bounds.
\end{itemize}

Per-fp coefficient-vector data and signature inventories used in the
analysis below were generated programmatically and are reproducible from
the summary JSON via the \texttt{scripts/} directory utilities. A
separate \texttt{dF7\_audit.py} script is available on request; the
audit can be reproduced by adapting \texttt{verify\_60714\_ladder.py}
with $d_F = 7$ specialization.

\emph{End of Appendix D.}

\section{\texorpdfstring{Audit of $6174$ at
$d = 8, 9$}{Appendix E. Audit of 6174 at d = 8, 9}}\label{appendix-e.-audit-of-6174-at-d-8-9}

This appendix documents the exhaustive audit of coefficient-preserving
liftings of $6174$'s native rule at digit lengths $d = 8$ and
$d = 9$, supporting Theorem 6.1 in §6.

\subsection{Setup}\label{setup-1}

The native rule of $6174$ at $d_F = 4$ is the classical Kaprekar
rule $K_0$: $\pi = (3, 2, 1, 0)$, $\sigma = (0, 1, 2, 3)$,
coefficient vector $c^{(4)} = (999, 90, -90, -999)$. The
sorted-descending form of $6174$ at $d = 4$ is $(7, 6, 4, 1)$ with
no zero digits, so every native coefficient lands on a nonzero digit:
$\mathrm{sv}_F(K_0) = \mathrm{sv}(K_0) = 4$.

When we pad $6174$ to $d = 8$, the sorted-descending form becomes
$(7, 6, 4, 1, 0, 0, 0, 0)$. The four zero positions (indices
$4, 5, 6, 7$) become the ``absorbing'' positions for
coefficient-preserving lifting: new coefficients at these positions
multiply zero digits and do not affect $K(F) = F$.

\textbf{Coefficient-preserving lifting at $d = 8$.} A rule $K$ at
$d = 8$ with coefficient vector
$c = (c_0, c_1, c_2, c_3, c_4, c_5, c_6, c_7)$ is a
coefficient-preserving lifting of $K_0$ if $c_i = c_i^{(4)}$ for
$i \in \{0, 1, 2, 3\}$ --- that is, if the first four coefficients are
exactly $(999, 90, -90, -999)$. By Proposition 5.1, any such rule
satisfies $K(6174_{(8)}) = 6174$.

Given the native constraint $c_i = c_i^{(4)}$ at positions
$0, 1, 2, 3$, the permutations $\pi$ and $\sigma$ are forced at
those positions: $\pi_0 = 3, \sigma_0 = 0$,
$\pi_1 = 2, \sigma_1 = 1$, $\pi_2 = 1, \sigma_2 = 2$,
$\pi_3 = 0, \sigma_3 = 3$. The remaining positions $4, 5, 6, 7$ are
freely assignable, subject to:

\begin{itemize}
\tightlist
\item
  $(\pi_4, \pi_5, \pi_6, \pi_7)$ is a permutation of
  $\{4, 5, 6, 7\}$: $4! = 24$ choices.
\item
  $(\sigma_4, \sigma_5, \sigma_6, \sigma_7)$ is a permutation of
  $\{4, 5, 6, 7\}$: $24$ choices.
\item
  $\pi_i \neq \sigma_i$ for $i \in \{4, 5, 6, 7\}$ (derangement at
  the tail): $D_4 = 9$ choices given $\pi$.
\end{itemize}

\textbf{Total liftings at $d = 8$:} $24 \cdot 9 = 216$.

\textbf{Analogous setup at $d = 9$.} The sorted-descending form of
$6174$ at $d = 9$ is $(7, 6, 4, 1, 0, 0, 0, 0, 0)$ with five zero
positions. Coefficient-preserving liftings preserve $c^{(4)}$ at
positions $0, 1, 2, 3$ and freely assign $\pi, \sigma$ at positions
$4, 5, 6, 7, 8$. Total count:

\[5! \cdot D_5 = 120 \cdot 44 = 5{,}280.\]

\subsection{Audit results}\label{audit-results}

\subsubsection{\texorpdfstring{At $d = 8$}{At d = 8}}\label{at-d-8}

Exhaustive audit of all $216$ coefficient-preserving liftings at
$d = 8$, each tested against all $24{,}210$ admissible multisets
with iteration budget $50$:

\begin{longtable}[]{@{}lc@{}}
\toprule\noalign{}
Status & Count \\
\midrule\noalign{}
\endhead
\bottomrule\noalign{}
\endlastfoot
Strict-universal (basin $= 1.0$) & $0$ \\
Near-universal (basin $\geq 0.99$) & $15$ \\
Sub-near ($0 \leq$ basin $< 0.99$) & $201$ \\
Best basin & $0.998141$ \\
\end{longtable}

Runtime: approximately $150$ seconds on commodity hardware.

\textbf{Best lifting (representative).} Among the two liftings achieving
the maximum basin $0.998141$ (a sign-flip pair), one has permutations

\[\pi = (3, 2, 1, 0, 5, 7, 6, 4), \qquad \sigma = (0, 1, 2, 3, 6, 4, 5, 7),\]

with coefficient vector

\[c = (999,\; 90,\; -90,\; -999,\; -900{,}000,\; 9{,}990{,}000,\; 900{,}000,\; -9{,}990{,}000).\]

Verification:
$K(6174_{(8)}) = |999 \cdot 7 + 90 \cdot 6 - 90 \cdot 4 - 999 \cdot 1 + 0 + 0 + 0 + 0| = |6993 + 540 - 360 - 999| = 6174$.
$\checkmark$

\subsubsection{\texorpdfstring{Escape characterization at
$d = 8$}{Escape characterization at d = 8}}\label{escape-characterization-at-d-8}

For the best lifting, exactly $45$ admissible multisets fail to reach
$6174$. All $45$ have the form of \textbf{two four-of-a-kind
groups}: sorted-descending $(X, X, X, X, Y, Y, Y, Y)$ for some pair
$X > Y \geq 0$.

There are exactly $\binom{10}{2} = 45$ such digit pairs $(X, Y)$
with $X > Y \geq 0$, accounting for all $45$ escape multisets.

\textbf{Complete list of escape multisets at $d = 8$} (listed as
sorted-descending tuples):

$(1, 1, 1, 1, 0, 0, 0, 0), \; (2, 2, 2, 2, 0, 0, 0, 0), \; \ldots, \; (9, 9, 9, 9, 0, 0, 0, 0)$
--- the $9$ multisets with $Y = 0$.

$(2, 2, 2, 2, 1, 1, 1, 1), \; (3, 3, 3, 3, 1, 1, 1, 1), \; \ldots, \; (9, 9, 9, 9, 1, 1, 1, 1)$
--- the $8$ multisets with $Y = 1$.

$\vdots$

$(9, 9, 9, 9, 8, 8, 8, 8)$ --- the $1$ multiset with $Y = 8$.

\textbf{Total: $9 + 8 + 7 + 6 + 5 + 4 + 3 + 2 + 1 = 45$ multisets.}

\textbf{Why these escape.} Each escape multiset
$(X, X, X, X, Y, Y, Y, Y)$ has only two distinct digit values, with
exactly four copies of each. Applied to the best lifting:

\[K(X^4 Y^4) = |c_0 \cdot X + c_1 \cdot X + c_2 \cdot X + c_3 \cdot X + c_4 \cdot Y + c_5 \cdot Y + c_6 \cdot Y + c_7 \cdot Y|\]
\[= |(c_0 + c_1 + c_2 + c_3) \cdot X + (c_4 + c_5 + c_6 + c_7) \cdot Y|\]
\[= |(999 + 90 - 90 - 999) \cdot X + (-900{,}000 + 9{,}990{,}000 + 900{,}000 - 9{,}990{,}000) \cdot Y|\]
\[= |0 \cdot X + 0 \cdot Y| = 0.\]

The first four coefficients sum to $0$ (since $K_0$ has sum-zero
structure), and the last four coefficients sum to $0$ (by the
coefficient-preserving lifting construction). When the input has only
two distinct values, the rule's action collapses algebraically to zero.

This is not a feature of the specific best lifting --- it is a
structural consequence of sum-zero constraints on coefficient-preserving
liftings. Every coefficient-preserving lifting at $d = 8$ maps every
four-of-a-kind-paired multiset to $0$.

\subsubsection{\texorpdfstring{At $d = 9$}{At d = 9}}\label{at-d-9}

Exhaustive audit of all $5{,}280$ coefficient-preserving liftings at
$d = 9$, each tested against all $48{,}520$ admissible multisets
with iteration budget $60$:

\begin{longtable}[]{@{}lc@{}}
\toprule\noalign{}
Status & Count \\
\midrule\noalign{}
\endhead
\bottomrule\noalign{}
\endlastfoot
Strict-universal (basin $= 1.0$) & $0$ \\
Near-universal (basin $\geq 0.99$) & $302$ \\
Best basin & $0.999073$ \\
\end{longtable}

Runtime: approximately $60$ minutes on commodity hardware.

\subsubsection{\texorpdfstring{Escape characterization at
$d = 9$}{Escape characterization at d = 9}}\label{escape-characterization-at-d-9}

For the best lifting at $d = 9$, the $45$ escape multisets have the
same structural form as at $d = 8$, extended with trailing zeros.
Specifically: sorted-descending $(X, X, X, X, Y, Y, Y, Y, 0)$ for
$X > Y \geq 0$ (when $Y \geq 1$) or $(X, X, X, X, 0, 0, 0, 0, 0)$
(when $Y = 0$). There are again $\binom{10}{2} = 45$ such multisets.

\textbf{The escape count is independent of $d$.} The $45$-multiset
structural form is preserved across $d = 8, 9$ and is expected to
continue at $d \geq 10$: at higher $d$, the escape multisets are
obtained by appending trailing zeros to the $d = 8$ escape forms.

\subsection{Asymptotic behavior}\label{asymptotic-behavior}

\textbf{Basin growth.} The escape count at $d = 8, 9$ is fixed at
$45$. The admissible multiset count grows with $d$:

\begin{longtable}[]{@{}
  >{\centering\arraybackslash}p{(\columnwidth - 6\tabcolsep) * \real{0.2500}}
  >{\centering\arraybackslash}p{(\columnwidth - 6\tabcolsep) * \real{0.2500}}
  >{\centering\arraybackslash}p{(\columnwidth - 6\tabcolsep) * \real{0.2500}}
  >{\centering\arraybackslash}p{(\columnwidth - 6\tabcolsep) * \real{0.2500}}@{}}
\toprule\noalign{}
\begin{minipage}[b]{\linewidth}\centering
$d$
\end{minipage} & \begin{minipage}[b]{\linewidth}\centering
\# admissible multisets
\end{minipage} & \begin{minipage}[b]{\linewidth}\centering
\# escapes
\end{minipage} & \begin{minipage}[b]{\linewidth}\centering
basin (best)
\end{minipage} \\
\midrule\noalign{}
\endhead
\bottomrule\noalign{}
\endlastfoot
$8$ & $24{,}210$ & $45$ & $0.998141$ \\
$9$ & $48{,}520$ & $45$ & $0.999073$ \\
$10$ & expected $\sim 92{,}000$ & expected $45$ & expected
$\sim 0.9995$ \\
$\cdots$ & $\cdots$ & $\cdots$ & $\cdots$ \\
$d \to \infty$ & $\to \infty$ & $45$ & $\to 1$ \\
\end{longtable}

\textbf{Asymptotic claim.} The best-lifting basin for $6174$
approaches $1$ as $d \to \infty$. Formally, if the $45$-escape
structural form persists at all $d \geq 8$, then

\[\text{best basin at } d = 1 - \frac{45}{|A_d|} = 1 - O(10^{-d/9}).\]

The constant-count escape class means the ratio $45 / |A_d|$ decays as
$|A_d|$ grows, which for admissible multisets at digit length $d$
grows combinatorially.

\textbf{Strict-universal at any $d \geq 8$?} Audit at $d = 8, 9$
finds $0$ strict-universal liftings. Whether any $d \geq 10$ admits
a strict-universal coefficient-preserving lifting of $6174$ remains
open. The structural argument of §E.2.2 suggests the answer is no: the
four-of-a-kind-paired escape class is algebraically forced by sum-zero
on the native and appended coefficients, and appears in every
coefficient-preserving lifting.

\subsection{Conjectural extension}\label{conjectural-extension}

\begin{conjecture}[E.1]
For every $d \geq 8$, the
$45$-multiset escape class persists at $6174$'s
coefficient-preserving liftings, and no coefficient-preserving lifting
is strict-universal. The best basin asymptotes to $1$ monotonically.
\end{conjecture}

Resolving Conjecture E.1 requires a structural argument at general $d$
beyond the direct audit at $d = 8, 9$. The algebraic proof of the
$45$-multiset escape class's persistence (that it is forced by
sum-zero constraints regardless of which coefficient-preserving lifting
is chosen) is outlined in §E.2.2 and could likely be extended to all
$d$, but we do not pursue this here.

\emph{End of Appendix E.}



\begin{thebibliography}{99}

\bibitem{Kaprekar1955} D. R. Kaprekar, ``An interesting property of the
number 6174,'' \emph{Scripta Mathematica}, vol.~15, pp.~244--245, 1955.

\bibitem{Trigg1974} C. W. Trigg, ``All three-digit integers lead to 495,''
\emph{The Mathematics Teacher}, vol.~67, pp.~41--45, 1974.

\bibitem{Prichett1981} G. D. Prichett, A. L. Ludington, and J. F. Lapenta,
``The determination of all decadic Kaprekar constants,'' \emph{The
Fibonacci Quarterly}, vol.~19, pp.~45--52, 1981. (Precursor: G. D.
Prichett, ``Kaprekar's routine with five-digit integers,'' \emph{The
Fibonacci Quarterly}, vol.~17, pp.~154--159, 1979.)

\bibitem{Peterson2008} A. Peterson and J. Pulapaka, ``Kaprekar's routine
and other number sequences in different bases,'' \emph{Virginia
Mathematics Teacher}, vol.~35, no. 1, 2008.

\bibitem{Nuez2021} F. Nuez, ``Parametric transformation functions in the
Kaprekar routine I,'' arXiv:2110.15123, 2021.

\bibitem{Nuez2021b} F. Nuez, ``Algebraic Kaprekar routine architecture
II,'' arXiv:2111.00932, 2021.

\bibitem{Nuez2022} F. Nuez, ``Symmetries in generalized Kaprekar's
routine,'' \emph{Symmetry}, vol.~14, no. 1, p.~37, 2022. (Also
arXiv:2110.13730.)

\bibitem{Devlin2021} M. Devlin and T. Zeng, ``Maximum distances in the
four-digit Kaprekar process,'' \emph{Integers}, vol.~21, A52, 2021.

\bibitem{Iwasaki2024} H. Iwasaki, ``A New Classification of the Kaprekar
Numbers,'' \emph{The Fibonacci Quarterly}, vol.~62, no. 4, pp.~275--281,
2024.

\bibitem{Kay2022} A. Kay and K. Downes-Ward, ``Fixed Points and Cycles of
the Kaprekar Transformation: 1. Odd Bases,'' \emph{Journal of Integer
Sequences}, vol.~25, Article 22.6.7, 2022.

\bibitem{Kay2024} A. Kay and K. Downes-Ward, ``Fixed Points and Cycles of
the Kaprekar Transformation: 2. Even bases,'' arXiv:2408.12257, 2024.

\bibitem{Dahl2026} C. D. Dahl, ``Coarse-Grained Drift Fields and
Attractor-Basin Entropy in Kaprekar's Routine,'' \emph{Entropy},
vol.~28, p.~92, 2026.

\emph{End of paper.}

\end{thebibliography}

\end{document}
